%XeLaTeX
\documentclass{article}
\setlength{\emergencystretch}{10pt}
\usepackage{microtype}
\usepackage{float}
\usepackage{fontspec}
\usepackage{polyglossia}

\setdefaultlanguage{english}
\setotherlanguages{tamil,greek}

% Set main font (used for English, the default language)
\setmainfont{Charis SIL}
\newfontfamily{\englishfont}{Charis SIL}

% Set specific fonts for other languages
\newfontfamily{\greekfont}[Script=Greek]{Gentium}
\newfontfamily{\tamilfont}[Script=Tamil,ItalicFont={NotoSerifTamilSlanted-Regular},BoldFont={NotoSerifTamil-Bold}]{FreeSerif}
\newcommand{\TamilText}[1]{\texttamil{#1}}
\newcommand{\TamilItalic}[1]{\texttamil{\emph{#1}}}
\newcommand{\TamilBold}[1]{\texttamil{\textbf{#1}}}
\newcommand{\tam}[1]{\TamilText{#1}}
\newcommand{\tamitalic}[1]{\TamilItalic{#1}}
\newcommand{\tambold}[1]{\TamilBold{#1}}
\newcommand{\TamilHead}[1]{\TamilItalic{#1}}
\newcommand{\TamilEtymon}[1]{\TamilText{#1}}
\newcommand{\Translit}[1]{\emph{#1}}
\defaultfontfeatures{Scale=MatchLowercase}
\def\arraystretch{1.15}

\usepackage{graphicx}
\graphicspath{ {./ } }
\usepackage[figurename=]{caption}
\usepackage{float}

\usepackage{tocloft}
\usepackage[dvipsnames]{xcolor}
\usepackage[object=vectorian]{pgfornament} %% http://altermundus.com/pages/tkz/ornament/index.html

\begin{document}
\begin{titlepage} % Suppresses headers and footers on the title page
	\centering % Centre everything on the title page
	%\scshape % Use small caps for all text on the title page

	%————————————————
	%	Title
	%————————————————

	\rule{\textwidth}{1.6pt}\vspace*{-\baselineskip}\vspace*{2pt} % Thick horizontal rule
	\rule{\textwidth}{0.4pt} % Thin horizontal rule
	
	\vspace{1\baselineskip} % Whitespace above the title
	
	{\scshape\Huge The Tamil Plutarch}
	
	\vspace{1\baselineskip} % Whitespace above the title

	\rule{\textwidth}{0.4pt}\vspace*{-\baselineskip}\vspace{3.2pt} % Thin horizontal rule
	\rule{\textwidth}{1.6pt} % Thick horizontal rule
	
	\vspace{1\baselineskip} % Whitespace after the title block
	
	%————————————————
	%	Subtitle
	%————————————————
	

    {\scshape\small Containing a Summary Account of the Lives of the Poets and Poetesses of Southern India and Ceylon\\From the earliest to the present Times, with select Specimens of their Compositions} 

    \vspace{1.0\baselineskip}

    {\scshape\small \emph{Aut prodesse volunt aut delectare poetæ.} --- Horace} % Subtitle or further description

    \vspace{1.0\baselineskip}

    {\scshape\small By \large Simon Casie Chetty,\small Esquire\\Author of the Ceylon Gazetteer}
‌
	%————————————————
	%	Editor(s)
	%————————————————


    \vspace*{\fill}

	\vspace{1\baselineskip}

    {\scshape\small Ripley \& Strong --- Printers\\ Jaffna 1859}
		
	\vspace{0.25\baselineskip} % Whitespace after the title block

    {\scshape\small Solar Anamnesis Edition}% Publication year}

	{\scshape\footnotesize CC0 1.0 Universal} % Publisher
\end{titlepage}
\clearpage
\tableofcontents
\clearpage
\section*{Preface.}
\paragraph{}
Of the languages of the Scythian family the Tamil confessedly occupies the most distinguished rank, and it is peculiar to the people of that part of India, which was formerly under the sway of the \emph{Chera}, \emph{Chola} and \emph{Pandiya} kings and of those of the eastern and northern provinces of Ceylon. The name Tamil, signifying ``sweet,'' is characteristic of the language. Indeed it is one of the most copious, refined, and polished languages spoken by man, as correctly observed by an accomplished Orientalist.\footnote{Taylor's Preface to Dr. Rotler's Tamil and English Dictionary, part 4.} Villiputt\={u}rer, in his Introductory stanza to the \emph{Baradam} which we quote below, metaphorically represents the Tamil language as the daughter of Agastiyer, who being born of the mountain (\emph{Pothiya}) and cradled in the glory of Pandiyen, sat upon the College bench (at Madura), crawled with her breast on the writing (of Sampanter) ascending against the stream of the \emph{Vygai} river, stood up amidst the fire unburned (when that writing was cast into it), and afterwards walked on the minds of the students, and now adorns the hip of the goddess of the earth whom (\emph{Vishnu}, in his metamorphosis as) a boar supported on the point of his tusks.
\begin{quotation}\small
பொருப்பிலேபிறந்துதென்னன்‌ புகழிலேகிடந்துசங்கத்‌

திருப்‌பிலேயிருந்துவைகை யாற்றிலே தவழ்ந்தபேதை

நெருப்பிலே நின்றுகற்போர்‌ நினைவிலேநடந்தோரேன

மருப்பிலேபயின்றபாவை மருங்கிலேவளருகின்றாள்‌.
\end{quotation}
\paragraph{}
Few nations on earth can perhaps boast of so many poets as the Tamils. Poetry appears to have been the first fixed form of language amongst them; for as has been remarked by Abbe Dubois, ``they have not a single ancient book that is written in prose, not even the books on medicine;''\footnote{Abbe Dubois' Description of the People of India, p. 260.} and hence the poet formed the inalienable part of the philologist, the theologian, the philosopher, the astronomer, the physician, etc. The inducements held out to poets and the rewards bestowed on them by the long line of \emph{Pandiya} kings, who graced the throne of Madura from the ninth century before to the fourteenth century after Christ, were most liberal, and might have done honour even to the court of Augustus. These kings had three different \emph{Sangams}, or Colleges established in their capital at three different periods, for the promotion of literature, more or less corresponding in character with the Royal Academy of Sciences founded by Louis 14th at Paris, and made it a rule that every literary production should be submitted to their Senatus Academicus, before it was allowed to circulate in the country, for the purpose of preserving the purity and integrity of the language. It may be well imagined how favourably these \emph{Sangams} operated on the talent and genius of the nation. From every part of Southern India poets crowded into the \emph{Sanga-mandapam}, or College hall to recite their compositions and the successful candidate besides winning the smiles of Royalty was rewarded with something more enduring and substantial as will appear from Vamshas\={u}d\={a}mani Pandiyen presenting a purse of gold to the poet Tarumi, and Kul\={e}sa Pandiyen honouring the poet Iddeik\={a}der by the gift of a young elephant and a horse, besides gold, and fertile lands. Neither were the kings of \emph{Ch\={e}ra} and \emph{Ch\={o}la} backward in patronizing poets; for they had a certain number of them always attached to their courts, and the names of Perum\={a}k\={o}thei Ch\={e}ra and Kul\={o}tunga Ch\={o}la are still celebrated, the former for his unbounded munificence towards the poets in general, and the latter towards the poet Ottak\={u}ter, the author of the war-chant called \emph{Kalingattu Parani}. There can be no doubt that an infinite number of works in the different departments of sciences and literature was composed during this brilliant age; but in the early part of the fourteenth century when the Muhammedan hordes poured into Southern India, and Pr\={a}krama Pandiyen was led away captive to Delhi, the Tamils had to deplore the loss of almost all their literature; for those ruthless fanatics amongst other outrages ransacked all the libraries in the country, and committed to the flames ``all that genius had reared for ages.'' But auspicious days were yet to dawn upon the Tamils. The power by which they were oppressed has passed away and is succeeded by one destined by Providence to ameliorate their social condition, and advance their intellectual status. With the occupation of India by the British, that nation has introduced into the country everything that has tended to make its own great. Not only has every facility been afforded for the improvement of the Tamil youth by the establishment of vernacular schools by Government as well as Christian Missionaries, but that mighty engine, the Press, has proved in India as elsewhere the greatest blessing, on the one hand by rescuing from oblivion the remnants of the ancient Native literature, and on the other by enriching it with numerous publications on European arts and sciences.

It had long been my intention to offer to the public an account of the Lives of our Poets, but the difficulties which presented themselves in procuring materials for the work owing to the absence of biographical records in Tamil either ancient or modern, was more than I had calculated upon, and I almost resolved to give up the undertaking in despair. Finding however that the traditions current among the people, carefully collected and scrupulously detached from fictitious and ornamental additions such as oriental imagination delights in, would afford the desired information, I was induced to direct my attention to that source of knowledge, and the result with all its imperfections is now before the public.

This work being the first attempt of the kind, must be notwithstanding all my vigilance necessarily subject to many inaccuracies and defects, for which I solicit the indulgence of the reader and the critic. Should my life be spared and should my friends be kind enough to communicate to me such additions or corrections as they think the work needs, I may enlarge and improve it for a new edition at some future period.

\bigskip

S. C. C.

\emph{Chilaw, August} 15, 1859.
\clearpage
\section{Agastiyer. --- அகஸ்தியர்.}
\paragraph{}
The history of Agastiyer being blended with the Hindu Mythology, it is now difficult to separate the real particulars of his life from those that are fabulous. This much, however, is certain, that he was a native of \emph{\={A}ry\={a}varta}, otherwise called \emph{Puṇyabh\={u}mi}, or the holy land, situated between the \emph{Imaus} and \emph{Vindhya} mountains, and that he led the first colony of \emph{Br\={a}hmans} which settled in Southern India from the north under the auspices of \textbf{Kulasekhara}, the founder of the \emph{P\={a}ndiyen} dynasty at Madura. The fable of his levelling the \emph{Vindhya} mountain, which had elevated its summits even to the skies and impeded the progress of the sun and moon, implies, his passing over it into the south with the \emph{Br\={a}hmans} from the north; while the other fable that on the occasion when all the gods and sages assembled together on the \emph{Imaus} mountain, to celebrate the nuptials of \emph{Siva} and \emph{P\={a}rvati} the northern part of the earth became depressed, and the southern part elevated, and \emph{Siva}, dismissed Agastiyer sending him to the \emph{Pothiya} mountain to reduce both parts to the same level, implies, his civilizing the people of the south and rendering them equal to the people of the north in point of civilization. Not only is Agastiyer believed to have diffused the worship of \emph{Siva} amongst the Tamils, but also to have taught them Medicine, Chemistry, Astronomy, and the other sciences. The legend, however, which attributes to him the invention of the Tamil alphabet, is far from probable. Perhaps he merely enlarged it by inserting in it such of the letters as are now used to represent the sounds of Sanskrit derivatives. Although it cannot be denied that the Tamils owe much of their civilization to Agastiyer, it is nevertheless not to be supposed, that before he arrived amongst them they were altogether immersed in barbarism, and that they were not acquainted with the art of writing. Such a supposition is disproved by the existence of a word for ``\emph{a letter}'' (எழத்து), and another for ``\emph{a book}'' (சுவடி), in the pure Tamil as it was spoken in the primitive times. However it be, there is no doubt that Agastiyer was the first, who arranged the Tamil grammar on the principles of the Sanskrit, and moulded the Tamil language itself into its present elegant form, and, therefore, justly earned for himself the title of \emph{Tamil-muni}, the Tamil sage. His treatise on Tamil grammar is said to have contained no less than 12,000 S\={u}tras or aphorisms; but with the exception of some fragments which have been preserved in quotations by \textbf{Tolk\={a}piyan\={a}r}, in his work on the same subject, it is not now extant. Agastiyer wrote both in Sanskrit, and in Tamil. He is reputed to have written a lakh of verses on theology, an equal number on alchemy, and two lakhs on medicine, but most of the Tamil works commonly ascribed to him are in reality not his productions. In the medical treatise, entitled \emph{Irrun\={u}ttanju} (இருநூற்றஞ்சு), lues veneria, is described under the name of \emph{Paranghiviy\={a}dhi} or ``the Portuguese disease,'' and again in the alchemical treatise, entitled \emph{P\={u}ranas\={u}tram} (பூரணசூத்திரம்), sublimate of mercury is described under the name of \emph{Paranghip\={a}sh\={a}nam}, or ``the Portuguese mineral''; hence these two works bear undoubted marks of having been written after the arrival of the Portuguese in India.

As a physician, Agastiyer occupies the same eminence amongst the Tamils as \textbf{Hippocrates} does amongst the Greeks, and it is deserving of notice, that there are some very curious coincidences between the doctrines of the former, and those of the latter, especially as regards the prognosis and diagnosis of diseases, the critical days, and premonitory symptoms of death. His knowledge of medicine was confessedly extensive, nor was he less remarkable for his skill in surgery. The story of his having once opened the cranium of a patient, in order to remove a toad, which had lodged itself in the brain, when disentangled from the twist of fable, at once proves that he was acquainted with trepanning; and it should be remarked that the existence of seminal animalcules, which was discovered by \textbf{Ludwig Hamm} in Europe only in \textsc{ce} 1677\footnote{See English Cyclopædia, Biography, Vol. 3. p. 871.} is mentioned by Agastiyer in one of his medical works, entitled \emph{Kurun\={a}dich\={u}tram} (குருநாடிச்சூத்தீரம்).\footnote{Roberts' Oriental Illustrations of the Sacred Scriptures, p. 281.} He was also well skilled in chemistry, especially in that part of it which relates to the analysis and synthesis of minerals, and the vehicles by the aid of which the transmutation of metals is effected.

Agastiyer was a married man. It is said that his wife \textbf{Lop\={a}mudra} was a pattern of conjugal affection, who attracted universal admiration on that account, and that he had a son by her named S\={a}garen. At what period Agastiyer established himself in Southern India is not known, and it must always remain so until we shall have been able to ascertain the real date of the existence of the king \textbf{Kulasekhara Pandiyen}, by whom he was patronized. All accounts concur in assigning the foundation of the \emph{Pandiya} kingdom at Madura to \textbf{Kulasekhara Pandiyen}; but they are at considerable variance with regard to the time when that event happened. Some place it as high up as \textsc{bce} 1500,\footnote{Taylor's Oriental Historical Manuscripts, vol. 1. p. 135.} while others bring it down to a later period; but we have reasons to believe that it could not have been later than, at least, the ninth century \textsc{bce}; for from the fact of \textbf{Wijaya}, who founded the \emph{Singhalese} dynasty in Ceylon in \textsc{bce} 543, marrying the daughter of a \emph{Pandiyen}, and the \emph{Pandiyen} sending along with her to Ceylon a retinue of seven hundred daughters of his nobility, together with eighteen officers of state, and seventy five menial servants, besides elephants, horses, and chariots, as recorded in the \emph{Mahawanso}, it would appear that the \emph{Pandiya} kingdom was not then in its infancy, but fully organized and at the lowest calculation a few centuries old. When Agastiyer left the court of \textbf{Kulasekhara Pandiyen}, he is stated to have assumed the ascetic life, and to have retired to the \emph{Pothiya} mountain, where the Hindus absurdly believe him to be still alive, although invisible to ordinary eyes.

Agastiyer was a theist (ஏகேச்சுரவாதி) in philosophy, and the Rev. \textbf{Mr. Taylor}, who calls him ``the prince of Indian doctors,'' and ``one of India's greatest Philosophers,'' justly observes, that in ``so far as a superiority to popular errors is concerned, he may perhaps claim to rank on a level with \textbf{Socrates}, or \textbf{Plato}.''\footnote{Taylors' Oriental Historical Manuscripts, Vol. 1. p. 172. p. 175.}

The following Sanskrit hymns or prayers written by Agastiyer, and which form a part of the \emph{Agastiyer Ashtakam} as quoted by \textbf{Mr. Brown}, in his Essay on the Creed, Customs, and Literature of the Jangamas,\footnote{Madras Journal of Literature and Science, Vol. 11. p. 161.} will afford an idea of the sentiments he held regarding the attributes of God.
\begin{quotation}\small
பரமானந்தம்‌ பரமசுகதம்‌ கேவலஞானமூர்த்திம்‌

துவந்துவாதீதம்‌ ககனசதிரிசம்‌ தத்வம்‌ அசியாதிலக்ஷியம்‌

ஏகம்‌ நித்யம்‌ விமலமசலம்‌ சர்வதாசாக்ஷிபூதம்‌

பாவாதீதம்‌ திரிகுணரகிதம்‌ சற்குரும்‌ தம்நமாமி.

வடவிதாபிசமீபேபூமிபாகேநிஷண்ணம்‌

சகலமுனிஜனானம்‌ ஞானதாதரம்‌ ஆராத்‌

திரிபுவனகுரும்‌ ஈசாம்‌ தக்ஷணாமூர்த்திதேவம்‌

ஜனனமரணதுக்க க்ஷேததக்ஷ்ம்நமாமி.
\end{quotation}
\begin{quotation}\small
1.\emph{I salute the great teacher, the bestower of divine happiness and supreme bliss; the image of perfect wisdom; who is removed from all griefs; who is represented by the sky; who is denoted by the ``Truth'' and other names. The one, Eternal, stainless, stable, and Omniscient, the incomprehensible, who knoweth neither Passion, Partiality nor Folly.}

2.\emph{He who sits on earth at the foot of the fig-tree: who bestows wisdom on all the devout hermits who surround him: Lord and teacher of the universe, the God who is embodied Goodness, him do I salute as the releaser from the bonds of life and death.}
\end{quotation}
\paragraph{}
Agastiyer is said to have had twelve disciples to whom he taught the different arts and sciences, and who were afterwards employed by him in instructing the people. The names of these disciples are \textbf{Tolk\={a}ppiyan\={a}r}, \textbf{Adankot\={a}s\={i}r\={i}yan\={a}r}, \textbf{Tur\={a}lingan\={a}r}, \textbf{Semp\={u}tcheyan\={a}r}, \textbf{Vaiy\={a}bigan\={a}r}, \textbf{V\={a}ippiyan\={a}r}, \textbf{Panamb\={a}ran\={a}r}, \textbf{Kazh\={a}ramban\={a}r}, \textbf{Avinayan\={a}r}, \textbf{K\={a}keip\={a}diniyan\={a}r}, \textbf{Nattattan\={a}r} and \textbf{V\={a}man\={a}r}; but few particulars are known respecting them.

\section{Akk\={a}rakkani Nachuman\={a}r. --- அக்காரக்கனிநச்சுமனார்.}
\paragraph{}
This poet, was one of the forty-nine professors of the Madura College, but we are not in possession of any particulars of his life, and of his poetical compositions nothing remains, save the following impromptu, which he uttered on the \emph{Kural} when that work was submitted to the judgment of the College bench.
\begin{quotation}\small
கலைநிரம்பிக்காண்டற்கினிதாகிக்கண்ணி

னிலைநிரம்புநீர்மைத்தெனினுந் --- தொலைவிலா

வானூர்மதியந்தனக்குண்டோவள்ளுவர்முப்‌

பானூனயத்தின்பயன்‌.
\end{quotation}
\begin{quotation}\small
\emph{The moon full of} kalei \emph{(the whole of her face being illuminated) pleaseth the external eyes, in like manner as the Kural full of} kalei \emph{(knowledge) pleaseth the intellectual eyes; but nevertheless she cannot be compared to Valluver's production; for she is neither spotless, nor does she retain her form and splendour unchanged like it.}
\end{quotation}
\section{Amirtakavir\={a}yer. --- அமிர்தகவிராயர்.}
\paragraph{}
Amirtakavir\={a}yer was a native of \emph{Ponn\={a}nk\={a}l}, in the Carnatic. His early display of poetical talents, which were of no mean order, having attracted the notice of \textbf{Reghun\={a}tha Setupati}, who reigned at Ramnad between \textsc{ce} 1649 and 1685, he obtained a place in the court of that king as his encomiast. He composed an erotic poem, under the title of \emph{Orutureikk\={o}vei} (ஒருதுறைக்கோவை), in honour of his royal benefactor, which is admired as a matchless production of the kind on account of its splendid diction, harmonious versification, and copious poetical imagery; but as the words used throughout the work have two or more meanings, it puzzles the interpreters.

\section{Amirtas\={a}garer. --- அமிர்தசாகரர்.}
\paragraph{}
Amirtas\={a}garer was a poet of great attainments, but we possess few particulars respecting his life. It appears from the invocation at the beginning of a treatise, which he wrote on Tamil versification, under the title of \emph{K\={a}rigai} (காரிகை), that he was a follower of the \emph{Jaina} religion; and \textbf{Gunas\={a}garer}, his commentator, represents his proficiency in Sanskrit to have been equally as considerable as in Tamil. The above mentioned treatise consists of forty stanzas, which are composed throughout in the \emph{Kalitturei} metre, with much elegance, and is inscribed to a lady, whose name, however, the author has not mentioned.

\section{Anandabh\={a}rati Aiyang\={a}r. --- அநந்தபாரதிஜயங்கார்.}
\paragraph{}
This poet was the son of \textbf{Sr\={i}niv\={a}sa Aiyang\={a}r, a \emph{Mirasidar},\footnote{Hereditary landholder.} who lived at \emph{Umeiyamm\={a}lpuram}, near Tanjore. He was born in \textsc{ce} 1786. Having commenced his studies at the early age of five before his thirteenth year he gave promise of his future celebrity, by composing a drama of the class, called \emph{Nondi}, in praise of \emph{Y\={a}neimelazhagar}, the tutelary deity of his village. Soon afterwards he went with his father to Trivalore, and remained there three years, spending his time in teaching a school. In his fifteenth year he entered the public service, and was successively employed as \emph{Karnam},\footnote{Accountant.} and \emph{Samprati}\footnote{Auditor of accounts.} of different temples in the Tanjore district, until his twenty-fifth year, when he resigned the office and establishing himself at Tiruvadamarud\={u}r, continued to devote the remainder of his life to literary pursuits. Having composed and recited a drama at the shrine at Tiruvadamarud\={u}r, \textbf{Aiy\={a}rappa Tambir\={a}n}, the \emph{Dharmakarta}\footnote{Manager or Supervisor.} of the shrine, was so much pleased with its exquisite beauty, that he honoured him with the gift of a house and garden, and conferred on him the title of \emph{Kavir\={a}jaswami}. Amongst his other productions the \emph{Uttara R\={a}m\={a}yana K\={i}rtana} (உத்தரராமாயணகீர்த்தண), \emph{D\={e}sikaprabandham} (தேசிகப்பிரபந்தம்), \emph{Bh\={a}gavata Dasamaskanda N\={a}dagam} (பாகவத்தசமஸ்கந்தநாடகம்), \emph{Mar\={u}d\={u}r Venba} (மருதூர்வெண்பா) and \emph{Mupp\={a}ttirattu} (முப்பாற்றிரட்டு) are much admired, and there are also some hymns composed by him in honour of the \emph{Saiva} shrines at Trichinapoly, Tirukudantei, Tiruvavaduturei, etc. which do him great credit. He died in \textsc{ce} 1846, in the sixtieth year of his age.

\section{Andaka Kaviv\={i}rar\={a}gava Mudeli\={a}r. --- அந்தகக்கவிவீரராகவமுதலியார்.}
\paragraph{}
Andaka Kaviv\={i}rar\={a}gava Mudeli\={a}r was an eminent poet, who is said to have been born blind, and by a special gift to have become inspired with poetic spirit. Whether the story of his inspiration be founded on fact or not he has certainly rendered himself famous by his two poems, one called \emph{Kazhukkunda Pur\={a}nam} (கழக்குன்றப்புராணம்) and the other \emph{Kazhukkunda M\={a}lei} (கழக்குன்றமாலை), which he composed in praise of the \emph{Saiva} shrine on the \emph{Kazhukkunda} mountain, in the Carnatic. His \emph{Kazhukkunda M\={a}lei} is confessedly a splendid production; but like other Tamil poems it is replete with hyperboles. We here extract a stanza from it as a specimen: 
\begin{quotation}\small
மாடேறுதாளுமதியேறுசென்னியுமாமறையோ

னோடேறுகையுமுடையார்தமக்கிடமூருழவர்‌

சூடேறுசங்கஞ்சொரிமுத்தைமுட்டையென்றேகமலக்‌

காடேறுமன்னஞ்‌ சிறகாலணைக்குங்‌ கழுக்குன்றமே.
\end{quotation}
\begin{quotation}\small
\emph{He who rides the bull, wearing the moon on his head, and holding the skull of Br\={a}hma in his hand, has his dwelling in Kazhukkundam, where the swan mistaking for her eggs the pearls scattered by the chanks\footnote{A species of Voluta found in corn-fields.} on the sheaves of corn, cut by the husbandmen, gathers them under her wings and broods over them.}
\end{quotation}
\paragraph{}
The following is a poetical epistle, which he addressed to a nobleman soliciting a boon;
\begin{quotation}\small
ஏடாயிரங்கோடியெழுதாமறன்மனத்‌ தெழுதிப்படித்தவிரகன்‌ --- இமசேதுபரியந்த மெதிரிலாக்கவிவீரராகவன்விடுக்குமோலை --- சேடாதிபன்சிரமசைத்திடும்புகழ்பெற்ற திரிபதகைகுலசேகரன்‌ --- தென்பாலைசேலம்புரந்துதாகந்தீர்த்த செழியனெதிர்கொண்டுகாண்க --- பாடாதகந்தருவமெறியாதகந்துகம்பற்றிக்கொலாதகோணம்‌ --- பறவாதகொக்‌கனற்‌ பண்ணாதகோடைவெம்படையிற்றொடாதகுந்தம்‌ --- சூடாதபாடலம்‌ பூவாதமாவொடுதொடுத்து முடியாதசடிலம்‌ --- சொன்னசொற்‌சொல்லாதகிள்ளையொன்றெங்குந் துதிக்கவரவிடல்வேண்டுமே.
\end{quotation}

\section{Anthonikutty Ann\={a}viy\={a}r. --- அந்தோனிக்குட்டியண்ணாவியார்.}
\paragraph{}
Anthonikutty Ann\={a}viyar was born at \emph{Manapar}, in the Tinnevelly district, and was contemporary with the celebrated Father \textbf{Beschi}. Some say that he was a \emph{Parava}, and others that he was a \emph{Yiluva} by caste; but be this as it may, he was the best of the poets that ever flourished amongst the Tamil Catholics. Having in the early part of his life indulged himself in licentious pleasures without restraint, and despised the admonitions of the clergy, he was excommunicated from the church, but by a miraculous interference as in the case of St. \textbf{Augustine}, he forsook his dissolute career and for the remainder of his life employed all his poetical talents in singing the praises of the Redeemer. All his poems have been recently collected in a single volume, and edited at Jaffna, under the title of ``Christian Psalmody'' and, as the Editor in his introduction justly observes, they are ``without parallel in piety, in strains, and in melody,'' and they ``form not only one of the most efficient helps to piety and devotion, but a pleasing instructor in the doctrines, precepts, and duties of our holy religion.''

\section{App\={a}vaiyer. --- அப்பாவையர்.}
\paragraph{}
App\={a}vaiyer, a poet and astronomer of considerable celebrity of whose personal history however nothing is known, except that he was a \emph{Jaina Br\={a}hman}. The \emph{T\={a}ndava M\={a}lei} (தாண்டவமாலை), an Astronomical treatise which he wrote is in great repute in Southern India, and it comprises in sixty-eight stanzas the substance of the whole of the \emph{Hor\={a}s\={a}ra}, \emph{Par\={a}s\={a}rya}, and other Sanskrit works on astronomy. It was written under the patronage of \emph{T\={a}ndav\={a}r\={a}yen}, a wealthy merchant of Nellikuppam after whom it has been named.

\section{Apper. --- அப்பர்.}
\paragraph{}
Apper, otherwise called \emph{V\={a}kiser}, was born at \emph{Tiruv\={a}mur}, a town situated on the Pennei river in the Carnatic. His father \textbf{Pugazhan\={a}r}, and his mother \textbf{M\={a}dini} were both of the tribe of \emph{Kurukkei Vel\={a}lers}. Being staunch \emph{Saivas}, they brought up Apper and his sister \textbf{Tilakavati}, in their own religion; but no sooner had they died than he entered a \emph{Jaina} monastery at P\={a}taliputra, and having studied the \emph{Jaina} books embraced their creed, and became their high priest, under the title of \emph{Dharmasena}. Sometime afterwards he was attacked with arthritis in his stomach which continued to baffle all medical efforts. His sister \textbf{Tilakavati}, who zealously adhered to the worship of \emph{Siva}, availing herself of the opportunity, endeavoured to persuade him into a belief, that his disease was a punishment from \emph{Siva} for his apostasy. He yielded to her persuasions, and in the dead of night, flying from the \emph{Jaina} monastery went to Tiruvathigai, where his sister resided, and throwing off the robe, resumed the \emph{Saiva} marks, by rubbing his forehead with ashes, and wearing a necklace of \emph{Ruddr\={a}ksha} beads around his neck. His subsequent history is entirely involved in fable. It is said that when his relapse was reported by the \emph{Jainas} to their king \textbf{Pallavar\={a}yen}, he sent for him, and after having tried to kill him by successively casting him into a burning lime-kiln, exposing him to an infuriated elephant, and poisoning his food, without any effect, had him thrown into the sea with a stone tied to his neck, but he was miraculously preserved from being drowned, the stone which was tied to his neck serving him as a float to reach the port of Padripuliy\={u}r in safety. One thing however is certain, that after his return to \emph{Saivaism}, he laboured assiduously to supplant the \emph{Jaina} religion throughout Southern India, in conjunction with \textbf{Sampanter} and \textbf{Sundarer}. He is reputed, to have composed no less than 49,000 \emph{Pathigams}, (each consisting of ten stanzas), in praise of the different \emph{Saiva} shrines, which he visited, but they have all perished, except 312. His epoch is uncertain, but as he and his collaborators are eulogized in the Introductory stanzas of the \emph{Kandapur\={a}nam}, a work written in \textsc{ce} 778, as deified heroes, we are inclined to place it, at least, two or three centuries before that period; for the author of the \emph{Kandapur\={a}nam} could not have imposed on the credulity of the people, by representing them as objects of divine worship, at a period close to their existence. In the latter part of his life he fixed his abode at Tiruppugal\={u}r, where he continued to minister in the temple until his death, which occurred when he had attained the advanced age of eighty-one.

\section{Arasak\={e}sari. --- அரசகேசரி.}
\paragraph{}
Arasakesari was no less a personage than the younger brother of the king \textbf{Parar\={a}ja Sekhara Chakravarti}, who reigned at Jaffna towards the close of the fifteenth century of the Christian era. He was not only skilled in Tamil, but also in Sanskrit, which enabled him to produce a very highly finished imitation of \textbf{K\={a}lidasa's} \emph{Reghuvamsa} in Tamil under the same title, in 2,444 stanzas, arranged in 26 chapters.

\section{Arisitkizhar. --- அரிசிற்கிழார்.}
\paragraph{}
No particulars of the life of this poet have been preserved. All that is known of him is that he was one of the forty-nine professors of the Madura College, and that the following impromptu was uttered by him on the \emph{Kural}.
\begin{quotation}\small
பரந்தபொருளெல்லாம்பாரறியவேறு

தெரிந்துதிறந்தோறுஞ்சேரச்சுருங்கிய

சொல்லால்விரித்துப்‌ பொருள்விளங்கச்சொல்லுதல்‌

வல்லாரார்வள்ளுவரல்லால்.
\end{quotation}
\begin{quotation}\small
\emph{Who else but Valluvar is able to separate according to their order all the things which were blended together in the Vedas, and to impart them to the world in a condensed form and with due amplification.}
\end{quotation}
\section{Arunanti. --- அருணந்தி.}
\paragraph{}
Arunanti, styled \emph{Sakal\={a}gama Panditer},\footnote{A title of honour somewhat analogous to that of Doctor of the Law among the Jews.} was a \emph{Saiva} ascetic of \emph{Tureiy\={u}r} in the Carnatic, equally distinguished as a poet and philosopher. The date of his birth and death is not known, but it appears that he was a pupil of \textbf{Meikanda Dever}, who flourished at Venneinell\={u}r in the eleventh century of the Christian era. He has transmitted his name to posterity by a voluminous treatise on the \emph{Saiva} system of metaphysics and theology under the title of \emph{Sivagn\={a}na Sitti} (சிவஞானசித்தி), and which is generally consulted as the best authority on the subject.

\section{Arun\={a}chala Kavir\={a}yer. --- அருணாசலக்கவிராயர்.}
\paragraph{}
This poet was born at \emph{Tilleiy\={a}di}, near Tranquebar \textsc{ce} 1712. His parents were both of the \emph{Velala} caste and appear to have occupied a respectable position in society. From his fifth to his twelfth year he studied Tamil in a school in his own village, and on the death of his parents which happened about this time, he removed to the \emph{Matam}\footnote{A sort of convent, whore Hindu Devotees dwell together.} at Dharmapuram, where under the tuition of several learned \emph{Tambir\={a}ns}, he was enabled to pass through a distinguished course in the Sanskrit and Telugu languages, and also became very proficient in the \emph{Saia} system of metaphysics and theology. His great abilities naturally attracting the attention of the \emph{Tambir\={a}ns}, induced them to try to secure his services, by persuading him to become like one of themselves, with an offer at the same time to make him the superior of their \emph{Matam}; but he excused himself on the ground that matrimony was preferable to celibacy. He however did not enter into the married state until his thirtieth year, remembering the old adage ``\emph{V\={i}v\={a}ham Vidhya n\={a}sti},'' marriage injures learning. When he did marry, he commenced business as a banker, but his love of literature was so great that he never allowed the cares of his calling to interfere with his studies, and he devoted all his leisure hours in reciting and expounding to his friends \textbf{Kamben's} \emph{R\={a}m\={a}yanam}, which he esteemed as surpassing all other works, not even excepting \textbf{Tiruvalluver's} \emph{Kural}. In 1754 having had an occasion to go to Pondicherry, he put up for a night on the way in the \emph{Matam} at Shiyally. It so happened that the \emph{Tambir\={a}n} this \emph{Matam}, named \textbf{Sidambara Pillei}, was one of his fellow students at Dharmapuram, and who therefore quite rejoiced at meeting him after a long separation was loth to part with him on the morrow. The \emph{Tambir\={a}n} had recourse to an artifice, with a view to delay his journey. Shewing him an unfinished \emph{Pallu}\footnote{A poem in imitation of the action of a \emph{Palla}, one of the lowest orders of agriculturists among the Tamils.} which he had commenced years ago, and laid aside for want of leisure, requested him to finish it for him before he left Shiyally, and the request being acceded to, they both parted for the night. Arunachala Kavirayer finished the poem in the course of the same night, and leaving the manuscript with a servant, to be delivered to the \emph{Tamir\={a}n}, when he awoke, resumed his journey before day-break. Though the \emph{Tambir\={a}n} was at first vexed at his friend's abrupt departure, yet when the servant handed over to him the poem, which he found quite complete, his admiration for him increased two-fold, and he determined to make him settle altogether in his neighbourhood. He had a house erected, and by means of a stratagem, which can only be pardoned when we consider the purity of his intention, he sent word to the poet's wife, that her husband had on his way to Pondicherry requested him to send for her and her children, and settle them in the house; and the deception proved entirely successful. When Arunachala Kavir\={a}yer on his way back from Pondicherry called again at Shiyally, and sought a lodging in the \emph{Matam}, the \emph{Tambir\={a}n} after feigning to be very angry with him for his unceremonious departure on the former occasion, took him out for a walk, and at once introduced him to his family. When the first moment of surprise had worn away he felt very gratified at all that the kindness of the \emph{Tambir\={a}n} had prompted him to do, and consented to make Shiyally his permanent residence for the future.

The earliest poetical compositions of Arun\={a}chala Kavir\={a}yer were the \emph{Asomuki N\={a}dagam} (அசோமுகீநாடகம்), \emph{Sh\={i}rk\={a}zhi Pur\={a}nam} (சீர்காழிப்புராணம்), \emph{Sh\={i}rk\={a}zhi K\={o}vei} (சீர்காழிக்கோவை) and \emph{Anum\={a}r Pilleitamil} (அனுமார்பிள்ளைத்தமிழ்); but what has secured for him a lasting fame is the \emph{R\={a}ma N\={a}dagam} (இராமநாடகம்), a drama founded on the exploits of \emph{R\={a}ma Chandra}, and which he composed when he was sixty years old. After he had submitted the \emph{R\={a}ma N\={a}dagam} to the criticism of the savants at Sirangam, he purposed to recite it in the presence of \textbf{Tulaji Mah\={a}r\={a}ja}, the king of Tanjore, but as the place was then besieged by the army of \textbf{Nawab Mad\={a}r Mulk}, he could not gain access to him. He therefore went to Pondicherry to recite it before \textbf{Anandaranga Pillei}, a noble personage there, who however not wishing to hear it prior to its recital at the court, referred him to \textbf{Manali Muttukishna} \emph{Mudeliar}, one of the celebrities of Madras. On reaching Madras, he addressed the following poetical epistle to the \emph{Mudeliar} requesting him to be allowed to recite the drama before him:
\begin{quotation}\small
சித்திதருகிறவல்லி புத்திதருகிறவல்லி தேவர்கள்‌வணங்குவல்லிதில்லைநாயகவல்லி சிவகாமவல்லி யிருதிருவிழிக்கருணையாலே --- சத்தியவாசகனென்று பூமண்டலாதிபர்கள்‌ தாமெங்கணுந்துதிக்குஞ்‌சகலபாஷாநிபுண மணலிமுத்துக்கிஷணசதுரநீவாழிகண்டாய்‌ --- நத்துமலர்குவளைநான்‌ மதியம்நீசெஞ்சாலி நான்௧ருவிமழைமேகம்நீ நளினம்நான்பருதிநீ பிள்ளை நானன்னைநீ நான்கவிஞன்வழுதிநீகாண்‌ --- சுத்தமுளராமாயணந்தனைக் கொண்டுனதுசுமுகமதுபெறவருகிறேன்‌ சொற்பொருளறிந்தளவு சொல்லுவேன்‌ கேட்கவே துணையாகவேணும்நீயே.
\end{quotation}
\begin{quotation}\small
\emph{Hail! O Muttukishna, thou skilled in all languages! Under the benignant auspices of S\={i}vak\={a}mi, the goddess of Tillei (Chillambram), who imparts heavenly bliss and wisdom, and whom the gods adore, thou art everywhere commended by the kings of the earth, as one true to his word. I am the Nymphaea flower, which expands at night, and thou art the moon; I am the corn, and thou art the raining cloud; I am the lotus-flower (which expands at day), and thou art the sun; I am the child, and thou art the father; and I am the poet, and thou art the P\={a}ndya king. I am coming to wait on thee with the sacred R\={a}m\={a}yanam, to obtain thy favour. I shall explain its meaning as much as I know, and it is for thee to hear it.}
\end{quotation}
\paragraph{}
This epistle having been favourably received by the \emph{Mudeliar}, he had the satisfaction of reciting his work before him, and the following impromptu, which he uttered on the occasion, commemorates the high honours, and munificent presents bestowed on him by this distinguished individual.
\begin{quotation}\small
கனந்தந்தான் கனகாபிஷேகந்தந்தான்‌

களங்கமில்லாக்கருப்பொருளையழைத்துத்தந்தான்‌

மனந்தந்தான்‌ முடிசூட்டுமாலைதந்தான்‌

வானிசிங்காதானத்திருத்தி வரிசைதந்தான்‌

இனந்தந்தான்‌ ராமகதையெவர்க்குந்தந்தான்

எனைராமாயணக்கவிஞனெனப்பேர்தந்தான்‌

அனந்தந்தான்‌ மணலிமுத்துக்கிஷணபூபன்‌

அகந்தந்தானிருமையிலுஞ்‌ சுகந்தந்தானே.
\end{quotation}
\begin{quotation}\small
\emph{He gave me gold, and heaped riches upon me;}

\emph{He gave me the products of the different countries;}

\emph{He gave me encouragement, he gave me a wreath to be worn on the head;}

\emph{He gave me honour, by seating me on the poetic throne;}

\emph{He gave me his friendship, he gave the history of R\={a}ma to all;}

\emph{He gave me a name as the R\={a}m\={a}yanic poet;}

\emph{He gave me food, he gave me a house;}

\emph{He gave me pleasure both even in the present and future life,}

\emph{So gave he Manali Muttukishna, who cherishes the world.}
\end{quotation}
\paragraph{}
Sometime afterwards \textbf{Tulaji Mahar\={a}ja} having become reconciled with the \textbf{Nawab}, Arun\={a}chala Kavir\={a}yer found admittance into the king's presence, and had the honour of reciting the drama before a full durbar, and was handsomely rewarded by the king.

It appears that in the latter years of his life a gloom having overspread his mind, he assumed the garb of a \emph{Saiva Pand\={a}ram}, and lived in retirement at Sh\={i}yally, where he died in \textsc{ce} 1779, at the age of sixty-seven.

\section{Arunaghiri N\={a}ther. --- அருணகிரிநாதர்.}
\paragraph{}
Arunaghiri N\={a}ther was an eminent poet, who from his entire devotedness to the worship of \emph{Siva}, renounced the world and assuming the life of an ascetic, lived and died at the \emph{Siva} fane at Trinomali, in the Carnatic. His compositions are of a religious character, and in praise of \emph{Skanda}, the Hindu Mars. They are as follows: \emph{Kanden Alangk\={a}ram} (கந்தனலங்காரம்), \emph{Kanden Anubh\={u}di} (கந்தனனுபூதி), \emph{Kanden Ant\={a}di} (கந்தனந்தாதி), \emph{Tiruvaguppu} (திருவகுப்பு), and \emph{Tiruppugazh} (திருப்புகழ்). This last work containing a series of 100 stanzas, being enriched by a splendid and harmonious flow of diction, has gained great popularity amongst the \emph{Saivas}, who generally have it recited, when any member of their sect is at the point of death. Arunaghiri N\={a}ther was also the author of a short harmonic poem, entitled \emph{Udatk\={u}tu Vannam} (உடற்கூற்றுவண்ணம்), descriptive of the different stages of human life, which together with an abridged translation of it in English, by the Rev. Mr. \textbf{Roberts},\footnote{Roberts Oriental Illustrations of the Sacred Scriptures, p. 348.} as revised by us we here subjoin:
\begin{quotation}\small
உடறகூற்றுவண்ணம்‌.

ஒருமடமாது மொருவனுமாகி யின்பசுகந்தரு மன்புபொருந்தி, உணர்வுகலங்கி யொழுகியவிந்து வூறுசுறோணித மீதுகலந்து, பனியிலோர்பாதி சிறுதுளிமாது, பண்டியில்வந்து புகுந்துதிரண்டு, பதுமவரும்பு கமடமிதென்று பார்வை மெய்‌ வாய்‌ செவி கால்கைகளென்‌ற, உருவமுமாகி யுயிர்‌வளர்மாத மொன்பதுமொன்று நிறைந்துமடந்தை, உதரமகன்று புவியில்விழுந்து யோகமும்‌ வாரமும்‌ நாளுமறிந்து, மகளிர்கள்‌ சேனை தரவணையாடை மண்படவுந்தி யுதைந்துகவிழ்‌ந்து, மடமயில்கொங்கை யமுதமருந்தி, யோரறிவீரறிவாகிவளர்ந்து, ஒளிநகையூறலிதழ்மடவாரு வந்துமுகந்திட வந்துதவழ்ந்து, மடியிலிருந்து மழலைமொழிந்து வாவிருபோவென நாமம்விளம்ப, உடைமணியாடை யரைவடமாட வுண்பவர்‌ தின்பவர்‌ தங்களோடுண்டு, தெருவிலிருந்து புழுதியளைந்து தேடியபாலரோடோடிநடந்து, அஞ்சுவயதாகிவிளையாடியே.

உயர்தருஞான குருவுபதேச முந்தமிழின்கலையுங்கரைகண்டு, வளர்பிறையென்று பலரும்விளம்ப வாழ்பதினாறுபிராயமும்வந்து, மயிர்‌முடிகோதியறுபதநீல வண்டிமிர்தண்டொடை கொண்டைபுனைந்து, மணிபொனிலங்கு பணிகளணிந்து, மாகதர்போகதர்‌ கூடிவணங்க, மதனசொரூபனிவனெனமோக மங்கையர்கண்டு மருண்டுதிரண்டு, வரிவிழிகொண்டு சுழியவெறிந்து, மாமயில்போலவர்‌போவதுகண்டு, மனதுபொறாமலவர்பிறகோடி மங்கலசெங்கலசந்திகள்‌கொங்கை, மருவமயங்கியிதழமுதுண்டு தேடியமாமுதல்‌ சேரவழங்‌கி, ஒருமுதலாகி முதுபொருளாயிருந்ததனங்களும்‌ வம்பிலிழந்து, மதனசுகந்தவிதனமிதென்று வாலிபகோலமும்‌ வெறுபிரிந்து, வளமையுமாறி யிளமையுமாறி வன்பல்விழுந்திரு கண்களிருண்டு, வயதுமுதிர்ந்து நரைதிரைவந்து வாதவிரோத குரோதமடைந்து, செங்‌கையினிலோர்தடியுமாகியே.

வருவதுபோவதொருமுதுகூனு மந்தியெனும்படி குந்திநடந்து, மதியுமழிந்து செவிதிமிர்வந்து வாயறியாமல்‌ விடாமன்மொழிந்து, துயில்வருநேரமிருமல்பொறாது தொண்டையுநெஞ்சுமுலர்‌ந்துவரண்டு, துகிலுமிழந்து சுணையுமழிந்து, தோகையர்பாலர்கள்‌ கோறணிகொண்டு, கலியுகமீதிலிவர்மரியாதை கண்டிடுமென்பவர்‌ சஞ்சலமிஞ்ச, கலகலவென்று மலசலம்வந்து கால்வழிமேல்வழி சாரநடந்து, தெளிவுமிராமலுரைதடுமாறி சிந்தையுநெஞ்சுமுலைந்துமருண்டு, திடமுமலைந்து மிகவுமலைந்து தேறினலாதரவேதெனநொந்து, மறையவன்வேதனெழுதியவாறு வந்ததுகண்டமுமென்றுதெளிந்து, இனியெனகண்ட மினியெனதொந்த, மேதினிவாழ்வுமிலாதினிநின்ற, கடன்‌முறைபேசுமெனவுரைநாவுறங்கிவிழுந்துகைகொண்டுமருந்து, ௧டைவிழிகஞ்சியொழுகிடவந்து பூதமுநாலு சுவாசமுநின்று, நெஞ்‌சிதடுமாறிவருநேரமே.

வளர்பிறைபோலவெயிறுமுரோம முஞ்சடையுஞ்சிறுகுஞ்சியும்விஞ்சு, மனதுமிருண்ட வடிவுமிலங்க, மாமலைபோல்‌ யமதூதர்கள்‌வந்துவலைகொடுவீசி வுயிர்கொடுபோக மைக்தரும்‌உந்து குனிந்தழநொந்து, மடியில்விழுந்து மனைவிபுலம்‌ப மாழ்கினரேயிவர்காலமறிந்து, பழையவர்காணுமெனுமயலோர்கள்‌ பஞ்சுபறந்திட நின்றவர்பந்தர்‌, இடுமெனவந்து பறையிடமுந்த வேபிணம்‌வேக விசாரியுமென்று, பலரையுமேவி முதியவர்தாமிருந்தசவங்கழுவுஞ்சிலரென்று. பணிதுகிறொங்கல்‌ களபமணிந்து பாவகமேசெய்து நாறுமுட ம்பை, வரிசைகெடாமலெடுமெனவோடி, வந்திளமைந்தர்‌ குனிந்துசுமந்து, கடுகிநடந்து சுடலையடைந்து மானிடவாழ்வெனவாழ்‌வெனநொந்து, விறகிடைமூடி யழல்‌கொடுபோட வெந்துவிழுந்து முறிந்துநிணங்க, ளுருகியெலும்பு கருகியடங்கி யோர்பிடிநீறுமிலாதவுடம்பை, நம்புமடியேனையினியாளுமே.
\end{quotation}
\begin{quotation}\small
\emph{The beginning of man is as a dew-drop falling from the tip of a blade of grass: he assumes his corporeal form in the womb of his mother, in the course of ten months and is then brought forth; he lies down, crawls, prattles, walks, and becomes acquainted with science. At sixteen he is in the bloom of youth; goes forth richly dressed, and adorned with jewels, courts the society of young women, is ravished by their eyes, lets himself loose to all irregularities and squanders his wealth. His youth now passes away and old age creeps on; his hair turns grey, his teeth drop, his eyes grow dim, his organs of hearing are impaired, his body becomes dry, his back bends, and he is troubled with incessant cough, which disturbs his sleep. He is no longer able to walk, and has recourse to a cane, and is gazed at by the young with derision. While in this condition, the ministers of Yama (Pluto), with their shaggy hair and frightful countenance, approach him, and seize and bear away his life amidst the screams of his wife and children. His kinsfolk and friends then assembled, talk of his good or evil deeds, convey his body to the burning ground with the sound of tom-toms, and commit it to the flame which consumes and reduces it to a handful of ashes.}
\end{quotation}
\section{Asht\={a}vad\={a}ni. --- அஷடாவதானி.}
\paragraph{}
The real name of this poet is not known, and the one under which he is here noticed being merely titular, alluding to the extraordinary retentiveness of his memory. He was born of a wealthy \emph{Br\={a}hman} family of the \emph{Vaishnava} sect long settled at \emph{Sirangam}, near Trichinopoly, and having finished his course of studies in Sanskrit and Tamil before he had attained his sixteenth year, he married and established himself at Tirupati; but quarreling with his wife he left home and set out on a pilgrimage to the \emph{Saiva} fanes in the south. When at Madura, he was so unfortunate as to fall into the snares of a courtesan, who after having stript him of everything by her sweet blandishments, turned him out destitute into the streets. Grieved at his losses, and wiser in his misfortune, he wrote a satirical poem, under the title of \emph{Viralividud\={u}tu} (விறலிவிடுதூது), exposing the wiles of the courtesan and warning the youth against them, and dedicated it to \textbf{N\={a}gama K\={u}lappa N\={a}yaken}, Poligar of Neleikkottei in the Dindigal district, by whom he appears to have been patronized in the later period of his life. This poem consists of 1144 couplets, and is not undeserving of the high reputation, which it has obtained for the beauty of its versification; but the author has completely demolished its moral tendency by the licentious description of the personal charms of the courtesan, and her enticing deportment.

\section{Athigam\={a}n. --- அதிகமான்.}
\paragraph{}
Athigam\={a}n was one of the brothers of \textbf{Tiruvalluver}, and having been exposed by his mother soon after his birth, in a wood at \emph{Karuv\={u}r}, he was found and brought up by a \emph{Chera} king. He is represented to have enjoyed great reputation as a poet, and the patron of poets, and to have composed many poems on various subjects, but none of his compositions is now forthcoming. According to some accounts, on the death of his foster-father he succeeded him on the throne, and that he was the same with \textbf{Cheram\={a}n Perum\={a}n}, who patronized the minstrel \textbf{Panapatirer}, but this is manifestly untrue; for \textbf{Cheram\={a}n Perum\={a}n} was not a foundling, but actually the son of \textbf{Sankaran\={a}tha Cheran}, and was crowned in the Kali year 3538 (\textsc{ce} 437\footnote{Madras Journal of Literature and Science, vol. 7. p. 284.}) which is four centuries anterior to the epoch assigned to \textbf{Tiruvalluver}.\footnote{Asiatic Researches, vol. 7. p. 349.}

\section{Athiv\={i}ra R\={a}ma P\={a}ndiyen. --- அதிவீரராமபாண்டியன்.}
\paragraph{}
Athiv\={i}ra R\={a}ma P\={a}ndiyen was one of the kings of the \emph{P\={a}ndiyen} dynasty, who reigned at Madura between the eleventh and twelfth century of the Christian era. Though the history of his life is not now forthcoming, yet his name has been handed down to posterity as that of an eminent poet, by the numerous poems which he composed. His principal work is the \emph{Naidatham} (நைடதம்), an imitation in Tamil of \textbf{Sriharsha's} \emph{Naishadham}, reciting the adventures of \emph{N\={a}lah} and \emph{Damayanti}. It comprises 1171 stanzas arranged in 29 cantos, and Mr. \textbf{Ellis} observes, that ``its high and courteous tone, notwithstanding the frequent occurrence of those gaudy images and far-sought allusions which European taste will denominate conceits, is worthy of the princely author.''\footnote{Ellis' Commentary on the Kural, p. 163.} It is so highly valued by the Tamils for the sweetness of its poetry that they call it ``the nectar of poets,'' but it should be observed that the extreme passion of the author for hyperbole has led him to extravagance, as will be seen by the following stanza in which he celebrates the beauty of \emph{Damayantis's} face.
\begin{quotation}\small
ஒழுகொளிவிரிந்த கதிர்மதிவண்ணனுந்தியந்தாமரையுறைவோன்‌ --- முழுமதிக்கலையுணிறைந்தபேரழகை மொண்டுகொண்டரிபரந்தகன்ற --- மழைமதர்நெடுங்கட்டமயந்திவதனம்‌ வகுத்தனன்கரத்தினாலெடுத்த --- புழைவழித்தோன்றும்‌ விசும்பினையன்றோகளங்கமென்றின்னமும்புகல்வார.
\end{quotation}
\begin{quotation}\small
\emph{He (Brahm\={a}) who dwells in the lotus-flower, blooming from the navel of the bright sapphire coloured (Vishnu), formed the face of the long-eyed Damayanti by taking his materials from the great beauty which is concentrated in the full moon; and it is the blue sky, which is seen through the aperture made in the moon for this purpose, that people still call the macula.}
\end{quotation}
\paragraph{}
His other works are the imitations of the \emph{K\={a}s\={i}k\={a}ndam} (காசீகாண்டம்), \emph{Linga Pur\={a}nam} (இலிங்கபுராணம்), and \emph{K\={u}rma Pur\={a}nam} (கூர்மபுராணம்), besides four \emph{Antadis},\footnote{A sort of poem, in which the last word of a stanza is repeated in the beginning of the next.} each comprising a centum of verses in praise of the \emph{Saiva} shrine at Karuvei and a collection of sententious maxims, under the title of \emph{Vettiv\={e}tkai} (வெற்றிவேற்கை), of which the Rev. Mr. \textbf{Taylor} has given an English version in the second volume of his Oriental Historical Manuscripts. Athiv\={i}ra R\={a}ma P\={a}ndiyen is also said to have been the writer of the poem, entitled \emph{Madana N\={u}l} (மதனநூல்); but we doubt whether he did really compose it. Its materials appear to have been derived from a work on the same subject in Sanskrit, called \emph{Kokva S\={a}stra}, and though its poetry is good yet it is very exceptionable in point of morality and the tenor of it shews that the pleasure of the fair sex was the uppermost thought in the mind of the author. It is the same work, which \textbf{Abbe Dubois}, in his Description of the People of India, Chap. 14., notices, and justly condemns, as abominable; yet we are sorry to state, that some Hindus at Madras have been licentious enough to class it with their \emph{\={A}gamas} under the title of \emph{K\={a}m\={a}gamam}, and to prostitute the press to its publication.

\section{Auveiy\={a}r. --- ஔவையார்.}
\paragraph{}
Auveiy\={a}r was the first of her sex in the Tamil country, who acquired a reputation as a poetess. The name Auveiy\={a}r is merely titular, signifying ``\emph{the elderly lady}'': what may have been her real name is not known. She was one of the sisters of \textbf{Tiruvalluver}, and her mother, who gave birth to her in a Choultry near \emph{Ureiy\={u}r}, having exposed her there, she was found and brought up by a \emph{P\={a}nan} or minstrel. How, and in what manner, she acquired her prodigious learning cannot be ascertained, but it is very probable she was in some measure indebted for it to her foster-father himself. Though she taught others that marriage was preferable to celibacy, yet she continued unmarried all her life. It is stated that by the virtue of a panacea, which she had discovered, she prolonged her existence to the period of 240 years, and then, according to the custom of her times, made the \emph{Mah\={a}prasth\={a}nagamana}, or great journey to the \emph{Imaus} and died there.

``She sang like \textbf{Sappho}; yet not of love, but of virtue.''\footnote{Calcutta Quarterly Review, vol. 25. p. 190.} Her two books of moral aphorisms usually called \emph{\={A}ttish\={u}di} (ஆத்திசூடி) and \emph{Kondei V\={e}ynden} (கொன்றைவேந்தன்) have been considered by the learned \textbf{Beschi} as being ``worthy of \textbf{Seneca} himself''\footnote{Babington's Shen Tamil Grammar, Introduction, p. 11.}; for they are as the Rev. Mr. \textbf{Percival} has justly remarked, ``of great beauty and value, replete with lessons of wisdom,'' and ``they have never been surpassed for sententious brevity; and generally they are equally distinguished by purity of principle.''\footnote{Percival's Land of the Veda, p. 100.} In both the books the aphorisms are written in the order of the Tamil alphabet; hence they have been most appropriately called ``the Golden Alphabet of the Tamils.''\footnote{Asiatic Researches, vol. 7. p. 350. London Edition.} No less than three English versions of Auveiyar's aphorisms have been published, the first by the Rev. Dr. \textbf{John} in the Asiatic Researches, the second by the Rev. Mr. \textbf{Bisset}, the author of the History of Ceylon under the name of Philalethes and the third by the Rev. Mr. \textbf{Sudgen}, and there is also a version of them in German, and another in Dutch, but it is hardly necessary to remark that they all merely convey the sense of the original, and that too often not very correctly. The other works of Auveiyar are the \emph{M\={u}durei} (மூதுரை), \emph{Nalvazhi} (நல்வழி), \emph{Kalviozhukam} (கல்வியொழுக்கம்), \emph{Asadik\={o}vei} (அசதிக்கோவை), \emph{Nann\={u}tk\={o}vei} (நன்னூற்கோவை), \emph{N\={a}nmanik\={o}vei} (நான்மணிக்கோவை), \emph{Aruntamilmalei} (அருந்தமிழ்மாலை), \emph{Tarisanapattu} (தரிசனப்பத்து), \emph{Gn\={a}na Kural} (ஞானக்குறள்), and \emph{Pandenant\={a}di} (பந்தனந்தாதி), and of these the first eight are on ethics, and religion, the ninth on metaphysics, and the tenth is a panegyric on a wealthy merchant of the name of \textbf{Panden}, who resided at Kaveripatnam. There is also extant under her name a \emph{Nigundu}, or dictionary in verse, containing the Tamil synonyma of the materia medica, but it is not in general use.

Auveiyar was more keen and clever than even her brother \textbf{Tiruvalluver}. The Hindu system of Moral philosophy involves four questions, viz: 1. What is \emph{virtue}, and how is it to be exercised? 2. What is \emph{Wealth}, and how is it to be acquired? 3. What is \emph{pleasure}, and how is it to be enjoyed? 4. What is \emph{heavenly bliss} and how is it to be obtained?\footnote{See Rottler's Tamil and English Dictionary under the word கேள்வி.} \emph{Tiruvalluver} however, in his \emph{Kural}, containing 1333 couplets, treated only of the three first questions, and when Auveiyar heard of it she compressed all the four questions, and their answers, within the compass of the following single stanza, which she uttered extempore:---
\begin{quotation}\small
ஈதலறந்தீவினைவிட்டீட்டல்பொருளெஞ்ஞான்றுங்‌

காதலிருவர்‌ கருத்தொக்க --- ஆதரவு

பட்டதேயின்பம்‌ பரனைநினைந்திம்மூன்றும்‌

விட்டதே பேரின்பலீடு.
\end{quotation}
\paragraph{}
\emph{To give is Virtue. That which is acquired without sin, is Wealth. The constant mutual affection of two, whose tastes agree, is Pleasure. To forsake these three, in the contemplation of the Supreme Being, is Heavenly Bliss.}

When \textbf{Tiruvalluver} went to the Madura College to submit the \emph{Kural} to its judgment, Auveiyar accompanied him thither, and making some signs to the professors, asked them what she meant by those signs upon which one of them improvised the following stanza, in reply:---
\begin{quotation}\small
இவ்வளவுகண்ணினாளிவ்‌வளவுசிற்றிடையா

ளிவ்‌வளவுபோன்றவிளமுலையா --- ளிவ்வளவாய்க்

காமத்தலைவனையுங்‌ காணாதுகன்றினா

ணாமத்தைக்காட்டுகுறிநன்கு.
\end{quotation}
\paragraph{}
Not satisfied with the explanation, she interpreted the signs herself, thus:
\begin{quotation}\small
ஐயமிடுமினறநெறியைக் கைப்பிடுமி

னிவ்வளவேனும்‌ மனத்தையிட்டுண்மின்‌ --- றெய்வ

மொருவனேயென்னவுணரவல்லீரே

லருவினைகளைந்துமறும்‌.
\end{quotation}
\begin{quotation}\small
\emph{Give alms, hold fast the true virtue,}

\emph{Eat rice, after giving to others even as much as a pinch.}

\emph{If you are able to know that God is but one,}

\emph{Then the five great sins will be destroyed.}
\end{quotation}
\paragraph{}
After the \emph{Kural} had passed the College bench, and each of the professors had pronounced a stanza on its merits, Auveiyar hearing \emph{Ideik\={a}der} praise it in a couplet, likening it to a mustard seed hollowed and filled with the waters of the seven seas, in as much as it contained much knowledge in a little compass, she observed to him, that it would have been more appropriate to liken it to an atom, which is even minuter than a mustard seed.

Auveiyar had the nickname of K\={u}zhukup\={a}di (கூழுக்குப்பாடி), or she who sung her verse for a cup of pap, and the reason assigned for it is, that a dancing girl, named \textbf{Chilambi}, who lived at Ambel, asked the poet \textbf{Kamber} to compose a verse on her, giving him 500 \emph{pons} for it, but as he never composed a stanza on anyone for less than 1000 \emph{pons}, he scribbled these first two lines of a stanza on the wall of her house with charcoal and went away.
\begin{quotation}\small
தண்ணீருங்காவேரி தார்வேந்தரிற்சோழன்

மண்ணாவதுஞ்சோழ மண்டலமே --- பெண்ணாவாள்
\end{quotation}
\paragraph{}
Immediately afterwards Auveiyar happening to call at the dancing girl's house and being informed of the affair, added the two following lines to the unfinished stanza of \textbf{Kamber}, accepting as a reward merely a cup of pap.
\begin{quotation}\small
அம்பற்சிலம்பியரவிந்தத்தாளணியுஞ்

செம்பொற்சிலம்பேசிலம்பு.
\end{quotation}
\paragraph{}
The following is the translation of the above stanza:---
\begin{quotation}\small
\emph{The best of rivers is the} K\={a}veri, \emph{the best of kings is} Chola,

\emph{The best of lands is Cholamandalam, the best of Chilambus,}\footnote{A tinkling foot ornament, worn by Tamil women.}

\emph{Are the gold Chilambus on the lotus feet,}

\emph{Of the damsel Chilambi of Ambel.}
\end{quotation}
\paragraph{}
On one occasion when some men reviled the character of women in Auveiyar's hearing, she is said to have retorted upon them with the following impromptu:
\begin{quotation}\small
நல்லார்களெல்லாரு நல்லாரேதன்மையால்‌

வல்லாராற்கேடு படாராயின்‌ --- நல்லறி

வாண்மக்கட்பற்பலர்க்கேயுண்டாகும்‌ பெண்டீரும்‌

மாண்புக்கெடுக்காவிடின்‌.
\end{quotation}
\begin{quotation}\small
\emph{All women would be good by nature, if men did not spoil them; and most men would have a tolerable stock of sense, if the women did 
not make fools of them.} \textbf{Ellis}.
\end{quotation}
\paragraph{}
The subjoined stanza eulogizing the city of \emph{K\={a}njipuram} is attributed to Auveiyar:---
\begin{quotation}\small
வையகமெல்லாங்‌ கழனியாய்‌ வையகத்துட்‌

செய்யகமேநாற்றிசையின்‌ றேயங்கள்‌ --- செய்யகத்துள்‌

வான்கரும்பேதொண்டை வழநாடுவான்கரும்பின்‌

சாறேயந்நாட்டிற்றலையூர்கள்சாறட்ட

கட்டியேகச்சிப்புறமெல்லாங்கட்டியுட்‌

டானெற்றமான்‌ சருக்கரைமாமணியே

யானேற்றான்‌ கச்சியகம்‌.
\end{quotation}
\begin{quotation}\small
\emph{The whole earth may be compared to an expanse of wet-land:}

\emph{The several countries of the earth, marked by their boundaries on the four sides, are cultivated fields within it;}

\emph{Tondei-vala-nadu is a lofty sugar-cane in one of these fields;}

\emph{The chief towns of this country are the crude juice of this sugar-cane;}

\emph{Cachchipuram (Conjiveram) and its vicinity are like a cake of unrefined sugar obtained by boiling this juice;}

\emph{And a large concretion of refined sugar in the midst of this cake}

\emph{Represents the interior of Cachchi, where the bull-borne Deity resides.}

\hspace*{15mm}\textbf{Ellis}.
\end{quotation}
\paragraph{}
\section{Avinayan\={a}r. --- அவிநயனார்.}
\paragraph{}
Little is known of the life of Avinayan\={a}r, except that he was one of the twelve disciples of \textbf{Agastiyer} and the author of a treatise on versification, which is called after his own name, \emph{Avinayam} (அவிநயம்). Father \textbf{Beschi}, in his Introduction to the Shen Tamil Grammar, notices Avinayan\={a}r as one of the persons, who wrote on the Tamil Grammar after \textbf{Agastiyer's} time and whose works have perished.

\section{Avir\={a}mipatter. --- அவிராமிப்பட்டர்.}
\paragraph{}
Avir\={a}mipatter, corrigo, \emph{Abir\={a}mi Bhatta}, was a poet equally well versed in Sanskrit and Tamil literature. He belonged to a class of \emph{Brahmans} at \emph{Tirukadav\={u}r} in the Carnatic, and was a staunch adherent of the sect which worship \emph{P\={a}rvati}, the female energy of \emph{Siva}, under the name of \emph{Abir\={a}mi}. Of his poetical compositions we have only a centum of hymns, addressed to \emph{Abir\={a}mi}, and which are remarkably harmonious in their versification. Such is the faith that \emph{Saivas} have in the efficacy of these hymns, that it is not unusual for them to commit them to memory, and recite them privately with great veneration. The following is one of the hymns referred to:
\begin{quotation}\small
மணியேமணியினொளியேயொளிருமணிபுனைந்த

வணியேயணியுமணிக்கழகேயணுகாதவர்க்குப்‌

பிணியேபிணிக்குமருந்தேயமரர்‌ பெருவிருந்தே

பணியேனொருவரை நின்பத்மபாதம்‌ பணிந்தபின்னே.
\end{quotation}
\begin{quotation}\small
\emph{Thou art the gem, thou art the brightness of that gem; thou art the necklace formed of that gem, and it is thou that givest it its beauty. Thou art the disease of those who do not approach thee, and the cure of those who approach thee; and thou art the great feast of the celestials. I shall not adore the feet of any after I have adored thy lotus-feet.}
\end{quotation}

\section{\={A}lankudi Vangan\={a}r. --- ஆலங்குடிவங்கனார்.}
\paragraph{}
Nothing is known of this poet, except that he was a native of \emph{\={A}lankudi} and that his high attainments had obtained for him a place in the Madura College as one among the forty-nine professors. His works have all perished, and even his own name would have long ago fallen into oblivion, but for the following impromptu which he uttered when the \emph{Kural} passed the College bench.
\begin{quotation}\small
வள்ளுவர்பாட்டின்‌ வளமுரைக்கின்‌ வாய்மடுக்குந்

தெள்ளமுதின்றீஞ்சுவையுமொவ்வாதாற்‌ --- றெள்ளமுத

முண்டறிவார்தேவருலகடையவுண்ணுமால்‌

வண்டமிழின்முப்பான்மகிழ்ந்து.
\end{quotation}
\begin{quotation}\small
\emph{Even the sweetness of ambrosia cannot be compared to the sweetness of Valluver's poem. The gods have known the taste of ambrosia by having partaken of it, but men will know it when they imbibe the milk (instruction), which issues from the three udders (the three parts) of the Kural.}
\end{quotation}

\section{\={A}lavand\={a}r. --- ஆளவந்தார்.}
\paragraph{}
\={A}lavand\={a}r was both a poet and philosopher. He belonged to a class of \emph{Br\={a}hmans}, styled \emph{Bhattas}, and traced his pedigree to \textbf{M\={a}dhava Bhatta}, an eminent \emph{Guru}, who lived at \emph{V\={i}rei} in the Carnatic, in ancient times. His poetical version of the \emph{Gn\={a}na V\={a}shistha} from Sanskrit into Tamil in 2,055 stanzas is an elaborate performance, and greatly redounds to his credit. The \emph{Gn\={a}na V\={a}shitha} treats of the \emph{Vedic} philosophy, and is interspersed with legends of various real or fabulous personages. The following stanza which we quote from the work contains a description of the attributes of God:
\begin{quotation}\small
அத்தமியாதசோதியருக்கனாயளவிலாதாய்‌

மெத்தியபிணியிலாதாய்‌ மேவரும்பிறப்பிலாதா

யெத்தலங்களினுமென்றுமியாவுமாயாவுஞ்செய்து

முத்தியுமாகியோங்கு முதற்பரம்பொருளொன்றுண்டு.

\hspace*{15mm}ஞானவிண்மகன்கதை 𑌨𑍍𑌧 ம்‌ பாட்டு.
\end{quotation}
\begin{quotation}\small
\emph{There is a supreme being, --- he is the sun who shines with unreflected light, without setting, --- he who is infinite, --- he who is free from infirmity, --- he who is not subject to birth, --- he who is everywhere and at all times, --- he who is all and the creator of all, and he who is the final beatitude itself.}
\end{quotation}

\section{\={A}ndi Pulavar. --- ஆண்டிப்புலவர்.}
\paragraph{}
This poet, who distinguished himself by his eminent attainments in philology, was the son of \textbf{P\={a}v\={a}dei V\={a}dhy\={a}r} of \emph{\={U}tankal}, a village in the Gingee country. He had a great partiality for writing verses of the kind called \emph{\={A}s\={i}r\={i}yam}, and we have a commentary of his on the first two parts of the \emph{Nann\={u}l}, under the title of \emph{Ureiyari Nann\={u}l} (உரையறிநந்நூல்), and a \emph{Nigandu}, or dictionary of the Tamil synonyms, both written in the same metre.

\section{\={A}siriyer Nallanduvan\={a}r. --- ஆசிரியர்நல்லந்துவனார்.}
\paragraph{}
This poet was one of the forty-nine professors of the Madura College. He had the title of \emph{\={A}siriyer} or doctor, conferred on him as an honorary mark of literary distinction. Neither his parentage nor personal history is known, and even of his compositions we have nothing more than the following impromptu in which he expressed his approbation of the \emph{Kural}, and eulogized its author when that work was submitted to the criticism of the College bench.
\begin{quotation}\small
சாற்றியபலகலையுந் தப்பாவருமறையும்‌

போற்றியுரைத்தபொருளெல்லாந் --- தோற்றவே

முப்பான்மொழிந்த முதற்பாவலரொப்பா

ரெப்பாவலரினுமில்‌.
\end{quotation}
\begin{quotation}\small
\emph{There is none amongst the learned that can compete with the first of poets (meaning Tiruvalluver) who composed the Kural, which has in itself all that the works of the sages and the infallible Vedas have preserved.}
\end{quotation}

\section{Chatteimuni. --- சட்டைமுனி.}
\paragraph{}
Nothing is known of Chatteimuni, except that he was a \emph{Saiva} ascetic, who lived at the same time with \textbf{Agastiyer}, and enjoyed some reputation as a poet, philosopher and physician. He wrote a work on theology, called \emph{Gn\={a}na N\={u}ru} (ஞானநூறு), and another on medicine, called \emph{Kalpa N\={u}ru} (கற்பநூறு). He is also believed to have been the author of the \emph{V\={a}da Nigandu} (வாதநிகண்டு), a treatise on alchemy which passes under his name.

\section{Ch\={e}ram\={a}n Perum\={a}n. --- சேரமான்பெருமான்.}
\paragraph{}
Ch\={e}ram\={a}n Perum\={a}n, also called \emph{Perum\={a}kodeiyar}, was both a king and poet. Some say that he was the son of \textbf{Sankaran\={a}tha Chera} and others of \textbf{Senkotporei Chera}, and in the \emph{Cholap\={u}rva Patayam}, one of the Mackenzie Mss. the date of his accession is placed in the \emph{Kali} year 3538, (\textsc{ce} 437). He became an ascetic while yet a youth, and retiring from Cranganore, where his father held his court established himself in the \emph{Saiva} temple at Tiruvanjikalam, but on the abdication of the throne by his father, he returned to the capital and assumed the reins of government upon the importunity of the ministers. He appears to have endeavoured to advance the interests of literature, by patronizing the poets, who resorted to his court; but being a superstitious votary of \emph{Siva}, spent most of his time in making pilgrimages. He was a friend of \textbf{Sundarer}, and died simultaneously with him as may be inferred from the fable of their repairing to \emph{Kail\={a}sam} together and remaining there in a state of beatification. He composed three poems in honour of \emph{Siva}, namely, the \emph{Mummani-k\={o}vei} (மும்மணிக்கோவை), \emph{Kail\={a}savul\={a}} (கைலாசவுலா), and \emph{Ponvannatand\={a}di} (பொன்வண்ணத்தந்தாதி) and they are still in much repute.

\section{Deivasih\={a}mani Pillei. --- தெய்வசிகாமணிப்பிள்ளை.}
\paragraph{}
Deivasih\={a}mani Pillei, was head \emph{Sirashtadar} of the Hussoor Kutcherry at Calicut. He was born of a \emph{Vel\={a}la} family at \emph{Dindigal}, in the Carnatic, on the 10th of April, \textsc{ce} 1804, and his father \textbf{Anthonimuttu Ubadesi}, a Christian, took every care to secure him a sound religious education. He had from his early years a predilection for poetry, which he afterwards cultivated with great success. In 1842 he composed two beautiful poems in praise of our Lord, respectively entitled \emph{Inbamanimalei} (இன்பமணிமாலை), and \emph{St\={o}trokummi} (ஸ்தோத்திரக்கும்மி), and they were both inserted in the \emph{Rising Sun}, a monthly Tamil periodical then edited at Colombo by ourselves. His poetical career was however not very long, for it was cut short by death on the 30th of November 1846.

\section{Ellappa N\={a}valer. --- எல்லப்பநாவலர்.}
\paragraph{}
Ellappa N\={a}valer was well versed in the Sanskrit and Tamil languages, and well read in the classics of both. His poetry was distinguished by superior excellence, as his imitation of \textbf{Vy\={a}s\={a}'s} \emph{Arun\={a}chala Mah\={a}tmya} (the legends of the \emph{Saiva} shrine at Trinomali in Tamil, under the title of \emph{Arun\={a}chala Pur\={a}nam} (அருணாசலபுராணம்), manifestly proves. This work comprises 586 stanzas, arranged in twelve cantos, and it has obtained an unusually wide circulation amongst the \emph{Saivas}. We quote the following three stanzas from the fifth canto (பாகம்பெற்றசருக்கம்), as a specimen, and which accounts for \emph{Siva} being represented as androgynous, or half man and half woman:---
\begin{quotation}\small
பிறப்‌பிலாமுதன்மறையெலாமளவிடாப்பெரியோ

னிறப்‌பிலாவிறைவிடையின்மேலினிதெழுந்தருள

வறத்தினாயகியடிபணிந்திறைஞ்சினளனைத்துந்

துறக்குமண்ணலுமவளுடனிவையிவைசொல்லும்‌.

பெண்ணுமாணுமாயுலகுயிர்ப்பொருளெலாம்பிரித்தோ

மண்ணிதாகிய பீடமுமிலிங்கமுமானோ

மெண்ணுநீயெனநாமெனவிரண்டிலைமரமு

முண்ணிலாவிய வயிரமுநிகர்த்துடனுறைவோம்‌.

அங்கம்யாவுநம்பொருட்டுவிட்டிமையவெற்படைந்தா

யிங்குநாமுனக்கொருபுறமளிப்பதேயியற்கை

மங்கையேநமதிடப்புறத்துறையெனமகிழ்வுற்‌

றங்கையாலணைத்தருளினனுருகியொன்றானார்‌.
\end{quotation}
\begin{quotation}\small
\emph{When the Eternal, who is beyond the compass of the various Vedas manifested himself riding on the bull, the virtuous (P\={a}rvati) casting herself at his feet, adored him: thereupon he addressed her thus; Come, O P\={a}rvati! I and thou becoming male and female produced all the living beings in the world, we became the Yoni and the Linga, there is no difference between us as I and thou; let us then be united together as the tree and the core.}

\emph{Having quitted the whole of thy body for my sake thou hast betaken thyself to the mountain Imaus; hence it is proper that I should yield to thee one half of my body. O Damsel! take thy place with me on my left. So saying he embraced her with his arms, and then their bodies melted and became united together.}
\end{quotation}

\section{Erichal\={u}r Mal\={a}dan\={a}r. --- எரிச்சலூர்மலாடனார்.}
\paragraph{}
Erichal\={u}r Mal\={a}dan\={a}r was a native of \emph{Erichal\={u}r}, and one of the forty-nine professors of the Madura College. His personal history like that of his colleagues has been lost, and the only relic of his poetical compositions is the following impromptu on the \emph{Kural}, which merely enumerates the divisions and subdivisions of that work.
\begin{quotation}\small
பாயிரநான்கில்லறமிருபான்பன்மூன்றே

தூயதுறவறமொன்றூழாக --- வாய

வறத்துப்பானால்வகையாவாய்ந்துரைத்தார்நூலின்‌

றிறத்துப்பால்‌வள்‌ளுவனார்தேர்‌ந்து.
\end{quotation}

\section{Ganapati Aiyer. --- கணபதி ஐயர்.}
\paragraph{}
Ganapati Aiyer was a \emph{Br\={a}hman} of \emph{Batticotta} in Jaffna, who enjoyed the reputation both of a poet and physician. The date of his birth is not known; but he died in \textsc{ce} 1803. He is chiefly known as the author of two dramas, the \emph{Alank\={a}rar\={u}pa N\={a}dakam} (அலங்காரரூபநாடகம்), and the \emph{V\={a}lab\={i}ma N\={a}dakam} (வாளபீமநாடகம்), which have gained much popularity and are often acted.

\section{Gn\={a}na K\={u}tur. --- ஞானக்கூத்தர்.}
\paragraph{}
Gn\={a}na K\={u}tur was a \emph{Saiva} ascetic and poet, who lived at \emph{Sivenpakam}, in the Carnatic; but of whose personal history scarcely any thing is known. He is however remembered as the author of the \emph{Vriddh\={a}chala Pur\={a}nam} (விருத்தாசலப்புராணம்), a poem consisting of 435 stanzas, arranged in eighteen cantos, and recounting the legends of the \emph{Saiva} shrine on the Vriddh\={a}chala Hill, in the Carnatic.

\section{Gn\={a}naprak\={a}sa D\={e}siker. --- ஞானப்பிரகாசதேசிகர்.}
\paragraph{}
Gn\={a}naprak\={a}sa Desiker was a native of \emph{Tinnevelly}, in Jaffna. Having gone over to Southern India in his early youth, and perfected himself there in the knowledge of the Tamil Grammar and Classics, under the tuition of several Pundits, he repaired to the Gauda country for the purpose of studying Sanskrit, in order to read the \emph{Saiva Agamas} in their original; but the \emph{Br\={a}hmans} of Gauda objected to it on the ground that he being a \emph{S\={u}dra} could not have access to those sacred books. One of them however, who was more liberal minded than the rest, perceiving the acuteness of his intellect, received him as his pupil, and in the course of a few years made him an able Sanskrit scholar. He afterwards returned to Southern India, and having been invested with the office of \emph{Tambir\={a}n}, established himself at Trinomali, and composed several works in Sanskrit for the use of the members of the \emph{Matam} at that place. Whilst he was thus employed at Trinomali, having heard some through envy remark that he was clever only in Sanskrit and not in Tamil, he wrote in Tamil a learned Commentary on the \emph{Sivagn\={a}na Sittiyar} which is now in great repute wherever the \emph{Saiva} religion prevails.

\section{Gunav\={i}ra Panditer. --- குணவீரபண்டிதர்.}
\paragraph{}
This poet was a native of \emph{Kalendai}, a town near Chingleput, and he attracted much notice in his day. He is known as the author of two treatises on the art of poetry, one of which styled \emph{Nemin\={a}dam} (நேமிநாதம்), consists of 96 stanzas, and the other styled \emph{Venb\={a}p\={a}ttiyel} (வெண்பாப்பாட்டியல்), consists of 100 stanzas. The last work is stated to have been composed under the patronage of \textbf{Vachchanandimuni}, an eminent ascetic, in the time of the king \textbf{Tribhuvana Deva}, perhaps the same with \textbf{Tribhuvana-malla}, who reigned at Anumakonda in the eleventh century of the Christian era.\footnote{Madras Journal of Literature and Science, vol. 10. p. 20.}

\section{Gurup\={a}dad\={a}sen. --- குருபாததாசன்.}
\paragraph{}
Nothing more is known of this poet than that he was the author of a poem, entitled \emph{Kumaresa Satagam} (குமரேசசதகம்), comprising a centum of stanzas on moral and social duties, and which has obtained great popularity. Each stanza in this work ends with an eulogistic address to \emph{Kumaresa}, the form of \emph{Skanda} worshipped at Tiruppulvaiyal, and we subjoin one of them as a specimen.
\begin{quotation}\small
தங்கமானது தழலினின்றுருகிமறுகினுந் தன்னொளிமழுங்கிடாது --- சந்தனக்குறடுதான்‌ மெலிந்துதேய்ந்தாலுமோ தன்மணங்குன்றிடாது --- பொங்கமிகுசங்குசெந்தழலில்வெந்தாலுமோ பொலிவெண்மைகுறைவுறாது --- போதவேகாய்ந்துநன்பால்குறுகினாலும்‌ பொருந்துசுவைபோய்விடாது --- துங்கமணிசாணையிற்றேய்ந்துவிட்டாலுந் துலங்குகுணமொளியாதுபின்‌ --- தொன்மைதருபெரியோர்‌ மடிந்தாலுமவர்களது தூயநிறைதவறாகுமோ --- மங்களகல்யாணிகுற மங்கைசுரகுஞ்சரியை மருவுதிண்புயவரசனே --- மயிலேறிவிளையாடு குகனேபுல்வயனீடு மலைமேவுகுமரேசனே.
\end{quotation}
\begin{quotation}\small
\emph{As gold though melted in the fire, yet will not lose its lustre; the sandal-wood though worn out, yet will not lose its fragrance; the chank though burnt, yet will not lose its whiteness; the milk though diminished in quantity by being boiled, yet will not lose its flavour; and the gem though ground, yet will not lose its brilliancy; so great men though impoverished, yet will not lose their magnanimity; O Kumaresa! the strong-shouldered king, who wedded the Kurava female and riding on the peacock, sports on the hill at Pulvaiyel.}
\end{quotation}
\section{Ideik\={a}der. --- இடைக்காடர்.}
\paragraph{}
Ideik\={a}der, so named after \emph{Ideikazhi-n\={a}du}, the land of his nativity, was a poet from his earliest youth. Having heard that \textbf{Kulesa P\={a}ndiyen}, who then reigned at Madura, added to the title of a king the still more enduring distinction of a poet, and had a seat on the College bench, Ideik\={a}der sought to secure his patronage by composing a panegyric on him, and reciting it in his presence; but the king took no notice of it. Stung to the quick at the neglect, he retired to a temple at Vada Madura, whither \textbf{Kapiler}, and the rest of the professors of the College, who did not approve of the king's conduct, followed him. The king on receiving this intelligence went to Vada Madura, and by much entreaty prevailed on him, and the professors to return to the College and having heard his poem again recited, honoured him by the gift of an elephant, and a horse, and by the donation of lands and gold. Satisfied with this he took leave of the king, and returned to his country. He again went to Madura to be present when his friend \textbf{Tiruvalluver} submitted the \emph{Kural} to the judgment of the College bench. After having heard the forty-nine professors each recite a stanza in praise of the \emph{Kural}, he expressed his own opinion of the merits of that work in the following distich:---
\begin{quotation}\small
கடுகைத்துளைத்தேழ்கடலைப்புகட்டிக்

குறுகத்தறித்தகுறள்.
\end{quotation}
\begin{quotation}\small
\emph{The Kural contains much in a little compass. Such is the ingenuity of its author, that he has compressed within its narrow limits all the branches of knowledge; as if he had hollowed a mustard seed, and inclosed all the waters of the seven seas in it.}
\end{quotation}
\paragraph{}
When the professors found that \textbf{Tiruvalluver} had readily answered the many puzzling questions which they proposed to him, they attempted to try Ideik\={a}der. He improvised the following stanza which they could not understand, and hence remained silent:---
\begin{quotation}\small
ஆற்றங்கரையினருகிருந்தமாமரத்திற்

காக்கையிருந்து ௧ஃகஃகெனக் --- காக்கைதனை

யெய்யக்கோலில்லாமலிச்சிச்செனவெய்தான்‌

வையக்கோனார்தன்மகன்‌.
\end{quotation}
\paragraph{}
Of Ideik\={a}der's works we have not met with any; but we find a grammatical treatise of his, under the title of \emph{Oosimuri} (ஊசிமுரி), quoted by \textbf{Gunas\={a}garer}, the Commentator of the \emph{K\={a}rigai}.

\section{Ilav\={e}ttan\={a}r. --- இளவேட்டனார்.}
\paragraph{}
This poet was one of the forty-nine professors of the Madura College. All that is known of his personal history is, that he was a native of Madura, and belonged to the caste of \emph{Aruvei-v\={a}nigars}, or linen-drapers. The following is the impromptu which he uttered in approval of the \emph{Kural}:
\begin{quotation}\small
இன்பமுந் துன்பமுமென்னுமிவையிரண்டு

மன்பதைக்கெல்லாமனமகிழ வன்பொழியா

துள்ளியுணரவுரைத்தாரேயோதுசீர்‌

வள்ளுவர்வாயுறைவாழ்த்து.
\end{quotation}
\begin{quotation}\small
\emph{Pleasure and pain are two different things (one constituting the whole of positive happiness and the other the whole of misery), and Valluver is to be praised for teaching all mankind in his Kural the means of securing the one, and avoiding the other.}
\end{quotation}

\section{Inbakavi. --- இன்பகவி.}
\paragraph{}
Inbakavi was a native of \emph{Manapar}, in the Tinnevelly District, and born of parents of the \emph{Parava} caste. His Christian name was \textbf{Xavier Henric Le\={a}m}. The juvenile productions of his muse, although not of the highest order, were nevertheless sufficient to secure him the patronage of \textbf{Don Gabriel Vaz Gomez}, the \emph{Parava} Chieftain of Tutikorin. Going once on a visit to the Raja of Ettiapore, the ambition of the poet was checked when he found that he had to stand, whilst the \emph{Samprati} or minister enjoyed a seat in the Raja's presence; but an impromptu stanza in which he praised the Raja, and respectfully claimed a seat, soon procured him the desired privilege, and also many presents. We subjoin the above mentioned stanza, and attempt to give an English translation of it, which we are sorry to say being stript of the melody and linked-sweetness of the original, appears but too trite in a foreign language.
\begin{quotation}\small
கங்கைகுலத்திற்‌ கனகமணியாயுதித்த

சங்‌கரக்குமார தயாநிதியே --- உங்களுடை

சம்பிரதிக்கந்தவிடந் தங்களுக்குமிந்தவிட

மின்பகவிக்கெங்கேயிடம்‌.
\end{quotation}
\begin{quotation}\small
\emph{O Shankara Kum\={a}ra, the bountiful, who has sprung up like a resplendent gem from the race of Ganga! That seat is for the Samprati, and this seat is for thee, but where is the seat for Inbakavi?}
\end{quotation}
\paragraph{}
Finding so much success attend him in his native District, Inbak\={a}vi resolved to visit the neighbouring Districts, and obtain rewards by singing the praises of the high personages there. He was well received wherever he went, and especially in Tanjore, where \textbf{Tattogee}, the minister of the king \textbf{Surfogee}, treated him with marked distinction by presenting him with a palanquin. He had been long while on his peregrination, when a severe attack of illness which threatened his life, obliged him to return to Tutikorin. On his presenting himself before the \emph{Parava} chieftain, and being asked where he had been so long absent, he answered in the following impromptu: 
\begin{quotation}\small
பெலத்ததொருவானவரை முனிவரரையுலகிலுள்ளபெரியோர்தம்‌மை --- யழைத்துவிடுமெமராசன்‌ முன்சென்றேனெனையும்பார்த்தெவர்காணென்றான்‌ --- தலத்திலுயர்‌ பரதசாதித்தலைவமோர்‌ வாசல்வித்வானென்றேன்‌ --- மலைத்தவனு முனதுநகர்க்கேகுமென்றான்‌ யானுமிதோ வந்திட்டேனே.
\end{quotation}
\begin{quotation}\small
\emph{Having appeared before Yama (the king of death), who summons before his presence the mighty celestials as well as the sages and magnates of this world, he asked me who I was, and I replied that I was the poet employed at the Gate of the eminent Parava Chieftain, and he thereupon was perplexed, and bade me go back to thy city, and behold I am here!}
\end{quotation}
\paragraph{}
Towards the latter part of his life, Inbakavi frequently visited Colombo, and once Jaffna. He then unfortunately indulged himself in spirituous liquors, and although this had in no measure impaired his intellectual powers, it tended in no small degree to diminish that respect, which was otherwise due to so clever a man. When his friends remonstrated with him on his vicious habit, he was wont to excuse himself almost in the same words with \textbf{Anacreon}:---
\begin{quotation}\small
\emph{When wine I quaff, before my eyes}

\emph{Dreams of poetic glory rise;}

\emph{And freshen'd by the goblet's dews,}

\emph{My soul invokes the heavenly Muse.}
\end{quotation}
\paragraph{}
We give a very well authenticated anecdote of Inbakavi, which will prove that he was not to be taken advantage of even in his weakest moments. A poet at Jaffna happening to meet Inbakavi when in a state of inebriety, in the house of Mr. \textbf{Philip Rodrigo Muttukishna}, the Kutcherry Mudeliar, tried to ascertain whether the powers of mind, or the strength of alcohol predominated in him. He asked Imbakavi to improvise a stanza, which should begin with the word \emph{mai} (black paint), and end with the word \emph{maiyir} (hair). Inbakavi hesitated not a moment, but surprised the rival poet with the following beautiful lines:---
\begin{quotation}\small
மையினமோகார்மணலோவண்டினமோ கொண்டலதோ

சைவலமோயாதெனநான்‌ சாற்றுவேன்‌ --- மையலுறு

கங்கைகுலமாலாபிறகாமுகில்‌ மேலாசைகொண்ட

மங்கையிவள்‌கொண்டைமயிர்.
\end{quotation}
\begin{quotation}\small
\emph{To what shall I liken the colour of the hair of the condé of the damsel, who is enamoured of the noble Abraham\footnote{The eldest son of the Kutcherry Mudeliar.} of the race of Ganga? Shall I liken it to the black paint, or to the black sand, or to the black beetle, or to the black cloud, or to the black moss?}
\end{quotation}
\paragraph{}
Before he quitted Jaffna he wrote a drama, called \emph{Kuravanji} (குறவஞ்சி), in honour of the Kutcherry Mudeliar, which elicited general admiration for its musical cadence, and its vivid colouring and imagery. This was his only composition of any bulk; but he has left several very singularly beautiful stray pieces, which he addressed to his different benefactors. One day when he was drunk, and attending the feast in the Hindu temple at Nellore as a spectator, he hymned the praises of \emph{Skanda} in an ode with the chorus (மயில்வாகனங்கொண்டாண்டிநல்லூர்க்கந்தன் மயில்வாகனங்கொண்டாண்டி), ``\emph{Skanda of Nellore has mounted his vehicle, the peacock, yea he has mounted his vehicle, the peacock},'' but in his sober moments feeling compunction for it, he instantly composed another ode commemorating the crucifixion of our Lord, the chorus of which is as follows:--- (பத்திராசனங்கொண்டாரேயேசுநாயகர் சிலுவைப்பத்திராசனங்கொண்டாரே) ``\emph{Lord Jesus has enthroned himself, yea he has enthroned himself on the Cross.}''

A supplicatory stanza of Inbakavi, addressed to our Lord in imitation of the one addressed to \emph{Siva} by \textbf{T\={a}yum\={a}naver},\footnote{Vide தாயுமான சுவாமி திருப்பாடற்றிரட்டு, p. 7. Madras Edition.} is too good to be omitted, and we therefore give it below with an English paraphrase.
\begin{quotation}\small
நாலாயிரஞ்செனர்களேழப்பமுண்டுபசிநலிவுதீர்ந்திடவுமிலையோ --- நாரமதுரசபானமாகவிலையோதுஞ்சி நாலுநாளானமனுவுந் --- தாலமிசைமீண்டுமுயிராவிபெறவில்லையோ சட்டையின்விளிம்புதொட்ட --- தையல்பிணிநோய்வாதை தீரவிலையோபிரசண்டமுமடங்கவிலையோ --- வேலைகயல்வாயினிதிதோணவிலையோ மூகர்விண்டுமுரைபேசவிலையோ --- விழியூனர்நேத்திரந்தெளியவிலையோ பங்கர்விரைவாய்‌நடக்கவிலையோ --- சோலியாகியகவலையாலுழலுமேழைகை தூக்குவதுனக்கருமையோ --- துங்கர்பன்னிருவர்குல வம்சமாதீன்றருள்சுபுத்‌திரனாகியயேசுவே.
\end{quotation}
\begin{quotation}\small
\emph{Have not four thousand persons satisfied their hunger with seven loaves? Was not water turned into wine, and the man who had been dead four days restored to life? Did not the woman, who touched the hem of thy garment, become cured of her disease, and the storm abate at thy word? Was not silver found in the mouth of the fish of the sea, and did not the dumb speak, the blind see, and the lame walk swiftly? Is it then hard for thee, Oh Jesus, the holy son whom the Virgin of Israel has brought forth! to help thy servant who is embroiled in trouble, and beset with many difficulties.}
\end{quotation}
\paragraph{}
Inbakavi died at Manapar in the year 1835.

\section{Ireiyen\={a}r --- இறையனார்.}
\paragraph{}
This poet was one of the forty-nine professors of the Madura College. He appears to have been a prodigy of learning, and the Hindus therefore absurdly believe that he was \emph{Siva} himself, who in disguise mingling with the professors of the College, contributed to the improvement of the Tamil literature. A treatise on the composition of erotic poems with examples, under the title of \emph{Ireiyen\={a}r Agapporul} (இறையனார் அகப்பொருள்), is ascribed to his pen; and tradition has preserved the following impromptu which he uttered on the merits of the \emph{Kural}:---
\begin{quotation}\small
என்றும்புலராதியாணர்நாட்‌செல்லுகினு

நின்றலர்ந்துதேன்‌பிலிற்று நீர்மையதாய்க் --- குன்றாத

செந்தளிர்க்கற்பகத்தின்‌ றெய்வத்திருமலர்போன்‌

மன்புலவன்வள்ளுவன்வாய்ச்சொல்‌.
\end{quotation}
\begin{quotation}\small
\emph{The Kural which has proceeded from the mouth of Valluver, the king of poets, will never lose its beauty by the lapse of time; it will be always in its bloom, shedding honey (of instruction) like the flower of the tree in Indras paradise.}
\end{quotation}

\section{Iretteiyer. --- இரட்டையர்.}
\paragraph{}
Two poets who were twin brothers, and who acquired great renown by their facility in extempore versification, during the reign of \textbf{Pr\={a}krama P\={a}ndiyen} in the fourteenth century of the Christian era. From their birth they suffered under great physical disadvantage, the elder being a cripple, and the younger blind. They depended for their subsistence solely on their powers of improvising verses, and when they went out for this purpose, the younger used to carry on his shoulders the elder, who led him in his way. Linked together by birth and misfortune, there was one tie which even bound them still closer. In whatever they improvised, the younger uttered the first two lines of a stanza and the elder the last two. Many of their extemporaneous stanzas are still extant, of which we will select but two as specimens. Going once to the house of a wealthy man, named \textbf{Tiruvenkatesen}, who was distinguished for his generous disposition, he was about to give them some presents, but his elder brother dissuaded him from doing it. Upon which the younger Iretteiyer observed.
\begin{quotation}\small
தேன்பொழியும்வாயான்திருவெங்கடேசனுட

னேன்பிறந்தானிந்தவினியான்காண்‌
\end{quotation}
\begin{quotation}\small
\emph{Why was this unpleasant man born of the same parents with the honey-mouthed Tiruvenkatesen?}
\end{quotation}
\paragraph{}
To which the elder Iretteiyer immediately replied thus:
\begin{quotation}\small
\hspace*{15mm}யான்சொலக்கேள்‌ ---

சீதேவியார்பிறந்த செய்யதிருப்பாற்கடலில்‌

மூதேவியேன்‌பிறந்தாள்‌முன்‌.
\end{quotation}
\begin{quotation}\small
\emph{Hear me!}

\emph{Was not M\={u}devi (the Goddess of misfortune) born in the same milky sea before Sidevi (the Goddess of fortune).}
\end{quotation}
\paragraph{}
On another occasion having lodged for a night in a \emph{Pilleiyar's} temple, they had taken the precaution to deposit under the altar the money, which they had earned during the day, wrapt in a cloth, and missing it the next morning, they railed against the god in the following stanza:
\begin{quotation}\small
தம்பியோபெண்திருடி தாயாருடன்பிறந்த

வம்பனோநெய்திருடுமாமாயன்‌ --- அம்புவியில்‌

மூத்தபிள்ளையாரே முடிச்சவிழ்த்தீர்போமோ

கோத்திரத்திலுள்ளகுணம்‌.
\end{quotation}
\begin{quotation}\small
\emph{Thy younger brother (Skanda) is a stealer of maidens, and that vile M\={a}yen (Vishnu), who was born with thy mother, is a stealer of butter, and O Pilleiy\={a}r (Ganesa)! thou hast now stolen the money which was wrapt in the cloth. Will the natural disposition of the family be ever changed.}
\end{quotation}
\paragraph{}
Neither the place, nor the date of the Iretteiyer's death is known, but a tradition would have it that they both died simultaneously. The Rev. Mr. \textbf{Taylor} in the 2d volume of his Oriental Historical Manuscripts, Appendix, D., quotes the following stanza in praise of \textbf{Pr\={a}krama Pandiyen}, as the production of the two brothers, and the reader will find that we have adopted his translation, with a few alterations in order to make it approximate to the original:---
\begin{quotation}\small
குறுமுனிவன்கொல்லியிலோ தமிழ்ப்பிறந்ததம்மானே

ஒப்பரியதிருவிளையாட்டுறையூரோவம்மானே

திருநெடுமாலவதாரஞ்‌ சிறுபுலியோவம்மானே

சிவன்முடிமேலணிவதுவுஞ்‌ செங்கதிரோவம்மானே

கரையழிந்தகாவிரியோ வுரைதெரிந்ததம்மானே

கடிப்பகைக்குவேம்போ கடுக்கையோவம்மானே

பரவைபடிந்ததுஞ்சோழன்‌ பதந்தனிலோவம்மானே

பராக்கிரமபாண்டியன்புகழைப்‌ பகரெளிதோவம்மானே
\end{quotation}
\begin{quotation}\small
\emph{O Lady! Was it from the mount Kolli the dwarfish sage (Agastiyer) produced the Tamil language?} (\emph{No, from the Pothiya in the P\={a}ndiya country.})

\emph{O Lady! Did the incomparable sports (of Siva) take place at Urei\={u}r?} (\emph{No, at Madura, the metropolis of the P\={a}ndiya country.}) 

\emph{O Lady! Was Vishnu ever metamorphosed into a tiger?} (\emph{No, into a fish, the device on the banner of the P\={a}ndiya kings.})

\emph{O Lady! Is it the sun which Siva wears over the tuft of hair upon the crown of his head?} (\emph{No, it is the moon the progenitor of the P\={a}ndiya race}).

\emph{O Lady! Was it the bank destroying K\={a}veri which saw the writing (of Sampanter) ascend up against its stream?} (\emph{No, it was Vaigai, the river of Madura.})

\emph{O Lady! What is more heroic, the \={A}tti or Veppa garland?} (\emph{the Veppa, worn by the P\={a}ndiya kings, vowing thereby to conquer or die.})

\emph{O Lady! Was it at the feet of the Chola the sea submitted?} (\emph{No, it was at the feet of Ugra P\={a}ndiyen.})

\emph{O Lady! Is it then a light thing to declare the praise of Pr\={a}krama P\={a}ndiyen?}
\end{quotation}
\paragraph{}
N. B. The answers within brackets are our own, placed for the purpose of rendering the sense of the questions intelligible. It would appear that the poets praise \emph{P\={a}ndiya}, at the expense of \emph{Chola}, by reminding that the latter cannot glory in such things as the former can.

\section{Kachiyapper. --- கச்சியப்பர்.}
\paragraph{}
Kachiyapper was a \emph{Br\={a}hman} of the \emph{Adi Saiva} sect, who officiated in the temple of \emph{Ek\={a}mbeswara} at Kanjipuram, in the Carnatic. He was no less conspicuous as a poet than as a \emph{Ved\={a}ntist}, as is evident from the poetical talents which he has displayed in his imitation of the \emph{Sankara Sanhita} of the \emph{Skanda Pur\={a}na} in Tamil, under the title of \emph{Kandap\={u}ranam} (கந்தப்புராணம்), comprising 13,305 stanzas in six books. He undertook this work upon the request of certain illustrious personages at K\={a}njipuram, and finished it, according to the date specified in the following stanza prefixed to \textbf{Parasir\={a}ma Mudeliar's Kandapur\={a}na V\={a}sakam}, in the \emph{Saka} year 700, corresponding with \textsc{ce} 778.
\begin{quotation}\small
வேதவேதாந்தமாகமத்தருமநூலதாம்‌ விரிந்திடுபுராணம்‌

போதமார்முனிவரிசைத்திடுபுராணமிருதிகள்சங்கிதைப்பொருந்து

நாதமார்சடங்கஞ்‌ சிப்பிகள்பிறவு நன்குணர்கச்சிமாமுனிவ

னேதமில்சகாத்த மேழெனுநூறிற்‌ காந்தத்தைத் தமிழினாலிசைத்தான்‌.
\end{quotation}
\begin{quotation}\small
The \emph{Saivas}, especially in Jaffna, consider it as one of their most sacred works next to the \emph{Vedas}, and regularly recite it in their temples once a year with great solemnity.
\end{quotation}

\section{Kadigaimuttu Pulaver. --- கடிகைமுத்துப்புலவர்.}
\paragraph{}
Kadigaimuttu Pulaver was a celebrated poet retained at the court of \textbf{Jagav\={i}rar\={a}ma Venkateswara Ettappa N\={a}yaker}, R\={a}ja of Etti\={a}puram, of which place he was a native. Of his compositions the \emph{S\={a}mudra Vil\={a}sam} (சமுத்திரவிலாசம்), a panegyric in 100 stanzas, is deservedly esteemed for its elegant style. It is inscribed to the R\={a}ja himself, between whom and the ocean many curious similes are drawn with an ingenuity scarcely surpassed. The \emph{Dikvijayam} (திக்குவிசயம்), another panegyric, consisting of 325 stanzas, has also much merit, and redounds to his credit. The hero of this poem is \textbf{Varagunar\={a}ma Vanniyanar}, R\={a}ja of Siv\={a}giri, the pomp and splendour of whose court, and his bravery and successful forays the author describes in the most enthusiastic strains.

\section{Kall\={a}der. --- கல்லாடர்.}
\paragraph{}
This poet was one of the forty-nine professors of the Madura College, and like many of the learned men of his time, he lived and died an ascetic. He appears to have devoted his attention chiefly to Grammar and Rhetoric, of the principles and rules of which two sciences he wrote a treatise, calling it after his own name \emph{Kall\={a}dam} (கல்லாடம்). This work, however, is not now in general use, but it nevertheless deserves the encomium bestowed on it in the following stanza, which one of its admirers has prefixed to it:
\begin{quotation}\small
கல்லாடர்செய்பனுவற்‌ கல்லாட நூறுநூல்‌

வல்லார்சங்கத்தில்‌ வதிந்தருளி --- சொல்லாயிம்‌

மாமதுரையீசர்‌ மனமுவந்துகேட்டுமுடி

தானசைத்தார்நூறுதரம்.
\end{quotation}
\begin{quotation}\small
\emph{When Kallader's work the Kall\={a}dam, consisting of a hundred stanzas, was recited before the College of the learned, the Lord of Madura having heard it, was so much pleased that he nodded approbation a hundred times.}
\end{quotation}
\paragraph{}
Kallader was present when \textbf{Tiruvalluver} submitted the \emph{Kural} to the College bench, and we subjoin here the impromptu, which he is said to have uttered on the incomparable merits of that production.
\begin{quotation}\small
ஒன்றேபொருளென்னின்‌ வேறென்பவேறென்னி

னன்றென்பவாறுசமயத்தார்‌ --- நன்றென்ன

வெப்பாலவருமியைபவேவள்ளுவனார்‌

முப்பான்மொழிந்தமொழி.
\end{quotation}
\begin{quotation}\small
\emph{Of the six sects, one will condemn the system of the other, but none of them will condemn the system propounded by Valluver in his Kural; as it has the merit of harmonizing the suffrages of them all, so that each sect would admit it to be its own.}
\end{quotation}

\section{Kalatt\={u}r Kizh\={a}r. --- களத்தூர்க்கிழார்.}
\paragraph{}
This poet was a native of \emph{Kalatt\={u}r}, a town near Madura, and he is numbered among the forty-nine professors of the Madura College. The only relic of his compositions is the following stanza on the \emph{Kural}.
\begin{quotation}\small
ஒருவரிருகுறளே முப்பாலினோதுந்

தருமமுதநான்குஞ்சாலு --- மருமறைக

ளைந்துஞ்சமயநூலாறுநம்வள்‌ளுவனார்‌

புந்திமொழிந்தபொருள்‌.
\end{quotation}
\begin{quotation}\small
\emph{He who studies the two-lined verses in the three divisions of Valluver's Kural will obtain the four things (virtue, wealth, pleasure and eternal happiness); for they contain the substance of the five Vedas (including the Mah\={a}bh\={a}rat), and the six systems of the six sects.}
\end{quotation}

\section{Kamber. --- கம்பர்.}
\paragraph{}
Kamber was the most eminent poet of his age. He was the son of a king of \emph{Kamban\={a}du}, a district in the Carnatic. His mother when pregnant of him was forced to flee in disguise in consequence of a rebellion, in which her husband was assassinated, and to take refuge in the house of an \emph{Ochchen}\footnote{One of the \emph{Kudi-makkels} or village servants whose particular duty it is to offer sacrifice to the goddess \emph{Pid\={a}ri} for the villagers.} at \emph{Tiruvalund\={u}r}, where she gave him birth. Under the hospitable roof of the \emph{Ochchen} he remained until his seventh year, when a wealthy farmer of the name of \textbf{Sadeiyen}, who lived at \emph{Vennei Nell\={u}r}, having been informed of the circumstances of his case, invited him, and his mother over to his village, and placing them in a house, which he had prepared for them, supplied them with all necessaries. He also carefully watched over the studies of his protégé, whose superior genius developed itself even in his childhood. The young poet it is said would not compose a stanza for less than the reward of a hundred \emph{pons}.\footnote{A coin, valued at 3 s. 6 d. sterling.} His fame reaching the ears of \textbf{R\={a}jendra Chola}, that king invited him to \emph{Urei\={u}r}, and appointed him as one of the poets of his court; conferring on him at the same time the title of \emph{Kavi Chakravati}, or the king of poets. The \emph{Ram\={a}yana} being in Sanscrit, Kamber, \textbf{Ottak\={u}ter}, and several others of the court poets, separately undertook to imitate it in Tamil. When their imitations were finished, they were recited in the presence of \textbf{Kulotunga Chola}, who had succeeded \textbf{R\={a}jendra Chola} on the throne. This king preferred Kamber's imitation above the rest, being of opinion that it was not only sublime in style and sense, but that it approached nearest to the original. It was then by the king's order submitted to the judgment of the learned assembly of \emph{Br\={a}hmans} at \emph{Sirangam}, who, although they fully approved of the work, nevertheless strenuously objected to the praise bestowed by Kamber on his patron \textbf{Sadeiyen} in every hundredth stanza of the poem, and which he did upon the request of his mother.\footnote{See \emph{Ratnasab\={a}pati M\={a}lei}, verse 26.} They said that human praise should not be mixed up with the divine; but on his representing to them the necessity of his paying at least some tribute to his patron, they allowed him to confine his praises to every thousandth stanza. Overjoyed at the concession, he rose up in the midst of the assembly and said, ``I considered my patron as one in a hundred, but this learned assembly has considered him as one in a thousand.'' The \emph{Br\={a}hmans} admiring his wit bestowed on him the title of \emph{Mah\={a}-Kavi}, or the great poet, and he in return complimented them with a centum of verses, expressive of his veneration for \emph{Sadagopa Azhw\={a}r}, one of the \emph{Vaishnava} saints. \textbf{Varaguna P\={a}ndiyen}, the king of Madura, having heard of him about this time, invited him to his court. He there composed the \emph{Sarasvati Ant\={a}di} (சரஸ்வதியந்தாதி), a poem in thirty stanzas in praise of \emph{Sarasvati}, the goddess of poetry, which obtained for him great honours, and munificent presents from that king. Kamber's \emph{R\={a}m\={a}y\={a}nam} consists of 12,016 stanzas, arranged in five books; and a writer in the Calcutta Review says that it ``may be compared to \textbf{Pope's} Iliad.''\footnote{Calcutta Review, Vol. 25. p. 108.} In one of the commendatory stanzas which is prefixed to the work the year of \emph{Saka} 808 (\textsc{ce} 886) is specified as the date of its publication by Kamber; but the Rev. Mr. \textbf{Caldwell}, the author of the Comparative Grammar of the Dravidian Languages, rejects this date as spurious from the evidence of certain inscriptions found at Cape Comorin, and in the old Chalukya country, according to which the Chola kings who patronized Kamber lived only in the eleventh century of the Christian era.

Other works ascribed to Kamber are the \emph{K\={a}nji Pur\={a}num} (காஞ்சிபுராணம்), \emph{K\={a}nji Pilleitamil} (காஞ்சிப்பிள்ப்பிள்ளைத்தமிழ்), \emph{Ch\={o}la Kuruvanji} (சோழக்குறவஞ்சி), \emph{Silei Yezhupadu} (சிலையெழுபது) and \emph{Er-Yezhupadu} (ஏரெழுபது). Of these however only the \emph{Silei-Yezhupadu}, containing seventy stanzas in praise of the bow, and \emph{Er-Yezhupadu}, containing seventy stanzas in praise of the plough, are within our reach, and as far as we can judge, their style and versification are worthy of the noble poet.

Less is known of Kamber's personal history than could be desired. All that we gather from traditions is, that the latter part of his life was a struggle between poverty, and ill-fortune. His son \textbf{Ambik\={a}pati}, who was likewise a poet at the court of \textbf{Kulotunga Chola}, having been carrying on an amorous intrigue with the king's daughter, and betraying it in one of his poetical effusions, was, by the king's command, impaled alive,\footnote{See \emph{J\={a}yankond\={a}n Sudagam}, verse 20.} and he died after having lingered on the stake for a couple of days, repeating incessantly the subjoined verse, in the midst of his sufferings. ---
\begin{quotation}\small
எரியுதெரியுதென்னுள்ளீரலைப்பற்றி

யெரியுதவியுதிலையென்செய்வேன்‌

நஞ்சிலேதோய்ந்த நளினவிழிப்பெண்ணணங்கு

நெஞ்சிலேயிட்டநெருப்பு.
\end{quotation}
\begin{quotation}\small
\emph{What can I do? The fire which the divine damsel with venomed eyes has kindled in my heart burns within me, it burns, yea it burns without extinguishing.}
\end{quotation}
\paragraph{}
Kamber was so overwhelmed with grief at this sad event, that he discontinued attending on the king, and confined himself entirely to his house. Whilst in this state of mind, the king's son, who was returning from school, having taken refuge in his house from a fierce elephant, which was at large in the street, he snatched the stylus which the prince held in his hand, and with it stabbed and killed him on the spot. The king immediately ordered Kamber to be brought up, and shot at him with an arrow; but the wound not proving mortal, and through some intervention which is not stated, he affected his escape into the \emph{P\={a}ndiyen} country. He has commemorated this event in the following impromptu:---
\begin{quotation}\small
வில்லம்பு சொல்லம்பு மேதினியிலிரண்டுண்டு

வில்லம்பு சொல்லம்பை வெல்லாது --- வில்லம்பு

பட்டுதெடாவென்மார்பிற்‌ பார்வேந்தேயுன்கிளையைச்‌

சுட்டுதெடாவென்வாயிற்சொல்‌.
\end{quotation}
\begin{quotation}\small
\emph{There are two kinds of arrows on earth, of which one is shot out from the bow, and the other from the mouth. O king! thine arrow from the bow has hit my breast, but mine, the word of my mouth, has burnt up thy race.}
\end{quotation}
\paragraph{}
After his flight from the Chola Court, he entirely lost his poetical inspiration, and being reduced to poverty was obliged to support himself by daily labour.\footnote{See \emph{Jayankond\={a}n Sadagam}, verse, 11.} He died in the sixtieth year of his age, at \emph{N\={a}ttarasenkottei}, in the Ramnad country, and was buried there in a sitting posture as is customary with the \emph{Vaishnava} sect, to which he belonged. It is said that his grave now forms an object of pilgrimage amongst the Hindus, and that the earth taken from it is given to children who are dull-headed, from a superstitious belief that it would sharpen their intellect.

When the intelligence of his death reached the ears of \textbf{Otteik\={u}tur}, who was always his rival when alive, he is said to have his sorrow and lamentation in the following impromptu:---
\begin{quotation}\small
இன்றைக்கோகம்பனிறந்தநாளிப்புவியி

லின்றைக்கோவென்கவிதையேற்குநாள்‌ --- இன்றைக்கோ

பூமடந்தைவாழப்‌ பொறைமடந்தைவீற்றிருக்க

நாமடந்தை நூலிழந்தநாள்‌.
\end{quotation}
\begin{quotation}\small
\emph{Is it to-day that Kamber has died?}

\emph{Is it to-day that he accepts my elegiac lay?}

\emph{Is it to-day that Sarasvati\footnote{\emph{Sarasvati}, the goddess of poetry.} has been widowed?}

\emph{Whilst Lakshmi\footnote{\emph{Lakshami}, the goddess fortune.} is blooming and P\={a}rvati\footnote{\emph{Parvati}, the consort of \emph{Siva}.} is sitting in majesty.}
\end{quotation}
\paragraph{}
There is a tradition that Kamber once visited a certain king of Ceylon, and that the king was so charmed with his poetry that he was treated with unusual affability, so much so, that the Court poets became jealous of him and formed a plot to keep him at a distance from the Palace, by bribing a barber to claim relationship with him; but he soon found means to expose their artifices and they had the mortification to see him afterwards even honoured with a seat on the throne by the side of the king. We cannot, however, vouch for the truth of this tradition; but the following stanza in praise of \emph{Sarasvati}, which is generally attributed to him is said to refer to it:---
\begin{quotation}\small
வெள்ளைக்கலையுடுத்து வெள்ளைப்பணிபூண்டு

வெள்ளைக்கமலத்தில்‌ வீற்றிருப்பாள்‌ --- வெள்ளை

யரியாசனத்திலரசனோடென்னைச்‌

சரியாயிருத்திவைத்ததாய்.
\end{quotation}
\begin{quotation}\small
\emph{It is my mother who sits in majesty on the white lotus, robed in white garments and adorned with white ornaments, that gave me an equal seat with the king on the white throne, borne by the lion.}
\end{quotation}

\section{Kapiler. --- கபிலர்.}
\paragraph{}
Kapiler was one of the forty-nine professors of the Madura College. He was the brother of \textbf{Tiruvalluver}, and being exposed by his mother soon after his birth at Trivalur, he was found and adopted by a \emph{Brahman} of that place. When he arrived at his seventh year, his foster-father invited the other \emph{Brahmans} at Trivalur in order to invest him with the sacred cord, but they all peremptorily refused to accept the invitation on the ground that he not being a \emph{Brahman} was not entitled to be initiated into the privileges of their order. Upon hearing this, Kapiler repaired to the place where the \emph{Brahmans} were assembled, and in an \emph{Agavel},\footnote{A short poem in the \emph{Asiriya} metre.} which he composed and recited before them, refuted their high pretensions by proving that caste was not innate but accidental, and they not being able to rebut his arguments, at once admitted him into their rank by investing him with the sacred cord. A metrical version of the \emph{Agavel} in English by the Rev. \textbf{C. J. Robinson} has been published in Rev. \textbf{P. Percival's} Land of the Veda.

Kapiler was present in the Madura College when his brother's \emph{Kural} was recited, and he conveyed his opinion of the merits of that work in the following impromptu:---
\begin{quotation}\small
தினையளவுபோதாச்சிறுபுன்னீர்நீண்ட

பனையளவுகாட்டும்படித்தான் --- மனையளகு

வள்ளைக்குறங்கும்‌ வளநாடவள்ளுவனார்

வெள்ளைக்குறட்பாவிரி.
\end{quotation}
\begin{quotation}\small
\emph{O king of the fertile country where the birds reared in houses are lulled by the women's chants! Valluver's Kural is short in words but extensive in sense, even as in a drop of water on the blade of the Tinei (millet) might be seen reflected the image of the tall Panei (Palmyra tree).}
\end{quotation}

\section{Karunaipprak\={a}sa D\={e}siker. --- கருணைப்பிறகாசதேசிகர்.}
\paragraph{}
Karunaipprak\={a}sa Desiker was the youngest son of \textbf{Kum\={a}rasw\={a}mi Pand\={a}ram} of Kanjipuram, and was like him a staunch adherent of the \emph{Vira Saiva} sect. Having studied under \textbf{Velliambala Tambir\={a}n}, the head ascetic of the \emph{Matam} at Sindupunturei, in Tinnevelly, and who was distinguished for his extensive erudition, he became a poet of the first order; but before he could fully display his poetical talents, he was hurried away to a premature grave in the eighteenth year of his age. He wrote a panegyric, under the title of \emph{Ishtalinga Agavel} (இஷடலிங்கவகவல்) on the \emph{Phallic} emblem, which the \emph{Vira Saivas} wear, and afterwards began to write an epic poem, called \emph{K\={a}latti Pur\={a}nam} (காளத்திப்புராணம்), treating of the legends of the \emph{Saiva} shrine on the mountain \emph{K\={a}lahasti}, in the Carnatic; but his death prevented him from finishing it.

\section{Kavir\={a}ja Pand\={i}ter. --- கவிராசபண்டிதர்.}
\paragraph{}
Kavir\={a}ja Pand\={i}ter was a native of \emph{V\={i}rei}, in the Carnatic, but nothing is known concerning his life, and the time in which he lived. He is chiefly remembered as the author of a poem, entitled \emph{Saundariyalaghiri} (செளந்தரியலகரி), consisting of 104 stanzas, and which is highly esteemed by the \emph{Saktyas} on account of its extolling the greatness of \emph{P\={a}rvati}, whom they worship as the mother of the worlds. So far as regards its versification this poem has considerable merits; but the licentious metaphors with which it abounds, do not at all suit its pretensions as a religious work.

\section{Kavuniyan\={a}r. --- கவுணியனார்.}
\paragraph{}
The name of this poet occurs among the forty-nine professors of the Madura College. Excepting the fact that he was the author of the following impromptu in commendation of the \emph{Kural}, we know nothing further about him:---
\begin{quotation}\small
சிந்தைக்கினிய செவிக்கினிய வாய்க்கினிய

வந்தவிருவினைக்கு மாமருந்து --- முந்திய

நன்னெறிநாமறிய நாம்புலமைவள்ளுவனார்

பன்னியவின்குறள்வெண்பா.
\end{quotation}
\begin{quotation}\small
\emph{The short Distichs which the learned poet Valluver has composed in order that we may know the ancient right way are sweet to the mind to meditate on, sweet to the ear to hear, and sweet to the mouth to repeat; and they moreover form a sovereign medicine to promote good and prevent evil actions.}
\end{quotation}

\section{K\={a}daver K\={o}n. --- காடவர்கோன்.}
\paragraph{}
This poet was one of the kings of the \emph{Pallava} dynasty, who reigned at Kanjipuram in former times, and he also figures among the sixty-three special votaries of \emph{Siva}, whose legends constitute the subject of the \emph{Tirutonder Pur\={a}nam}. It is said that after a short reign he abdicated the throne in favour of his son, and passed the remainder of his life in devotional exercises. He wrote a panegyrical stanza on each of the \emph{Saiva} shrines in Southern India, but they are rarely met with at present.

\section{K\={a}kaip\={a}diniyan\={a}r. --- காக்கைப்பாடினியனார்.}
\paragraph{}
This poet was one of the disciples of \textbf{Agastiyer}, and he appears to have been as learned as his fellow student \textbf{Tolk\={a}ppiyan\={a}r}. No other particulars concerning his life are known, and it is impossible to ascertain them at this distance of time. He wrote a treatise on grammar, called \emph{K\={a}keip\={a}diniyam} (காக்கைப்பாடினியம்), after his own name, and this has been cited by subsequent writers on the subject, as a standard authority.

\section{K\={a}lam\={e}gam. --- காளமேகம்.}
\paragraph{}
The time when this poet flourished is usually referred to the reign of \textbf{Tirumalar\={a}yen} in the fifteenth century of the Christian era. His real name is not known. He obtained the sobriquet of \emph{K\={a}lamegam}, or the hail-pouring cloud, for the sharpness of his invectives against those who failed to reward his poetry. There is no doubt that he was possessed of superior talents, but his character has been tainted by a deplorable egotism, in proof of which we subjoin one of his poetical epistles to \textbf{Tirumalar\={a}yen}, wherein he trumpets forth his own attainments, and defies all other poets.
\begin{quotation}\small
தூதைந்துநாழிகையிலாறுநாழிகைதனிற்‌ சொற்சந்தமாலைசொல்‌லத்‌ --- துகளிலாவந்தாதியேழுநாழிகைதனிற்‌ றோகைபடவிரித்துரைக்கப் --- பாதஞ்செய்மடல்கோவை பத்துநாழிகைதனிற்‌ பரணியொருநாண்முழுதுமே --- பாரகாவியமெலாமோரிருதினத்திலே பகரக்கொடிக்கட்டினேன்‌ --- சீதஞ்செயுந்திங்கண்மரபினா னீடுபுகழ்செய்யதிருமலைராயன்‌முன்‌ --- சீறுமாறென்று மிகுதாறுமாறுகள்‌ சொற்றிருட்டுக்கவிப்புலவரைக்‌ --- காதங்கறுத்துச்‌ சவுக்கிட்டடித்துக்கதுப்பிற்புடைத்துவெற்றிக் --- கல்லணையினொடுகொடிய கடிவாளமிட்டேறுகவிகாளமேகநானே.
\end{quotation}
\begin{quotation}\small
``\emph{It is I the poet K\={a}lamegam who have hoisted a flag, proclaiming thereby that I shall compose a D\={u}tu, in five Naligais\footnote{A \emph{N\={a}ligai} is equal to 24 minutes of an English hour.}; a Male\={i} in six Naligais; an Ant\={a}di in seven Naligais; a Madal, or Kovai, in ten Naligais; a Parani in a day; and all kinds of epics in a couple of days. I will cut off the ears of the poets, who purloining the compositions of others, parade them as their own, and bluster and puff in the presence of the ever renowned king Tirumalarayen, of the race of the placid moon; I will whip them on their backs, slap them on their cheeks, and saddle and ride on them, bridling them with a hard bridle.}''
\end{quotation}
\paragraph{}
K\={a}lamegam has however left no poems of any magnitude, and of his compositions the most admired is a series of stanzas, under the title of \emph{Yamakandam}, which being composed of words two-fold in meaning, each contains a description of two different things, as will be seen by the subjoined specimens.
\begin{quotation}\small
நஞ்சிருக்குந் தோலுரிக்கும்‌ நாதன்முடிமேலிருக்கும்‌

வெஞ்சினத்தாற்பற்பட்டால்‌ மீழாது --- விஞ்சுபுகழ்‌

தேன்பாயுஞ்சோலை திருமலைராயன்பதியிற்‌

பாம்பாகும்‌ வாழைப்பழம்‌.

வாரிக்களத்தடிக்கும்‌ வந்துபுகுங்கோட்டையிற்குட்‌

போருக்குமாங்கே புரியுமாம்‌ --- நேரே

செய்க்கோலஞ்செய்யுந்‌ திருமலைராயன்பதியில்‌

வைக்கோலுமால்யானையாம்‌.
\end{quotation}
\begin{quotation}\small
1.\emph{It has poison, it sheds its skin, it sits on the crown of Siva's head, its bite under irritation cannot be cured. O! Such is the case with the Cobra Capella in the city of the renowned Tirumalarayen, where the groves flow with honey.}

2.\emph{It is easily bruised, its skin peels off, it hangs from the top of the parent tree, and it will not recover its former state if it comes in contact with the teeth. O! Such is the case with the plantain fruit in the city of the renowned Tirumalarayen, where the groves flow with honey.}

1.\emph{It takes up with the hand (the branches of trees), and beats (them) against the neck, it enters forts, and delights in battle. O! Such is the case with the Elephant in the city of Tirumalarayen, adorned with fields.}

2.\emph{It is taken up with the hand and beaten out in the thrashing floor, it is then stacked and being twisted into cord is used for tying around the heap of corn. O! Such is the cast with the straw in the city of Tirumalarayen, adorned with fields.}
\end{quotation}
\paragraph{}
Of the anecdotes related of K\={a}lamegam, we think the following will interest the reader:---

One day as he was attending the temple at Tiruval\={u}r, his attention having been drawn to a diamond, which was sparkling on the breastplate of the image of \emph{Tiy\={a}gar} (\emph{Siva}), he cracked a joke at the expense of the god, in this impromptu:
\begin{quotation}\small
அன்னவயல்சூழுமாரூர்த்தியாகர்நெஞ்சி

லின்னம்‌வயிரமிருக்குதோ --- முன்னந்

தொண்டன்மகனைத்தின்றுஞ்‌ சோழன்மகனைக்கொன்றுஞ்

சண்டனெமனுதைத்‌ததாள்‌.
\end{quotation}
\begin{quotation}\small
\emph{Has Tiy\={a}gar of Ar\={u}r, which is surrounded by rice-fields, still Vairam (diamond, meaning also wrath) in his breast? Has he not been appeased, even after devouring the son of Sirutonden, killing the son of Chola, and striking Yama with his foot?}
\end{quotation}
\paragraph{}
Another day he went to a temple of \emph{Skanda} to witness the feast, but found it difficult to push his way through the mass of the people that thronged the temple. Vexed at this disappointment, he gave vent to his feelings in an impromptu as follows:---
\begin{quotation}\small
அப்பனிரந்துண்ணியாத்தாள்பெருநீலி

சொற்பெரியமாமன்‌ றொழிற்றிருடன்‌ --- இப்புவியில்‌

தமையன்பெருவயிறன்றாரணியெல்லாமறியக்‌

குமராவுனக்கேன்வீண்கூத்து.
\end{quotation}
\begin{quotation}\small
\emph{O Kumara! Why this vain parade for thee! Thy father (Siva) is a mendicant, thy mother (P\={a}rvati) is an ogress, thy uncle (Kishna) is a professed thief, and thy brother (Ganesa) is a glutton as all the world knows.}
\end{quotation}
\paragraph{}
When he once visited Kannapuram, the \emph{Br\={a}hmans}, who officiated in the \emph{Vaishnava} temple there, having importuned him to chant the praise of their god, he improvised the following stanza:
\begin{quotation}\small
கன்னபுரமாலே கடவுளிலுநீயதிக

முன்னிலுமோநானதிகமோதக்கேள்‌

உன்பிறப்போபத்து வோர்சிவனுக்கொன்றுமில்லை

யென்பிறப்போவெண்ணத்தொலையாது.
\end{quotation}
\begin{quotation}\small
\emph{O the Lord of Kannapuram! thou art greater than Siva, but I am greater than thee; for hear me, thou hadst ten births, and Siva had none, but my births are too many to be counted.}
\end{quotation}
\paragraph{}
An anonymous admirer of K\={a}lamegam, has commemorated him in the subjoined verse:
\begin{quotation}\small
கம்பனென்றுமம்பனென்றுங்‌ காளியொட்டக்கூத்தன்றுங்‌

கும்பனென்றும்பேர்கொடுப்பதுண்டோ --- அம்புவியில்‌

கன்னாவதாரன்‌ காளமேகப்புலவன்‌

அன்னாளிலேபிறந்தாக்கால்‌.
\end{quotation}
\begin{quotation}\small
\emph{Would Kamben, and Amben, K\={a}lidas and Ottak\={u}ten, nay even Kumben (Agastiyer), have acquired poetical fame on earth, if the poet K\={a}lamegam, who is the incarnation of Karna\footnote{An ancient king of Angadesa famous for his liberality.} himself, had been born in their days?}
\end{quotation}

\section{K\={a}ngayer. --- காங்கேயர்.}
\paragraph{}
K\={a}ngayer was a \emph{Siva} ascetic, and poet, who appears to have lived prior to \textbf{Mandalapuruder}; but the exact date of his existence cannot be ascertained. He composed a Tamil lexicon in verse, under the title of \emph{Urichol Nigandu} (உரிச்சொல் நிகண்டு), which unlike others of the kind, being written in the \emph{Venba} metre, and well adapted for beginners, is now generally used in schools. An edition of this work published at Pondicherry in 1840, comprises only 220 stanzas: but the one which issued from the Press at Jaffna in 1858, has 330 stanzas, and whether the addition is genuine, or not, it appears to us to be a great improvement.

\section{K\={a}reikal Ammeiy\={a}r. --- காரைக்காலம்மையார்.}
\paragraph{}
The real name of this celebrated poetess was \emph{Punithavati}. She was the only daughter of \textbf{Tanatattan}, a \emph{Chetty} who resided at Karikal. Her father married her early in life to the son of another \emph{Chetty} at Negapatam, but her husband sometime afterwards deserted her, and going over to the \emph{P\={a}ndi} country, took another wife, and established himself there. After having failed to effect a reconciliation with him she assumed the habit of an ascetic, and repaired to the \emph{Imaus}, where she died. Of her several poems, the most admired are the \emph{Atputhatiruvant\={a}di} (அற்புதத்திருவந்தாதி) \emph{Tiruviretteimanim\={a}lei} (திருவிரட்டைமணிமாலை) and \emph{M\={u}ttatiruppathikam} (மூத்ததிருப்பதிகம்), which are in high esteem amongst the \emph{Saivas}.

\section{K\={a}rikannan\={a}r. --- காரிக்கண்ணனார்.}
\paragraph{}
This poet was a native of \emph{K\={a}ver\={i}patnam}, in the Carnatic. He ranks as one among the forty-nine professors of the Madura College, but we are ignorant of the particulars of his life. His compositions have all perished, save the following impromptu on the \emph{Kural}:
\begin{quotation}\small
ஐயாறுநூறு மதிகாரமூன்றுமா

மெய்யாயவேதப்‌ பொருள்விளங்கப்‌ --- பொய்யாது

தந்தானுலகிற்குத்‌ தான்வள்ளுவராகி

யந்தாமரைமேலயன்‌.
\end{quotation}
\begin{quotation}\small
\emph{It is no other than Ayen (Br\={a}hma) himself, who is seated on the beautiful lotus-flower, assuming the form of Valluver, has given to the world the truths of the Vedas in three hundred and thirty chapters, that they may shine without being mixed up with falsehood.}
\end{quotation}

\section{K\={a}riy\={a}r. --- காரியார்.}
\paragraph{}
Nothing is known of this poet but what is contained in the preface to his \emph{Kanakkadik\={a}ram} (கணக்கதிகாரம்), according to which he was a native of \emph{Kurkai}, the ancient residence of the \emph{P\={a}ndiyans}, situated on the northern bank of the river \emph{Tambraparni}, and that he claimed descent from one of those kings. The \emph{Kanakkadik\={a}ram} is a treatise in verse on Arithmetic, compiled from materials derived from the \emph{Ganitas\={a}stra}, \emph{Bhuvanadipa}, and other Sanskrit works on the same subject. After the usual invocation to \emph{Ganesa} it first describes the system of weights and measures, and then the different operations of Arithmetic, and concludes with a series of arithmetical problems with their solution. It is on the whole a valuable work.

\section{K\={a}riy\={a}r. --- காரியார்.}
\paragraph{}
Little is known of the personal history of K\={a}riy\={a}r beyond the fact that he was a \emph{Br\={a}hman} of \emph{Tirukkadav\={u}r}, in the Carnatic, who distinguished himself by composing amatory poems called \emph{K\={o}vei} (கோவை), which being recited before the \textbf{Chera}, \textbf{Chola}, and \textbf{Pandiya} kings of his time, obtained for him their patronage. The date of his existence is uncertain; but it was no doubt anterior to the reign of \textbf{Anabh\={a}ya Chola}, as he figures as one of the sixty-three special votaries of \emph{Siva} in the \emph{Tirutonder Pur\={a}nam}, which was written by \textbf{Sekilar} under the auspices of that king. He is said to have lavished all his fortune in building, and endowing temples for the worship of \emph{Siva}, and in maintaining the \emph{Brahmans}.

\section{K\={i}randeiy\={a}r. --- கீரந்தையார்.}
\paragraph{}
K\={i}randeiy\={a}r was one of the forty-nine professors of the Madura College, but there is nothing extant of his compositions besides the subjoined impromptu, which he uttered in admiration of the \emph{Kural} 
\begin{quotation}\small
தப்பாமுதற்பாவாற்றாமாண்டபாடலினான்‌

முப்பாலினாற்பான்மொழிந்தவரெப்பாலும்‌

வைவைத்தகூர்வேல்‌ வழுதிமனமகிழத்‌

தெய்வத்திருவள்ளுவர்‌.
\end{quotation}
\begin{quotation}\small
\emph{It is the divine Valluver, who by his great skill in poetry, has imparted in a series of verses of the first class, arranged in three parts, the means of obtaining the four things (meaning virtue, wealth, pleasure, and eternal happiness), that the heart of P\={a}ndiyen, who wields the edged javelin may be delighted therewith.}
\end{quotation}

\section{Kodign\={a}zhanm\={a}ni P\={u}dan\={a}r. --- கொடிஞாழன் மானிபூதனார்.}
\paragraph{}
This poet was one of the forty-nine professors of the Madura College; but neither his birthplace, nor his personal history can be ascertained at present. His name, however, has survived the lapse of time by the following impromptu, which he uttered in approbation of the \emph{Kural}:
\begin{quotation}\small
அறனறிந்தேமான்ற பொருளறிந்தேமின்‌பின்‌

றிறனறிந்தேம்வீடு தெளிந்தேமறனெறிந்த

வாளார்நெடுமாற வள்ளுவனார்தம்வாயாற்‌

கேளாதனவெல்லாங்கேட்டு.
\end{quotation}
\begin{quotation}\small
\emph{O P\={a}ndiya, the mighty king, who wields the sword brandished among the enemy! Having heard to-day from the mouth of Valluver what we never heard before, and comprehended their meaning, we have been enlightened with the knowledge of eternal happiness.}
\end{quotation}

\section{Kongkaner. --- கொங்கணர்.}
\paragraph{}
Kongkaner was a celebrated philosopher, and physician, probably a native of (as his name would seem to imply) \emph{Kongan\={a}du}, the modern Coimbatore country. It is said that he became an ascetic in the prime of his youth, and having associated himself with the sages who then abode on the \emph{Pothiya} mountain, was instructed by them in all branches of knowledge. Some make him contemporary with \textbf{Agastiyer}, and some with \textbf{Tiruvalluver}, but nothing certain is known of his epoch. He wrote a treatise against polytheism, under the title of \emph{Kongkanar Gn\={a}nam} (கொங்கணர்ஞானம்), and two treatises on medicine, one of which is entitled \emph{Kadeikandam} (கடைக்காண்டம்), and the other \emph{Gunavagadam} (குணவாகடம்).

The following is a stanza taken from his \emph{Gn\={a}nam}.
\begin{quotation}\small
கடவுளோனொருவனுண்டே வேதமொன்றே காரணசற்குருதீட்‌சை தானுமொன்றே --- யடைவுடனே யவனருளும்பதவியொன்றேயம்‌புவியின்‌ மனுப்பிறவியானதொன்றே --- நடைவழியும்பலமனுவோர்க்கொன்றேயல்லால்‌ நால்வேதமறுசமய நடக்கைவேறாய்த் --- திடமுடையதேவர்பலருண்டென்போர்கடீய்நரகுக்குள்ளாவர்‌திண்ணந்தானே.
\end{quotation}
\begin{quotation}\small
\emph{There is but one God, there is but one Veda, there is but one way of initiation by the good spiritual guide, there is but one kind of bliss which he grants, and there is but one genus amongst mankind upon the earth. They who hold the contrary by asserting that there are four Vedas, and six sects, and many gods, will surely enter the fiery hell.}
\end{quotation}

\section{K\={o}raker. --- கோரக்கர்.}
\paragraph{}
K\={o}raker is reckoned among the nine \emph{Siddhas}, who had acquired supernatural qualifications and powers by intense austerities. Be this as it may, there is no doubt that he was well versed in poetry as well as in medicine, and chemistry. He has left a work in 100 verses, under the title of \emph{K\={o}raker Vyppu} (கோரக்கர்வைப்பு), treating of the composition and preparation of medicines, and it is said that it was he who first introduced the use of \emph{Ganja} (\emph{Cannabis indica}) in the Tamil practice; hence that weed is now generally called \emph{K\={o}raker-m\={u}li} (கோரக்கர்மூலி), or K\={o}raker's drug.

\section{K\={o}thamanar. --- கோதமனார்.}
\paragraph{}
K\={o}thamanar was one of the forty-nine professors of the Madura College. His personal history has been lost, and of his compositions there is nothing remaining, except the following impromptu on the \emph{Kural}:
\begin{quotation}\small
ஆற்றலழியுமென்றந்தணர்கணான்மறையைப்‌

போற்றியுரைத்தேட்டின்‌ புறத்தெழுதா --- ரேட்டெழுதி

வல்லுநரும்‌ வல்லாரும்‌ வள்ளுவனார்முப்பாலைச்‌

சொல்லிடினுமாற்றல்சோர்வின்று.
\end{quotation}
\begin{quotation}\small
\emph{The Br\={a}hmans preserve the four Vedas orally, and never commit them to writing for reason that if they were read by all, they would be less valued; but the Kural of Valluver though committed to writing, and read by all, would nevertheless not lose its value.}
\end{quotation}

\section{K\={o}v\={u}rkizhlar. --- கோவூர்க்கிழார்.}
\paragraph{}
Little more is known of K\={o}v\={u}rkizhlar than that he was one of the forty-nine professors of the Madura College. The subjoined impromptu on the \emph{Kural}, is the only relic of his compositions which has been preserved:
\begin{quotation}\small
அறமுதனான்குமகலிடத்தோரெல்லாந்

திறமுறத்தேர்ந்து தெளியக் --- குறள்வெண்பாப்‌

பன்னியவள்ளுவனார்‌ பான்முறைநேரோவ்வாதே

முன்னைமுதுவோர்மொழி.
\end{quotation}
\begin{quotation}\small
\emph{None of the works even of the sages of antiquity will bear comparison with the Kural, whereby Valluver has enlightened the world with the knowledge of virtue, wealth, pleasure, and eternal happiness.}
\end{quotation}
\section{Kuhei Namasiv\={a}yer. --- குகை நமச்சிவாயர்.}
\paragraph{}
This \emph{Saiva} ascetic was alike distinguished as a \emph{Vedantist} and a poet, and his memory is still held in high estimation by his co-religionists. He first established himself at Tirunamali, but afterwards went over to Chillambaram, and spent the remainder of his life there. His principal poem, the \emph{Sh\={o}naghirim\={a}lei} (சோணகிரிமாலை), is a soliloquy addressed to his own soul. In the following stanza, which we have selected from it as a specimen, the reader will find that the author in exhorting his soul to remember \emph{Siva}, before the calamities of old age and death shall have overtaken him, expresses himself in the same strain as \textbf{Solomon} in the Ecclesiastes, Ch. 12. 1. et seq.
\begin{quotation}\small
காயநெகிழாமுன்‌ கண்களிருளாமுன்‌

வாயிலுள்ளபல்லு வழுவாமுன்‌ --- தாயம்பார்த்‌

தோடிநமன்வாராமுனுள்ளமேசோணகிரி

நாடிநமவென்றேநட.
\end{quotation}
\begin{quotation}\small
\emph{Before the body decays, before the eyes grow dim, before the teeth drop, and before Yama (the god of death), watching the opportunity, swiftly approaches; O my soul! turn thy thoughts to the mountain Shona (Tirunamali) and walk on towards it, saying I adore thee.}
\end{quotation}
\section{Kulapatiy\={a}r. --- குலபதியார்.}
\paragraph{}
Kulapatiy\={a}r, called also \emph{Kulachireiy\={a}r}, was minister of \textbf{K\={u}na P\={a}ndiyen}, as well as one of the forty-nine professors of the Madura College. When the king embraced the \emph{Jaina} religion he did not follow his example, but continued a \emph{Saiva}, and carried his zeal for his own religion to such extremes that no sooner had he, in concert with \textbf{Sampanter}, succeeded in persuading the king to recant, than he raised a cruel persecution against the \emph{Jainas}, and caused eight thousand of them to be impaled alive, and the rest driven from the country. Of his poetical compositions there is nothing now forthcoming except the subjoined impromptu in praise of the \emph{Kural}.
\begin{quotation}\small
உள்ளக்கமலமலர்த்தியுளத்துள்ள

தள்ளற்கரியவிருடள்ளுதலால்‌ --- வள்ளுவனார்

வெள்ளைக்குறட்பாவும்‌ வெண்கதிருமொக்குமெனக்

கொள்ளத்தகுங்குணத்தைக்கொண்டு.
\end{quotation}
\begin{quotation}\small
\emph{As the Kural of Valluver causes the lotus-flower of the heart to expand, and dispels from it the darkness which cannot otherwise be dispelled, it may well be compared to the hot-rayed sun, which causes the lotus flower of the tank to expand, and dispels the darkness from the face of the earth.}
\end{quotation}
\section{Kulasekhara Perumal. --- குலசேகரப்பெருமாள்.}
\paragraph{}
Kulasekhara Perumal being considered as one of the twelve \emph{\={A}zhv\={a}rs}, or special votaries of \emph{Vishnu}, who have been deified by the Hindus, his history is consequently very much intermixed with fable; but we may nevertheless glean from it some facts which are worthy of credence. He was the son of the king \textbf{Dridharrata}, who ruled over Malay\={a}lam, holding his court at \emph{Tiruvanji}. On succeeding his father he reigned but a few years, when abdicating the throne, he became a devotee, and set out on a course of pilgrimages to the different \emph{Vaishnava} fanes in the Carnatic. Whilst staying at Mannan\={a}rkovil, after a visit to Sirangam, K\={a}njipuram and Tirupati, he was suddenly overtaken by death. He appears to have been well skilled both in Sanskrit and Tamil poetry as he wrote the \emph{Mukunda Malei} (முகுந்தமாலை) in the former, and one hundred and five stanzas in the latter, which form a part of the \emph{N\={a}l\={a}yira prabhandam} (நாலாயிரபிரபந்தம்) now extant amongst the \emph{Vaishnavas}. The period of his existence cannot now be determined, but the fable would have that he was born in the twenty-eighth year of the \emph{Kali} age.

\section{Kumaragurupara D\={e}siker. --- குமரகுருபரதேசிகர்.}
\paragraph{}
Kumaragurupara D\={e}siker was a \emph{Saiva} ascetic and poet of considerable eminence, born of \emph{Vel\={a}la} parents at \emph{Strivyguntam}, in the Tinnevelly district, about two hundred years ago. He established his fame as a poet while yet a child; for it is reported of him that when only five years old, he composed a poem, called \emph{Kalivenb\={a}} (கலிவெண்பா), in praise of the shrine of \emph{Skanda} at Tiruchendur, and a few years later another poem, called \emph{Pilleitamil} (பிள்ளைத்தமிழ்), in praise of the shrine of \emph{Men\={a}kshi} at Madura, which he recited before the king \textbf{Tirumalan\={a}yaker}. He soon afterwards entered the \emph{Matam} at Dharmapuram, and after studying there the \emph{Saiva} system of philosophy, and improving his knowledge of the Tamil and Sanskrit languages and poetry, assumed the habit of an ascetic; and the heads of the \emph{Matam} being convinced of his deep learning, invested him with the title of \emph{Tambiran}. He then repaired to Benares, and presided over a \emph{Matam} in that city until his death, which occurred in the prime of his manhood. Besides his two juvenile productions, he was also the author of twelve others, of which the last one entitled \emph{N\={i}dineri-vilakkam} (நீதிநெறிவிளக்கம்), consisting of 102 stanzas on moral subjects, is considered the best, and has been translated into English in prose by Mr. \textbf{Stokes} of the Madras Civil service. The following selections from the \emph{N\={i}dineri-vilakkam}, with Mr. \textbf{Stokes'} translation, may not be found uninteresting:---
\begin{quotation}\small
தொடங்குங்காற்றுன்பமாயின்பம்பயக்கு

மடங்கொன்றறிவகற்றுங்கல்‌வி --- நெடுங்காம

முற்பயக்குஞ்சின்னீரவின்பத்தைமுற்றிழாய்‌

பிற்பயக்கும்பீழைபெரிது.
\end{quotation}
\begin{quotation}\small
\emph{Learning at first painful, will afterwards afford pleasure. It will destroy ignorance and extend knowledge. But the pain which succeeds to the short-lived pleasure of immoderate lust, is great, O thou adorned with perfect jewels!}
\end{quotation}
\begin{quotation}\small
எத்துணையவாயினுங்‌ கல்வியிடமறிந்

துய்த்துணர்வில்லெனினில்லாகு --- முய்த்துணர்த்துஞ்

சொல்வன்மையின்றெனினென்னாமஃதுண்டேற்‌

பொன்மலர்நாற்றமுடைத்து.
\end{quotation}
\begin{quotation}\small
\emph{Learning, however extensive, will be useless, unless there be discretion to display it in the proper place; and even then, without the power of language, of what avail is it? With that, it is a flower of gold that possesses fragrance.}
\end{quotation}
\begin{quotation}\small
கற்றார்க்குக்கல்விநலனேகலனல்லான்‌

மற்றோரணிகலம்வேண்டாவா --- முற்ற

முழுமணிப்பூணுக்குப்‌ பூண்வேண்டாயாரே

யழகுக்கழகுசெய்வார்‌.
\end{quotation}
\begin{quotation}\small
\emph{The learned need no other ornament than the excellence of learning. Nothing is wanting to adorn an ornament perfectly set with every precious stone. Who would beautify beauty itself!}
\end{quotation}

\section{Kum\={a}rasinga Mudeliar. --- குமாரசிங்கமுதலியார்.}
\paragraph{}
Kum\={a}rasinga Mudeliar was descended from a respectable \emph{Vel\={a}la} family at \emph{N\={a}vetkuli}, in Mantotte, and enjoyed some reputation both as a clever poet, and an experienced physician. Having unfortunately taken an active part in the tumult which occurred in the district of Manaar, in consequence of the imposition of a tax upon ornaments by Governor \textbf{North} in \textsc{ce} 1800, he was tried for sedition, and being convicted, was sentenced to be whipped. There is a very pathetic address to the \textbf{Blessed Virgin}, in a series of stanzas, which he repeated impromptu while being whipped, craving her intercession that he may be endowed with fortitude to bear the affliction. He was also the author of several hymns now sung in the Catholic churches in Mantotte, at Novenas, and which exhibit his strong devotional feeling.

\section{K\={u}langkaiyer. --- கூழங்கையர்.}
\paragraph{}
K\={u}langkaiyer was a highly celebrated poet and philosopher, well skilled both in the Sanskrit and Tamil languages. No less a place of sanctity amongst the Hindus than \emph{Kanjipuram} claims him as its native, and his parents it would appear traced their lineage to the \emph{Velalers}, who colonized \emph{Tondamandalam} under the auspices of \textbf{Adonda Chakravarti}, before the commencement of the Christian era. He was educated at first in a country school, where he gave early indications of his superior genius, and afterwards in the \emph{Saiva Matam} at Tiruvatur, in the Tanjore country, under the tuition of the learned \emph{Tambirans} or ascetics, who resided in it, and whose profession it was to assist the youth in the prosecution of the study of the Hindu system of Metaphysics and Theology. Having distinguished himself while a student by his various attainments, and earned for himself the title of \emph{Tambiran}, he was selected in the first instance for the office of Junior teacher in the \emph{Matam}; the duties of which he performed in an admirable manner, evincing a great solicitude to promote the interests of his pupils. The superior of the \emph{Matam} soon becoming jealous of the young \emph{Tambiran}, who from his zeal and learning seemed to bear away the palm of praise, falsely accused him of some crime, in consequence of which he was required to take up a piece of red-hot iron, to prove his innocence. He submitted to the ordeal; but in its performance he was unfortunately deprived of the use of his right hand, and henceforth he passed under the sobriquet of \emph{K\={u}langkaiyer}, (one with a crippled hand), and his real name became forgotten by disuse. Subsequently, he came to Jaffna, and settled there under the kindness and hospitality of \textbf{Vytilinga Chettiar}, the principal merchant in the place, and it appears that teaching was the sole occupation in which he engaged himself for the remainder of his life. He not only taught the children of his benefactor, but likewise of many others who had sought his assistance, in order to secure for them a classical education. The fame of his deep erudition soon spread over to Colombo, from which place he received several invitations from various individuals to come and help them in their studies. With his characteristic desire to diffuse knowledge, he accepted their invitations, and paid several visits to Colombo, where he always met with a warm reception from all parties. To his guidance was the Reverend Father \textbf{Gabriel Pacheco} chiefly indebted for that thorough knowledge of Tamil, which afterwards enabled him to produce in that language those valuable religious works, especially the History of the People of God (தேவப்பிரசையின்திருக்கதை), in two ponderous volumes, whereby his name has been eminently distinguished.

K\={u}langkaiyer was the author of an exposition of the \emph{Nann\={u}l}, and he also composed some epic poems, and several minor pieces of poetry, remarkable for their elegant and harmonious versification. Amongst his epics may be mentioned \emph{Yosep Pur\={a}nam} (யோசேப்புராணம்), which treats of the history of the patriarch \textbf{Joseph}, in 1023 stanzas, arranged in twenty-one cantos. He inscribed it to the learned \textbf{De Melho}, who was his friend and literary associate, as a tribute of respect to his eminent worth and abilities; but unfortunately only a fragment of it is now extant. As a specimen of its poetry, we subjoin the five following stanzas, extracted from the second canto entitled \emph{\={A}ttupadalam}:
\begin{quotation}\small
பயம்புவிக்கருள்‌ பயோதரமியாவும்வெண்டிரை

பயங்கொள்வேலையிற்‌ பரந்துவாரிதியிடைபடிந்து

பயங்களானவை பருகியப்பரவையின்வடிவாய்ப்‌

பயங்கள்கோரகை கொண்டிடப்பரந்ததம்பரமேல்‌.

பரவுமால்ககுபங்களோரெட்டொடுவானைப்‌

புரவுவேந்தராள்‌ புவியினைப்போர்த்தெனவிருண்டு

விரவுகின்றது விண்ணெனும்பந்தரின்கண்ணே

யிரவுநேர்கரும்‌ படாத்தினால்‌ விதானஞ்செய்தென்ன.

திருவுலாங்கலைப்‌ பிலிப்புமேலோன்குருச்சிந்தை

மருவுஞானநல்லொளியென மின்னிமற்றொன்வாய்ச்‌

சுருதியாமெனத் தொனித்தின்னோனுயிர்க்கருடொலைமாக்‌

கருணையாமெனப்பொழிந்தது காரிருட்புயலே.

தேவதேவநம்பராபரன்‌ சீரணிசிறந்த

தாவிலாத்திருநாமமே தயங்குறுஞ்சுத்த

மேவிடுஞ்சியோன்வரைக்குமஞ்சனம்புரிவிதம்போ

லோவிலாமழை பெய்தனவுலகுளோருவப்ப.

இறும்பினுச்சியிற்‌ பெய்திடும்புனலெல்லாமிழிந்தே

யெறும்புசீயம்புல்லிருங்குடாவடியெலாமீர்த்து

நறும்பணைத்தரளங்களைக்கொழித்துடனடந்து

வறும்புனத்திடையுலாயதுகடல்புரைவாவி.
\end{quotation}
\begin{quotation}\small
\emph{The clouds which supply the earth with water, having overspread, and settled upon the sea abounding with silvery waves, imbibed the water of the sea, and assuming the form of the sea itself spread over the sky.}

\emph{When the darksome clouds rapidly moved along the sky, extending themselves to the eight points of the compass, and covering over the earth ruled by kings, they resembled a black canopy hanging as it were beneath the pavilion of heaven.}

\emph{The clouds darted lightnings, bright as the light of wisdom in the mind of the learned teacher Philip De Melho; sounded as the words of instruction proceeding from his mouth; and poured down rain with a benevolence like his to promote the happiness of mankind.}

\emph{Gladdening the inhabitants of the earth, the rains poured down incessantly; and it seemed as if they were bathing the mount Sion, whereon the holy name of our God, the God of gods is peculiarly distinguished.}

\emph{The waters falling on the summit of the mount descended along its sides, and having filled both the holes of ants, and the caves of lions tigers and bears, pursued their course towards the sea, traversing the hilly country, and carrying with them the gems which they had collected together.}
\end{quotation}
\paragraph{}
During the latter part of his life, K\={u}langkaiyer chiefly resided at Chiviateru, in the parish of Chundikuly; and he died at the same place advanced far in years, a short time after the surrender of Jaffna to the British. He was buried at Tinnevelly, in the parish of Nellore, and the Rev. Mr. Roberts in the ``Oriental Illustrations of the Sacred Scriptures,'' p. 205, notices his sepulture, thus: ``Some of those who are reported to be very holy are buried in salt, and in a sitting posture, so that they are considered to be still at their devotions. Thus \emph{Koonakaitambaran} (K\={u}langkai Tambiran), of Jaffna was buried.''

K\={u}langkaiyer was free from prejudice, liberal in his views, friendly to all, and identified himself with no party, hence he was universally respected and admired by people of every class and creed. His constant intercourse with De Melho enabled him to acquire a knowledge of the Bible, and whenever so disposed he was ready to argue very ably and warmly in favour of the divine origin of Christianity, which excited the astonishment of the votaries of Hinduism. It is said that when at the dedication of the \emph{Saiva} temple at Vannarponne, which was built and endowed by \textbf{Vytilinga Chettiar}; he went to witness the solemnities, the officiating \emph{Br\={a}hman} presented him with the holy ashes to be rubbed on the forehead; but he took and rubbed it on his slippers, observing at the same time ``\emph{let the like join the like},'' that is, that as the one was made of cow-dung, and the other of cow's hide, let them be joined together. We, however, much regret that, convinced as he was of the truths of Christianity and the falsity of Hinduism, and treating the one with respect, and the other with contempt, he should nevertheless have lived and died a Hindu.

It should be mentioned that K\={u}langkaiyer, after his arrival and settlement in Jaffna, directed his attention to the study of the Dutch, and Portuguese languages, and made himself so far acquainted with them as to be able to speak them with tolerable fluency, and obviated the necessity of an interpreter in his intercourse with the Dutch \emph{Commandeur}, and other high functionaries in the place, whom he visited, and who always received him with civility and courteousness.\footnote{We are indebted for the particulars of the life of K\={u}langkaiyer and those of De Melho to the kindness of M. P. Jurgen Ondatji, Esquire of the Colonial Office, Colombo.}

\section{Lorenzo Pulaver. --- லோரெஞ்சுப்புலவர்.}
\paragraph{}
Lorenzo Pulaver was born of \emph{Vel\={a}la} parents at \emph{P\={a}sikulam}, in Mantotte, and brought up in the Catholic religion which his family professed from the time of its introduction into Manaar by St. \textbf{Francis Xavier}, the Apostle of India. He lived in the middle of the last century, and though he did not write any large poems, yet his abilities are manifest from the many excellent hymns which he composed in honour of our \textbf{Saviour}, and the \textbf{Blessed Virgin}. It is said that while he was on a visit to Putlam, \textbf{Allah Pitchei Annaviy\={a}r}, a Moorish poet of some repute, who lived there, endeavoured to bewilder him in a contest respecting the \emph{Kanam} (கணம்), or foot of an invocation in a poem, by producing a forged rule in support of his false proposition; but he soon exposed the forgery, and carried away the palm in the contest. He was remarkably quick at making repartees. One day a poet from Colombo, meeting him for the first time near the church at Mantotte, and wishing to be facetious at his expense, asked him, ``\emph{Is the Mantotte buffalo very big?}'' (மாதோட்டத்தான் கடா மெத்தப்பெரிதோ) and the ready reply was, ``\emph{Yes it is little bigger than the Colombo ass}'' (கொழம்பான் கோவேறுகழதையிற் கொஞ்சம்பெரிது). \emph{Notandum est}, in Tamil both the words ``\emph{buffalo}'' and ``\emph{ass}'' are metaphorically applied for a blockhead.

\section{Machamuni. --- மச்சமுனி.}
\paragraph{}
Machamuni, who was both a distinguished poet and philosopher, lived at the same time with \textbf{Agastiyer}; but nothing further is known of him. None of his compositions had escaped the ravages of time, save the \emph{V\={a}da Negandu} (வாதநிகண்டு), a treatise on the nostrums and mysteries of alchemy, which science now justly held in contempt, appears to have been once much cultivated among the Tamils and an anecdote is still current of a merchant in the Carnatic, who after having spent all his wealth in trying to make gold, was reduced to beggary, and used to exclaim as he went about begging (வங்கமிறுகுதில்லை லிங்கமுருகுதில்லை) ``\emph{the lead would not be hardened and the cinnabar would not melt}.''

\section{Madura P\={a}l\={a}siriyan\={a}r. --- மதுரைப்பாலாசிரியனார்.}
\paragraph{}
This poet must have been a man of great erudition, as may be inferred from the fact of his having had the title of \emph{\={A}siriyan} or Doctor conferred on him, by his contemporaries. His name occurs among the forty-nine professors of the Madura College but none of his compositions has been preserved, save the following impromptu on the \emph{Kural}.
\begin{quotation}\small
வெள்ளிவியாழம்‌ விளங்கிரவிவெண்டிங்கள்‌

பொள்ளெனநீக்கும்‌ புறவிருளைத் --- தெள்ளிய

வள்ளுவரின்குறள்வெண்பாவகிலத்தோ

ருள்ளிருணீக்குமொளி.
\end{quotation}
\begin{quotation}\small
\emph{The planets Venus, Jupiter, Sun and Moon would swiftly dispel the darkness from the face of the earth; but the Kural of the learned Valluver would more swiftly dispel the darkness from the hearts of men.}
\end{quotation}

\section{Madura Perumaruthan\={a}r. --- மதுரைப்பெருமருதனார்.}
\paragraph{}
This poet was one of the forty-nine professors of the Madura College. Nothing further is known of him beyond that he uttered the following impromptu on the \emph{Kural}, which contains merely an enumeration of the chapters of that work, with an eulogy on \textbf{Tiruvalluver}, for having imparted through them the sublime truths of the \emph{Vedas}, and redeemed the world from its errors:---
\begin{quotation}\small
அறமுப்பத்தெட்டுப்பொருளெழுபதின்பத்‌

திறமிருபத்தைந்தாற்றெளிய --- முறைமையால்‌

வேதவிழுப்பொருளை வெண்குறளால்‌ வள்ளுவனா

ரோதவழுக்கற்றதுலகு.
\end{quotation}

\section{Madura Tamil-n\={a}yagan\={a}r. --- மதுரைத்தமிழ்நாயகனார்.}
\paragraph{}
This poet was one of the forty-nine professors of the Madura College. His real name is not known, and the one under which he is noticed is merely titular, referring to his high attainments in Tamil literature. When the \emph{Kural} was submitted to the judgment of the College Bench, he expressed his approbation of that work in the following impromptu:---
\begin{quotation}\small
எல்லாப்பொருளுமிதன்பாலுளவிதன்பா

லில்லாதவெப்பொருளுமில்லையாற்‌ --- சொல்லாற்

பரந்தபாவாலென்பயன்வள்ளுவனார்‌

சுரந்தபாவையத்துணை.
\end{quotation}
\begin{quotation}\small
\emph{What is the use of works of great length, when the short work of Valluver alone is enough to edify the world? It contains all things and there is nothing which it does not contain.}
\end{quotation}

\section{Manav\={a}lad\={a}ser. --- மனவாளதாசர்.}
\paragraph{}
Manav\={a}lad\={a}ser was a \emph{Br\={a}hman} poet of the \emph{Vaishnava} sect. His real name is said to have been \emph{Pilleiperum\={a}l Aiyankar}. He wrote a poem under the title of \emph{Tiruv\={e}nkata M\={a}lei} (திருவேங்கடமாலை), consisting of a centum of verses in praise of the shrine of \emph{Vishnu} at Tirupati. The first fifty verses of this poem is plain enough to be understood, but the last fifty have been purposely rendered obscure by the use of words which are susceptible of different interpretations.

\section{Mandalapuruder. --- மண்டலபுருடர்.}
\paragraph{}
Mandalapuruder was a \emph{Jaina} ascetic, equally eminent as a poet and lexicographer. He is commonly called \emph{Vira Mandalaver}, but why so is not known. His poetic lexicon, entitled \emph{S\={u}d\={a}mani Nigandu} (சூடாமணிநிகண்டு), outshines all similar productions in Tamil, and its execution certainly redounds to his credit. This work, according to his own averment, was composed by him at the suggestion of \textbf{Gunapattiren}, his spiritual preceptor, and it consists of 989 stanzas, arranged in twelve chapters, and treats of the synonyms of the Hindu Deities, and of the objects of the Animal, Vegetable and Mineral kingdoms as well as of the homonymous and generic terms in nearly the same order with the \emph{Amera Kosha}. He was also the author of an epic poem in honour of \emph{Arha}, the God of the \emph{Jaina} sect. He appears to have lived in the time of the king \textbf{Krishna R\={a}yer}, who reigned at Vijayanagaram from \textsc{ce} 1508 to \textsc{ce} 1530, and whom he celebrates in the 10th stanza of the ninth chapter of his \emph{Nigandu}.

\section{Maruttuven Tamotharanar. --- மருத்துவன்தாமோதரனார்.}
\paragraph{}
This poet was one of the forty-nine professors of the Madura College, but nothing further is known of him than that he was also well skilled in the medical art. His encomiastic impromptu on the \emph{Kural}, which we subjoin, is very curious, in as much as it contains at the same time a recipe for head-ache.
\begin{quotation}\small
சீந்திநீர்க்கண்டந் தெறிசுக்குத்தேனளாய்‌

மோந்தபின்யார்க்குந் தலைக்குத்தில்‌ --- காந்தி

மலைக்குத்துமால்யானை வள்ளுவர்முப்பாலாற்‌

தலைக்குத்துத்தீர்வுசாத்தற்கு.
\end{quotation}
\begin{quotation}\small
\emph{O king, who resembles the huge elephant which thrusts its tusk at the mountain (taking it for an enemy)! All are relieved of their head-ache by smelling the Sindil salt,\footnote{A kind of salt produced from the stalk of the \emph{Menispermum cordifolium}.} and sliced dry-ginger mixed with honey; but S\={a}tanar was relieved of his own head-ache by hearing the three parts of the Kural recited.}\footnote{See the Life of Sitaleichatanar [Chithalai Chathanar].}
\end{quotation}

\section{M\={a}m\={u}lan\={a}r. --- மாமூலனார்.}
\paragraph{}
This poet was one of the professors of the Madura College, and his memory has been preserved only by the following impromptu which he uttered on the \emph{Kural} and its author:---
\begin{quotation}\small
அறம்பொருளின்பம்‌ வீடென்னுமந்நான்கின்‌

றிறந்தெரிந்துசெப்பியதேவை --- மறந்தேயும்‌

வள்ளுவனென்பானோர்‌ பேதையவன்வாய்ச்சொற்‌

கொள்ளாரறிவுடையார்‌.
\end{quotation}
\begin{quotation}\small
\emph{Valluver who after studying the four topics (of the moral philosophy, viz: Virtue, Wealth, Pleasure and Eternal happiness) has imparted their substance (through the medium of the Kural), is in reality a god, and if any shall so far forget his divine character, and say that he is a mere mortal, not only will the learned reject his saying but likewise take him for an ignorant man.}
\end{quotation}

\section{M\={a}ngkudi Marudan\={a}r. --- மாங்குடிமருதனார்.}
\paragraph{}
Nothing more is known of this poet than that he was one of the forty-nine professors of the Madura College. When the \emph{Kural} passed the College bench he commended it in the following impromptu:---
\begin{quotation}\small
ஒதற்கெளிதாயுணர்தற்கரிதாகி

வேதப்பொருளாய்‌ மிகவிளங்கித் --- தீதற்றோ

ருள்ளுதோறுள்ளுதோறுள்ளமுாருக்குமே

வள்ளுவர்வாய்மொழிமாண்பு.
\end{quotation}
\begin{quotation}\small
\emph{The beauty of Valluver's Kural is that it not only illustrates the abstruse doctrines of the Vedas; but is itself a Veda, easy to be studied, and having the effect of melting the hearts of the righteous who study it.}
\end{quotation}

\section{M\={a}nikav\={a}sagar. --- மாணிக்கவாசகர்.}
\paragraph{}
M\={a}nikav\={a}sagar was a \emph{Saiva} devotee and poet of considerable eminence, who flourished during the reign of the king \textbf{Arim\={a}rta P\={a}ndiyen}, which some place between the fifth and eighth century after the Christian era. He belonged to a \emph{Br\={a}hman} family of the race of \textbf{Am\={a}tiyer} at \emph{V\={a}dav\={u}r}, a town on the banks of the \emph{Vaigai}; and when we strip off the marvellous with which the particulars of his life have been blended in the \emph{Tiruv\={a}dav\={u}rer Pur\={a}nam}, it would appear that early in life his superior learning, and qualifications having attracted the notice of the king he was summoned to the Court, and appointed Minister of state, under the title of \emph{Tenneven Piramar\={a}yen}, but sometime afterwards when sent on a commission to buy horses for the king, yielding to a delusion, he expended the money with which he was entrusted in founding and endowing a \emph{Saiva} temple at Tiruperunturei, and consequently incurred the royal displeasure and was put in irons and cast into prison. At this juncture the \emph{Vaigai} river happened to overflow its banks, threatening destruction to the capital, and the king considering it as a retributive punishment sent by \emph{Siva} for ill-treating M\={a}nikavasagar, set him at liberty. He thereupon assumed the ascetic life and after making pilgrimages to other \emph{Saiva} shrines, established himself at Chillambaram, where he died in the thirty-second year of his age. While he was residing at Chillambaram, he distinguished himself by defeating the \emph{Buddhist} priests, who went over there from Ceylon, and challenged the \emph{Br\={a}hmans} to a disputation, as related in the sixth Chapter of the \emph{Tiruv\={a}dav\={u}rer Pur\={a}nam}, of which we have published an English translation in the second number of the Journal of the Ceylon Branch of the Royal Asiatic Society. He is known as the author of a Tamil Grammar, under the title of \emph{Kuvaleiy\={a}nandam} (குவலையானந்தம்), besides two other works respectively called \emph{Sittambala-k\={o}vei} (சிற்றம்பலக்கோவை), and \emph{Tiruv\={a}sagam} (திருவாசகம்). The \emph{Sittambala-kovei}, consisting of 400 stanzas, is an epithalamium in which the mutual passion and love of \emph{Siva} and his consort \emph{P\={a}rvati} are described with glowing imagery, and the Hindus consider it an allegorical poem capable of a spiritual interpretation. The \emph{Tiruvasagam} contains a series of hymns addressed to \emph{Siva}, which the Hindus esteem highly as being most affecting; hence the proverb, (திருவாசகத்திலுருகார் ஒருவாசகத்திலுமுருகார்). Those who will not be moved by the \emph{Tiruv\={a}sagam} (the holy word) will not be moved by any other \emph{V\={a}sagam} (word). The hymns are certainly very impressive, as may be judged by the specimen which Mr. \textbf{Ellis} has quoted in his Commentary on the \emph{Kural}.

\section{M\={a}rgasah\={a}ya D\={e}ver. --- மார்க்கசகாயதேவர்.}
\paragraph{}
M\={a}rgasah\={a}ya D\={e}ver is said to have been born a poet, and hence he had the title of \emph{Vara Kavi}; but we are ignorant both of the place of his birth, and the time of his existence. Of his compositions we have met with only a poem of the class called \emph{Pillei tamil} (பிள்ளைத்தமிழ்), describing the gestures and amusements of the childhood of \emph{Skanda} in the character of \emph{Muruga}, as worshipped at Tiruvirinchei, and which has acquired some repute amongst the \emph{Saivas}.

\section{Meikanda D\={e}ver. --- மெய்கண்டதேவர்.}
\paragraph{}
Meikanda D\={e}ver, whose proper name was \emph{Svetavanaperumal}, claimed his descent from a \emph{Vel\={a}la} family at \emph{Vennei-nell\={u}r}, a town situated on the river \emph{Pennei}, in the Carnatic. He is reputed to have been a pupil of \textbf{Paranjoti Tambiran}, who flourished during the eleventh century of the Christian era. He was an able poet, well studied in the Hindu metaphysics and theology. Of his various works the most admired is the \emph{Sivagnana P\={o}tham} (சிவஞானபோதம்), a metaphysical and theological treatise, which has recently been translated into English by the Rev. \textbf{H. R. Hoisington}, and published in the fourth volume of the Journal of the American Oriental Society, ``as worthy of notice and of preservation'' covering as it does nearly the whole field of Hindu philosophy.

\section{M\={o}sikiranar. --- மோசிகீரனார்.}
\paragraph{}
This poet was one of the forty-nine professors of the Madura College; but there is nothing remaining of his compositions save the following impromptu on the \emph{Kural} merely describing the divisions of that work.
\begin{quotation}\small
ஆண்பாலேழாறிரண்டுபெண்பாலடுத்தன்பு

பூண்‌பாலிருபாலோராறாக --- மாண்பாய்‌

காமத்தின்பக்கமொரு மூன்றாகக்கட்டுரைத்தார்‌

நாமத்தின்வள்ளுவனார்நன்கு.
\end{quotation}

\section{Muttus\={a}mi Pillei. --- முத்துச்சாமிபிள்ளை.}
\paragraph{}
Muttus\={a}mi Pillei was born of Catholic parents who belonged to a respectable class of \emph{Vel\={a}las} at Pondicherry. Having early distinguished himself as a Tamil poet, and at the same time made himself familiar with the Sanskrit, Telugu, Latin and English languages, he obtained the offices of Manager of the College of Fort St. George and Moonshee to the Tamil Translator to Government at Madras; both which he held until his death. In 1816 at the desire of Mr. \textbf{Ellis}, the Senior member of the College Board, he made a journey to the south for the purpose of procuring useful Tamil works for the College; and he availed himself of the opportunity to visit the different places in the south which had been the scene of the labours of Father \textbf{Beschi} and collect materials for a brief history of the life of that celebrated Missionary, which he published both in Tamil and in English. Some time afterwards a \emph{Chanda} poet, named \textbf{Ponnambalam}, who lived at Pursew\={a}kam, having published a blasphemous lampoon on Christianity, which made a great sensation amongst both Catholics and Protestants, Muttusami Pillei undertook to refute it; and accordingly composed a very able work, under the title of \emph{Dhikk\={a}ram} (திக்காரம்), partly in poetry, and partly in prose; and when he had it recited before the assembly of learned men at Madras, it was not only received with great applause; but also obtained for him the present of a breast-plate set with precious stones, and a suit of Surat shawls; and the Pundits of the College each addressed to him a stanza in Sanskrit as well as in Tamil, commending his production for the sublimity of its language and sense. Besides the two works noticed above, he wrote a commentary on Mr. \textbf{Ellis'} stanzas, called \emph{Taravu Kochchagakalipp\={a}} (தரவுகொச்சகக்கலிப்பா); this he did in order to refute an idea which was prevalent amongst the Hindus that Mr. \textbf{Ellis} having ended each of his stanzas with the words \emph{Namasiv\={a}ya}, he had therefore become a convert to their religion: by proving that these words were never intended to represent the pentagrammaton, but only to convey the meaning ``reverence to the only God.'' He died on the 23d of September 1840.

\section{Muttu T\={a}ndaver. --- முத்துத்தாண்டவர்.}
\paragraph{}
Muttu T\={a}ndaver was a native of \emph{Shiyally}, near Chillambaram. His parents were both musicians of the \emph{Saiva} temple at \emph{Shiyally}; but a disease with which he was afflicted in his youth disabled him from following their profession. He however, it is said, became a poet by inspiration, and was all his life employed in chanting the praises of \emph{Siva} in the temple at Chillambaram, receiving daily five gold fanams for his services. He used to chant a song every day, beginning it with the word which first caught his ears as he entered the temple. It appears that on a certain day as he was chanting a song in the temple, he died suddenly. The \emph{Saivas} however would have it that whilst he chanted the song beginning with the chorus (மாணிக்கவாசகர் பேறெனக்குத்தரவல்லாயோவறியேன்) ``\emph{I know not whether thou art able to give me the same bliss which thou didst give to Manikav\={a}sagar or no},'' he was absorbed in the glory which appeared in the sanctuary.

All his chants have been carefully preserved and they form the best recitative poem in the Tamil language, but some of them abound in licentious imagery.

\section{Nachin\={a}rkiniyar. --- நச்சினார்க்கினியார்.}
\paragraph{}
Of the life of this poet we have no account; but he appears to have been a man of considerable attainments. His commentaries on the \emph{Tolk\={a}ppiyam} and \emph{Tirumurug\={a}ttupadei} are much esteemed, and they are certainly masterly productions of a logical mind. The exact period of his existence is very uncertain; but we think we shall not be far from the point in placing it before the tenth century of the Christian era.

\section{Nachuman\={a}r. --- நச்சுமனார்.}
\paragraph{}
Nachuman\={a}r is known as one of the forty-nine professors of the Madura College; but no particulars of his history are forthcoming. The subjoined impromptu on the merits of the \emph{Kural} is ascribed to him:---
\begin{quotation}\small
எழுத்தசைசீரடி சொற்பொருள்யாப்பு

வழுக்கில்வனப்பணிவண்ண --- மிழுக்கின்றி

யென்றெவர்செய்தனவெல்லாமியம்பின

வின்றிவரின்குறள்வெண்பா.
\end{quotation}
\begin{quotation}\small
\emph{When and by whom was a poem ever composed so entirely faultless both in versification and sense as the Kural which he (Tiruvalluver) has this day recited?}
\end{quotation}

\section{Nalk\={u}r-v\={e}lviy\={a}r. --- நல்கூர்வேள்வியார்.}
\paragraph{}
This poet claims notice as one of the forty-nine professors of the Madura College, and it would appear from his surname, which is a contracted form of the epithet \emph{V\={e}lviy\={a}ler}, usually applied to the \emph{Br\={a}hmans}, that he must have belonged to their class. His compositions are not now forthcoming, with the exception of the following impromptu in praise of \textbf{Tiruvalluver}.
\begin{quotation}\small
உப்பக்கநோக்கியுபகேசிதோண்‌மணந்தா

னுத்தரமாமதுரைக்கச்சென்ப --- விப்பக்க

மாதானுபங்கிமறுவில்புலச்செந்நாப்‌

போதார்புனற்கூடற்கச்சு.
\end{quotation}
\begin{quotation}\small
\emph{They say, that he (Siva), who reclines on the shoulders of Upakesi (P\={a}rvati), is the patron of North Madura whither he went (with Iddeik\={a}dar); but he (Tiruvalluver), the poet who pours out instruction in honeyed words with a parental solicitude, is the patron of South Madura abounding with water.}
\end{quotation}

\section{Nall\={a}pillei. --- நல்லாபிள்ளை.}
\paragraph{}
Nall\={a}pillei was born of a \emph{Karuniker} family at \emph{Muthalampedu}, in the Carnatic. Such was the precocity of his talents, that before he had passed his sixteenth year he mastered the Tamil, Sanskrit and Telugu languages and their literature, and obtained a distinguished rank among the poets of his time. He first composed an epic poem, entitled \emph{Deivay\={a}nei Pur\={a}nam} (தெய்வயானைபுராணம்), in 3000 stanzas, and then set about enlarging \textbf{Villiputturer's} imitation of the \emph{Mahabh\={a}rat} by adding to its 3372 stanzas, 14728 more. The latter work occupied him from \textsc{ce} 1732 to 1744, and he survived its completion but one year.

\section{Nall\={a}than\={a}r. --- நல்லாதனார்.}
\paragraph{}
Nall\={a}than\={a}r was a poet, who flourished during the Augustan age of Tamil literature. Though we have no records of the particulars of his life; yet a centum of epigrams, which he has left behind him, prove that his poetical talents were far above mediocrity. His epigrams are collectively called \emph{Trikadugam} (திரிகடுகம்), and are said to have had the unanimous approval of the professors of the Madura College. Each epigram contains a comparison of three things and their result: as,
\begin{quotation}\small
கொல்யானைக்கோடுங்‌ குணமிலியுமெல்லிற்

பிறன்கடைநின்றொழுகுவானு --- மறந்தெரியா

தாடும்பாம்பாட்டுமறிவிலியுமிம்மூவர்‌

நாடுங்காற்றூங்குபவர்‌.
\end{quotation}
\begin{quotation}\small
\emph{He who runs against the killing elephant;}

\emph{He who stands outside his neighbour's door in the nighttime;}

\emph{He who causes the snake to dance;}

\emph{These three do so at their peril.}
\end{quotation}

\section{Nambiy\={a}nd\={a}r Nambi. --- நம்பியாண்டார்நம்பி.}
\paragraph{}
Nambiy\={a}nd\={a}r Nambi was a poet, who flourished amongst the \emph{Br\={a}hmans} of the \emph{Adi Saiva} sect at \emph{N\={a}reiyur}, in the Carnatic, during the reign of the king \textbf{R\={a}jar\={a}ja-abhaya Chola}, under whose patronage he composed a poem, entitled \emph{Kalitturei Ant\={a}di} (கலித்துறையந்தாதி), treating of the legends of the sixty-three special votaries of \emph{Siva}; and which on being recited before the king and the learned assembly of \emph{Br\={a}hmans} was very much applauded as a masterly performance. He was an enthusiastic admirer of \textbf{Sampanter}, and wrote two poems, one entitled \emph{Mummanikk\={o}vei} (மும்மணிக்கோவை), and the other \emph{Sanbeiviruttam} (சண்பைவிருத்தம்), in praise of that ascetic.

\section{Nam\={a}zhv\={a}r. --- நம்மாழ்வார்.}
\paragraph{}
Nam \={A}zhv\={a}r deserves notice as a poet as well as one of the twelve \emph{\={A}zhvars} or special votaries of \emph{Vishnu}. He was born of a \emph{Vel\={a}la} family at \emph{Tirukurugai}, on the \emph{T\={a}mbraparni} river, and his genius having displayed itself while yet a child to the astonishment of his parents, they not being able to account for it otherwise persuaded themselves into the belief that he was an incarnation of \textbf{Sena Mudellar}, one of their saints. He was the author of four different poems in honour of \emph{Vishnu} containing in all 1296 stanzas, and they now form a part of the \emph{N\={a}l\={a}yiraprabhandam} (நாலாயிரபிரபந்தம்), which the \emph{Vaishnavas} regard next to the \emph{Vedas}.

\section{Napp\={a}lattan\={a}r. --- நப்பாலத்தனார்.}
\paragraph{}
Napp\={a}lattan\={a}r was one of the forty-nine professors of the Madura College, and we know nothing more about him. His impromptu on the \emph{Kural} is as follows:---
\begin{quotation}\small
அறந்தகளியான்றபொருடிரியின்பு

சிறந்தநெய்செஞ்சொற்றீத்தண்டு --- குறும்பாவா

வள்ளுவனாரேற்றினார்வையத்துவாழ்வார்க

ளுள்ளிருணீக்கும்விளக்கு.
\end{quotation}
\begin{quotation}\small
\emph{Valluver has lighted a lamp for dispelling the darkness from the hearts of those who live in the world; having virtue for its bowl, wealth for its wick, pleasure for its oil, the fire of expression for its flame, and the short stanza for its stand.}
\end{quotation}

\section{Nariveruttaleiy\={a}r. --- நரிவெருத்தலையார்.}
\paragraph{}
Though this poet occupied an honourable place among the forty-nine professors of the Madura College; yet we have neither the particulars of his life, nor are we aware of the existence of any of his compositions. He is, however, reported to have been the author of an impromptu on the \emph{Kural}, which we give below.
\begin{quotation}\small
இன்பம்பொருளறம்வீடென்னுமிந்நான்கு

முன்பறியச்சொன்னமுதுமொழிநூன்‌ --- மன்பதைகட்‌

குள்ளவரிதென்றவை வள்ளுவருலகங்‌

கொள்ளமொழிந்தார்குறள்‌.
\end{quotation}
\begin{quotation}\small
\emph{Virtue, wealth, pleasure, and eternal bliss; these four were obscurely revealed to mankind in the ancient Veda; but Valluver has now rendered them plain in his Kural, that they may be fully known by the world.}
\end{quotation}

\section{Nark\={i}rer. --- நக்கீரர்.}
\paragraph{}
Nark\={i}rer was a poet, who had the enviable distinction of occupying a seat in the Madura College as its head in the ninth century of the Christian era. The incidents of his life as narrated in the \emph{Tiruvilleiy\={a}del Puranam}, are involved in fiction. He, it is stated, having cavilled at a song, which \emph{Siva} himself had composed for a \emph{Brahman}, named \textbf{Tarumi}, \emph{Siva} cast his fiery eye on him in order to reduce him to ashes, and he only escaped from its effects by plunging himself in the sacred pool of \emph{Pottamerei} and chanting the praises of \emph{Siva}. He appears to have belonged to the caste of chank-cutters; for when \emph{Siva}, according to the fable, taunted him with his low birth by observing: ``Is he who cuts the chank able to scan my verse?'' he is represented to have retorted upon the god with the following impromptu:---
\begin{quotation}\small
சங்கறுப்பதெங்கள்குலஞ்‌ சங்கரனார்க்கேதுகுலம்‌

பங்கமறச்சொன்னாற்‌ பழுதாமே --- சங்கை

யரிந்துண்டுவாழ்வோமரனேநின்போல

விரந்துண்டுவாழ்வதிலை.
\end{quotation}
\begin{quotation}\small
\emph{Our caste is certainly that of chank-cutters, but what is the caste of Sankara (Siva)? It would be a disgrace to mention it: O Hara (Siva)! We live indeed by cutting chanks; but not by begging alms like thee.}
\end{quotation}
\paragraph{}
He was the author of an epic, entitled \emph{Sriharni-puranam} (சிறீகருணிப்புராணம்), which he composed at the request of \textbf{Kulachchirein\={a}yan\={a}r}, the prime minister of the king \textbf{K\={u}na Pandiyen}, and likewise of a hymn, entitled \emph{Tirumurug\={a}ttuppadei} (திருமுருகாற்றுப்படை), in honour of \emph{Skanda}.

The subjoined is a stanza which Nark\={i}rer uttered in praise of the \emph{Kural} after it was reviewed and sanctioned by the College bench:---
\begin{quotation}\small
தானேமுழுதுணர்ந்து தண்டமிழின்வெண்குறளா

லானவறமுதலாவந்நான்கு --- மேனோருக்‌

கூழினுரைத்தாற்குமொண்ணீர்முகிலுக்கும்‌

வாழியலகென்னாற்றுமற்று.
\end{quotation}
\begin{quotation}\small
\emph{What recompense can the world make to the poet (Tiruvalluver), who fully understanding the four subjects (virtue, wealth, pleasure and eternal happiness), has explained them in sweet Tamil distichs, that the ignorant may be enlightened, or to the cloud that showers down fertilizing rain?}
\end{quotation}

\section{Nattattan\={a}r. --- நத்தத்தனார்.}
\paragraph{}
This poet was one of the forty-nine professors of the Madura College, and he is known only by the following impromptu which he uttered in praise of the \emph{Kural}.
\begin{quotation}\small
ஆயிரத்துமுந்நூற்றுமுப்பதருங்குறளும்‌

பாயிரத்தினோடு பகர்ந்ததற்பின்‌ --- போயொருத்தர்

வாய்கேட்கநூலுளவோ மன்னுதமிழ்ப்புலவ

ராய்க்கேட்கவீற்றிருக்கலாம்‌.
\end{quotation}
\begin{quotation}\small
\emph{After studying the one thousand three hundred and thirty distichs of the Kural, including the introduction, there is no other work to be studied; for they are enough to make one a perfect poet in Tamil from whom others may seek instruction.}
\end{quotation}

\section{Nattattan\={a}r. --- நற்றத்தனார்.}
\paragraph{}
Nattattan\={a}r was one of the twelve disciples of \textbf{Agastiyer}; but no account of his life has reached the present times. He wrote a treatise on Prosody and Versification, called after his own name \emph{Nattattam} (நற்றற்றம்); but it is now scarcely to be met with.

\section{Nayanappa Mudeliar. --- நயணப்பமுதலியார்.}
\paragraph{}
Nayanappa Mudeliar was born of a \emph{Vel\={a}la} family at Pondicherry in \textsc{ce} 1779, and he acquired the reputation of a poet at an early period of his life. In his eighteenth year he left Pondicherry, and went to Madras, where he found employment as a Tamil Pundit in the College of Fort St. George. While he was thus employed he devoted his leisure hours in revising and editing the poems respectively called \emph{Sittambalak\={o}vei} (சித்தம்பலக்கோவை), \emph{Tanje\={i}v\={a}nan-k\={o}vei} (தஞ்சைவாணன்கோவை), and \emph{Reghun\={a}tha S\={e}tupati Oruturei-k\={o}vei} (இரகுநாதசேதுபதி ஒருதுறைக்கோவை), as well as the \emph{N\={a}ladiy\={a}r} (நாலடியார்), the ninth and tenth parts of the \emph{Div\={a}kara Nigandu} (திவாகரநிகண்டு), and the eleventh part of the \emph{Sud\={a}mani Nigandu} (சூடாமணிநிகண்டு), with their explanations in prose, which were written by himself, and which exhibit much critical acumen. When a Committee was formed at Madras amongst the principal Tamil Gentlemen in the place for the purpose of publishing an edition of \textbf{Villiputturer's} Imitation of the \emph{Mah\={a}bharat}, the Editorship was entrusted to Nayanappa Mudeliar; but whilst he was engaged in collating the different manuscripts of the work he died suddenly on the 25th of January 1845.

\section{N\={a}gand\={e}van\={a}r. --- நாகன்றேவனார்.}
\paragraph{}
This poet is now remembered chiefly on account of his connection with the Madura College, in which he was one among the forty-nine professors. Of his poetical compositions, only the following impromptu on the \emph{Kural} survives:---
\begin{quotation}\small
தாளார்மலர்ப்பொய்கைதாங்குடைவார்தண்ணீரை

வேளாதொழிதல்வியப்பன்று --- வாளாதா

மப்பாலொருபாவை யாய்பவவோவள்ளுவனார்‌

முப்பான்மொழிமூழ்குவார்‌.
\end{quotation}
\begin{quotation}\small
\emph{It is no wonder if those who have bathed in the water of a tank abounding with lotus-fowers will not desire to bathe in any other water; but it is a wonder indeed if they who have read Valluver's work will desire to read any other work.}
\end{quotation}

\section{N\={a}r\={a}yanabh\={a}rati. --- நாராயணபாரதி.}
\paragraph{}
This poet was a \emph{Brahman} who lived at \emph{Venmani}, the exact location of which in the Carnatic we have not been able to ascertain. He is chiefly known as the author of a very popular poem, called \emph{Tiruv\={e}nkata Sadagam} (திருவேங்கடசதகம்), containing a centum of stanzas on moral, and social duties, and which he composed under the patronage of one \textbf{Manav\={a}la Nar\={a}yanen} of \={A}mal\={u}r, and whom he celebrates at the end of every stanza as the favourite of \emph{Tiruvenkata}, the form of \emph{Vishnu} worshipped at Tirupati. The Rev. W. \textbf{Taylor} has translated ten stanzas of this poem into English, and the translation forms one of the Appendices to the second volume of his Oriental Historical Manuscripts.

\section{N\={a}tkavir\={a}ja Nambi. --- நாற்கவிராச நம்பி.}
\paragraph{}
Of the biography of this poet very few particulars have been preserved. All that is known of him is, that he was a native of \emph{Putiankudi}, in the Carnatic, born and bred up in a \emph{Jaina} family; and having made poetry the exclusive object of his study, acquired the high distinction of being considered the king of the poets, and as an improvisator he had scarcely any equal in his day. He has left behind him a treatise containing rules for composing amatory poems, under the title of \emph{Agapporal Illakkanam} (அகப்பொருளிலக்கணம்), which is now generally consulted as a standard authority on the subject.

\section{Ottak\={u}ter. --- ஒட்டக்கூத்தர்.}
\paragraph{}
Ottak\={u}ter was one of the poets who flourished at the court of \textbf{R\={a}jendra Chola}, and that of his son \textbf{Kulotunga Chola}. He excelled all his colleagues in the composition of war chants, called \emph{Parani}. His \emph{Kalingattu Parani} (கலிங்கத்துபபரணி), in which he celebrated the conquest of the \emph{Kalinga} country by \textbf{Kulotunga Chola} is a matchless production of the kind, and it is related that the hero of the poem was so much pleased with it, that he rewarded him by rolling a golden coconut at the end of every stanza, as it was recited; but this is manifestly an oriental exaggeration. The \emph{Uttarak\={a}nda} which forms the seventh book of the \emph{Ramayanam} was composed by Ottak\={u}ter, and was with the approval of \textbf{Kamber} himself incorporated with his own work. This though written in the same metre and style as the other six books, yet on a close inspection an inferiority is perceptible, as much as can be discovered between the first part of \textbf{Dryden's} \emph{Absolom} and \emph{Achitophel}, and the second part of it by \textbf{Tate}. A good deal of rivalry appears to have existed between Ottak\={u}ter and \textbf{Kamber}. On one occasion the king going out on a visit to one of his mistresses, found on his return that the door of the palace was closed against him by the Queen, and that no entreaty or persuasion would make her relent. The king at length sent for Ottak\={u}ter, and desired him to soothe her by his song, and cause her to open the door; and he thereupon addressed her in the following impromptu:
\begin{quotation}\small
நானேயினியுன்னை வேண்டுவதில்லை நளினமொழித்

தேனேகபாடந்‌ திறந்தருள்வாய்‌ திறவாதிடிலோ

வானேழ்புரக்குமிரவிகுலாதிபனின்‌ வாயில்வந்தாற்‌

றானேதிறக்குமுன்கையிதழாகியதாமரையே.
\end{quotation}
\begin{quotation}\small
\emph{O thou whose speech is as sweet as honey! No more shall I entreat thee. Deign to open the door, if not, the arrival at thy door of the king of the solar race who commands the seven clouds, will make thy lotus-like hand open it spontaneously.}
\end{quotation}
\paragraph{}
But this instead of soothing her had only the contrary effect. She bolted the door with an additional bolt, saying, (ஒட்டக்கூத்தன்பாட்டுக்கு இரட்டைத்தாப்பாள்), i. e. ``\emph{Ottak\={u}ter's song is worth two bolts.}'' The king next sent for \textbf{Kamber} who knowing full well the weak point in women flattered her self-conceit in the following impromptu, and obtained his request:---
\begin{quotation}\small
இழையொன்றிரண்டு வகிர்செய்தநுண்ணிடையேந்துசெம்பொற்‌

குழையொன்றிரண்டு விழிமடவீர்‌ கொண்டகோபந்தணி

மழையொன்றிரண்டுபெய்மாணாபரணனின்‌ வாயில்வந்தாற்‌

பிழையொன்றிரண்டு பொறுப்பதல்லோகடன்‌ பூவையற்கே.
\end{quotation}
\begin{quotation}\small
\emph{O thou with the slender waist, and dark blue eyes and ears adorned with golden jewels! Calm thine anger; for when the illustrious sovereign whose country is blessed with rain twice a month, waits at thy door, it is thy duty to forgive him once, nay even twice.}
\end{quotation}

\section{Padik\={a}su. --- படிக்காசு.}
\paragraph{}
Padik\={a}su was a native of \emph{Kalandei}, in the Carnatic. He flourished as a poet at the Court of \textbf{Kilaven Setupati}, otherwise called \textbf{Regunatha Setupati}, who reigned at Ramnad from \textsc{ce} 1686 to 1723. His principal work was a poem, entitled \emph{Tondamandala Sadagam} (தொண்டமண்டலசதகம்), consisting of a hundred stanzas in praise of the \emph{Tonda} country. He appears to have been very expert in the composition of poetical epistles, some of which are still extant. It should, however, be remarked that he was excessively conceited of his own accomplishments, and moreover a severe satirist of his time, calling himself in that respect the representative of \textbf{K\={a}lamegam}.

\section{Paduman\={a}r. --- பதுமனார்.}
\paragraph{}
This poet appears to have flourished at an early period; but no particulars of his life have reached us. He is however well known as the author of a learned commentary on the \emph{N\={a}lady\={a}r}, a poem containing 400 stanzas on moral subjects, and supposed to have been the joint composition of four hundred poets who visited the court of a certain \emph{P\={a}ndiya} king, each of them being the author of a stanza.

\section{Panamb\={a}ran\={a}r. --- பனம்பாரனார்.}
\paragraph{}
Panamb\={a}ran\={a}r, so called from the town \emph{Panamb\={a}ram}, where he lived, was one of the twelve disciples of \textbf{Agastiyer}. Having devoted himself with great assiduity to study, he acquired a high reputation as a grammarian and poet, and wrote an able treatise on Grammar under the title of \emph{Panamb\={a}ran\={a}r S\={u}tram} (பனம்பாரனார்சூத்திரம்). He also wrote the preface to the \emph{Tolk\={a}ppiyam}, a Grammar which his fellow student \textbf{Tolk\={a}ppiyan\={a}r} had composed.

\section{Paraner. --- பரணர்.}
\paragraph{}
Paraner was one of the forty-nine professors of the Madura College and primus inter pares; but scarcely anything is known of his biography. He was present on the College bench when the \emph{Kural} was submitted to its judgement, and is stated to have expressed his opinion of that work in the following impromptu:---
\begin{quotation}\small
மாலுங்குறளாய்‌ வளர்ந்திரண்டுமாணடியான்‌

ஞாலமுழுதுநயந்தளந்தான்‌ --- வாலறிவின்‌

வள்ளுவருந்தங்குறள்‌வெண்பாவடியால்வையத்தா

ருள்ளுவவெல்லாமளந்தாரோந்து.
\end{quotation}
\begin{quotation}\small
\emph{M\={a}l (Vishnu) in his Kural or dwarfish incarnation measured the whole earth, with his two expanded feet; but Valluver has measured the thoughts of all mankind, with his (stanza of) two short feet.}
\end{quotation}

\section{Paranj\={o}ti Tambir\={a}n. --- பரஞ்சோதித்தம்பிரான்.}
\paragraph{}
Paranj\={o}ti Tambir\={a}n, who was a \emph{Saiva} ascetic and superior of the \emph{Saiva Matam} at Madura during the reign of \textbf{Athivira P\={a}ndiyen}, signalized himself by his profound skill in the Sanskrit and Tamil languages and poetry, as well as in the Hindu metaphysical philosophy. At the request of the king, by whom he was patronized, he wrote for him a poem, under the title of \emph{Tiruvilliy\={a}del Pur\={a}nam} (திருவிளையாடற்புராணம்) comprising 3363 stanzas, in 72 cantos. It treats of the sixty-four sports of \emph{Siva} at Madura, as originally related in the Sanskrit legend called \emph{H\={a}la-sya Mahatmya}; and Professor \textbf{Wilson}, in his Historical Sketch of the Kingdom of \emph{P\={a}ndiya}, states that it is reputed to have been written in the \emph{Saka} year 973, corresponding with \textsc{ce} 1051.\footnote{See Journal of the Royal Asiatic Society, Vol. 3. p. 203.} Paranjoti Tambir\={a}n also wrote for himself another poem, entitled \emph{P\={o}tikali venb\={a}}, (போற்றிக்கலிவெண்பா) on the same subject as the above, but on a contracted scale, and in the form of hymns addressed to \emph{Siva}.

An analysis of the \emph{Tiruvilliy\={a}del Pur\={a}nam} has been published by the Rev. \textbf{W. Taylor}, in his Oriental Historical Manuscripts, Vol. 1.

\section{Parimel-azhagar. --- பரிமேலழகர்.}
\paragraph{}
Parimel-azhagar was a \emph{Brahman} poet who owes his celebrity to a Commentary which he wrote on the \emph{Kural}; but nothing is known of his personal history, or the era of his existence. There are nine other Commentaries on the \emph{Kural} written by different poets at different times, but his work has thrown them all in the shade, and fixed itself in popular esteem as the most classical production. Mr. \textbf{Ellis} has availed himself of the aid of this work in framing the valuable notes and illustrations which accompany his metrical version of the first twelve chapters of the \emph{Kural}; and the Rev. Mr. \textbf{Drew} has published a great portion of it with his prose translation of the sixty-three chapters of the \emph{Kural}.

\section{Patirakiriy\={a}r. --- பத்திரகிரியார்.}
\paragraph{}
Patirakiriy\={a}r is generally believed to have been a king, but of what country is not known. He was a contemporary with \textbf{Pattanattupilleiy\={a}r}, and like him abandoning all his worldly possessions became a naked \emph{Sanniy\={a}si}, and begged his bread from door to door, enduring the privation of everything that could in any way have served to gratify his senses. He left a series of 235 couplets, collectively known as his \emph{Pulambel} or laments, and of which we subjoin a few as being illustrative of his religious opinions.
\begin{quotation}\small
என்னையறியாதிருந்தாட்டுஞ்சூத்திரிநின்‌

றன்னையறிந்து தவம்பெறுவதெக்காலம்‌.

ஆதிமுதலாகிநின்ற வரியென்றவட்சரத்தை

யோதியறிந்துள்ளே யுணர்வதினியெக்காலம்‌.

சாத்திரத்தைச்சுட்டுச்‌ சதுமறையைப்பொய்யாக்கிச்‌

சூத்திரத்தைக்கண்டு சுகம்பெறுவதெக்காலம்‌.

உளியிட்டகல்லுமுருப்பிடித்தசெஞ்சாந்தும்‌

புளியிட்டசெம்பும்‌ பொருளாவதெக்காலம்‌.

ஆதிகபிலர்சொன்ன வாகமத்தின்சொற்படியே

சாதிவரையில்லாமற்‌ சஞ்சரிப்பதெக்காலம்‌.
\end{quotation}
\begin{quotation}\small
\emph{Oh! When will the time come that I shall know the secret mover of the machinery within me, and obtain the reward of my penance?}

\emph{Oh! When will the time come that I shall study the mystery of the letter A, which stands the first of all letters, and comprehend its meaning?}

\emph{Oh! When will the time come that I shall burn the Sastras, and prove the four Vedas to be false and he made whole by discovering the mystery?}

\emph{Oh! When will the time come that the chiselled stone, the moulded red clay, and the burnished copper shall be rendered profitable?}

\emph{Oh! When will the time come that men shall live together without any distinction of caste, according to the doctrine promulgated in the beginning by Kapiler?}
\end{quotation}

\section{Pattanattu-pilleiy\={a}r. --- பட்டணத்துப்பிள்ளையார்.}
\paragraph{}
Pattanattu-pilleiy\={a}r was a celebrated philosopher, who lived at \emph{K\={a}veripatnam}, in the Carnatic, some time about the 10th century of the Christian era. He belonged to the \emph{Chetty} caste, and his real name is said to have been \textbf{Venkata Chetty}. He was possessed of great wealth, which he acquired by trading with Ceylon and the neighbouring islands. One day hearing that some of his ships which had been missing for a time returned to the port laden with gold dust, he went to see them, and during his absence, a \emph{Saiva} mendicant called at his house, and asked alms of his wife; but she refused, saying that she could not give any when her husband was not at home. The \emph{Saiva} mendicant thereupon went away, leaving with her a slip of \emph{olah} wrapt in a rag, and requesting her to deliver it to her husband on his return. When Pattanattu-pilleiy\={a}r returned to his house, he looked into the slip of \emph{olah} and finding the words (காதற்றவூசியும் வாராதுகாணுங் கடைவழிக்கே) ``\emph{Mind that even a needle with a broken eye will not follow thee in thy last day},'' he at once imbibed an aversion to the worldly life, gave away all his wealth to the people around, and, abandoning his house, became a naked ascetic, and passed the remainder of his days, subsisting wholly on alms, and esteeming a potsherd and pure gold alike. He had a sister, who, ashamed of his conduct, attempted to poison him, but without success. He latterly took up his abode in a wood at Trivatur, where he caused the cowherds, who resorted to the wood to pasture their cattle, to bury him by day up to the shoulders, leaving only his neck and head above the ground, and take him up at night-fall; but one day having been inadvertently allowed to remain buried during the night, he was found dead the next morning. He appears to have been likewise a poet, as we have a collection of poems which he composed in praise of the different \emph{Saiva} shrines he had visited, besides a number of elegiac verses which he ejaculated extempore at the cremation of his mother's corpse.

Pattanattu-pilleiy\={a}r was a monotheist; but he identified the only one God with \emph{Siva}, and also maintained the doctrine of metempsychosis, teaching that emancipation from future birth, and absorption into the Divine essence was the supreme bliss. He at the same time denied that man was a free-agent, and held that he was incapable of doing either good or evil by himself, as may be seen by the following stanzas which we quote from one of his works:---
\begin{quotation}\small
தவியாதிருநெஞ்சமே தில்லைமேவியசங்கரனைப்‌

புவியார்ந்திருக்கின்ற ஞானாகரனைப்‌ புராந்தகனை

யவியாவிளக்கையென்னம்பலத்தாடியையைந்தெழுத்தாற்‌

செவியாமனீசெவித்தாற்‌ பிறவாமுத்திசித்திக்குமே.

அன்னையெத்தனையெத்தனையன்னையோவப்பனெத்தனையெத்தனையப்பனோ --- பின்னையெத்தனையெத்தனைபெண்டிரோ பிள்ளையெத்தனையெத்தனைபிள்ளையோ --- முன்னையெத்தனை யெத்தனைசென்மமோமூடனாயடியேனுமறிந்திலே --- னின்னமெத்தனை யெத்தனை சென்மமோயென்செய்வேன்‌கச்சியேகம்பநாதனே.

நன்னாரிற்பூட்டிய சூத்திரப்பாவை நன்னார்தப்‌பினாற்

றன்னாலுமாடிச்‌ சலித்திடுமோவந்தத்தன்மையைப்போ

லுன்னாலியானுந் திரிவதல்லான்‌ மற்றுனைப்பிரிந்தா

லென்னாலிங்காவதுண்டோவிறைவாகச்‌சியேகம்பனே.
\end{quotation}
\begin{quotation}\small
\emph{O my heart! do not be desponding. If thou wilt meditate on the five letters\footnote{The five letters of the Hindu Cabala, க, ம, சி, வா, ய.} and devoutly worship Sankara (Siva), who dwells at Tillei, who fills the earth with his wisdom; who destroyed the (rebellious) cities; who is the lamp which is never extinguished; and who dances at the Sacred Court, surely thou shalt obtain liberation from future birth and be absorbed into his divine essence.}

\emph{Oh Ekamba, (Siva) of Kachchi! Can I, an ignorant dog, know how many fathers or how many mothers, or how many wives, or how many children I have had in the course of my former transmigrations, and how many more I shall have in my future?}

\emph{Oh Ekamba, of Kachchi! Will the puppet which is moved by a string, move by itself after the string shall have been broken? Am I not moved by thee in like manner! Therefore can I do anything by myself, if I am separated from thee?}
\end{quotation}

\section{Pavananti. --- பவணந்தி.}
\paragraph{}
The name of Pavananti has been handed down as the author of the \emph{Nann\={u}l} (நன்னூல்), a treatise on Grammar which as Mr. \textbf{Stokely} has justly remarked, ``stands conspicuous among the grammatical treatises of all nations, for logical arrangement and comprehensive brevity;'' but nothing more is known of him than that he was the son of \textbf{Sanmathi Muni}, a \emph{Jaina} sage, who lived at Sanak\={a}puram, and that he wrote his work under the auspices of a king named \textbf{Kangkan}. He had in the \emph{Nann\={u}l} proposed to treat his subject under five heads, viz: (எழத்து) letters, (சொல்) words, (பொருள்) matter, (யாப்பு) versification, and (அணி) embellishment; but having died before he completed his design, the work therefore comprises only the first two heads.

\section{P\={a}napattiran. --- பாணபத்திரன்.}
\paragraph{}
P\={a}napattiran was a minstrel, who lived at the court of \textbf{Varaguna Pandiyen} at Madura. Another minstrel named \textbf{Yeman\={a}then}, who came from the north having exhibited his art before the king, and received rewards, became elated with success and challenged a contest in playing on the lute. P\={a}napattiren accepted the challenge, and by defeating the rival in the contest pleased the king so much that the king honoured him with a triumphant entry of the city on an elephant, and loaded him with presents. Shortly afterwards he incurred the displeasure of the king by purloining some valuables from the palace, and was consequently deprived of all his emoluments and forbidden the royal presence. Finding himself in great straits, he contrived a device to obtain his livelihood for the future. Having been informed that \textbf{Perum\={a}kkothei}, who then reigned over the \emph{Chera} country, was extremely devoted to \emph{Saiva} worship, he forged a poetical letter in the name and character of \emph{Siva}, and took it to the king, saying that it was given to him by \emph{Siva} to be delivered to the king. The purport of the letter was, ``I \emph{Siva}, who dwell in the city of \emph{Alavei} (Madura), which is surrounded by a wall, and adorned with palaces, say to the king \emph{Chera} that P\={a}napattiren, a minstrel and one as dearly beloved by me as thou art, is coming over to wait on thee. As thou hast acquired renown by lavishing thy wealth upon poets, give him what he wants and dismiss him.'' The king in his blind zeal to \emph{Siva} believing the letter to have been actually indited by the god himself, received it with great solemnity, and placing the bearer on the throne, presented him with considerable wealth. He then returned to Madura, and lived there during the remainder of his life, attending the temple thrice a day, and singing the praises of the god to the lute. It is stated that his wife also was well accomplished in singing, and that she acquired great fame by having carried away the palm in a contest in singing with a songstress whom the king \textbf{R\={a}jar\={a}ja P\={a}ndiyen} had sent for from Ceylon.

\section{Periya \={A}zhv\={a}r. --- பெரியாழ்வார்.}
\paragraph{}
This poet was a native of \emph{Villiputt\={u}r}, and has had the distinction of carrying a prize in a poetical contest in the court of \textbf{Vallabadeva P\={a}ndiyen}, king of Madura. He is considered as one of the \emph{\={A}zhv\={a}rs} or special votaries of \emph{Vishnu}, and that four-hundred stanzas of the \emph{N\={a}l\={a}yiraprabhandam} (நாலாயிரபிரபந்தம்) are said to have been contributed by him.

\section{Perunch\={i}tan\={a}r. --- பெருஞ்சீத்தனார்.}
\paragraph{}
No particulars of the life of this poet are known, except that he was one of the forty-nine professors of the Madura College. His impromptu on the merits of the \emph{Kural}, however, has been preserved, and we give it below.
\begin{quotation}\small
ஏதமில்வள்ளுவரின்குறள்வெண்பாவினா

லோதியவொண்பொருளெல்லாமுரைத்ததனாற்

றாதவிழ்தார்மாறதாமேதமைப்பயந்த

வேதமே மேதக்கன.
\end{quotation}
\begin{quotation}\small
\emph{O king P\={a}ndiya, who wears the garland of full-blown flowers! Say, what are more excellent, the Vedas, or the short sweet distichs, through the medium of which the faultless Valluver has divulged the whole contents of the Vedas.}
\end{quotation}

\section{Perund\={e}van\={a}r. --- பெருந்தேவனார்.}
\paragraph{}
There were two poets of this name, contemporaries, and both professors of the Madura College. One of them distinguished himself by writing an imitation of the \emph{Mah\={a}bh\={a}rat} in \emph{Venba} metre, and which is now known by the appellation of \emph{Sangattu B\={a}radam} (சங்கத்துப்பாரதம்), and the other by writing a treatise on versification, entitled \emph{Kavis\={a}garam} (கவிசாகரம்).

The following is the impromptu which the author of the \emph{Sangattu B\={a}radam} uttered on the merits of the \emph{Kural}:---
\begin{quotation}\small
எப்பொருளும்யாருமியல்பினறிவுறச்‌

செப்பியவள்ளுவர்தாஞ்செப்பவரு --- முப்பாற்குப்‌

பாரதஞ்சீராமகதைமனுப்பண்டைமறை

நேர்‌வனமற்றில்லைநிகர்‌.
\end{quotation}
\begin{quotation}\small
\emph{As the Kural which Valluver has composed contains in itself everything, clothed in language easy to be understood by all, so it may bear comparison with the Mah\={a}bh\={a}rat, R\={a}m\={a}yana, and Menu's Institutes, nay even with the Vedas of old.}
\end{quotation}
\paragraph{}
The author of the \emph{Kavis\={a}garam} carried his admiration of the \emph{Kural} to the highest pitch, as will be seen by the subjoined impromptu which he uttered.
\begin{quotation}\small
பூவிற்குத்தாமரையே பொன்னுக்குச்சாம்புனத

மாவிற்கருமுனியாயானைக்கமரரும்ப

றேவிற்றிருமாலெனச்சிறந்ததென்பவே

பாவிற்குவள்ளுவர்‌ வெண்பா.
\end{quotation}
\begin{quotation}\small
\emph{The best among flowers is the lotus; the best among metals is gold; the best among kine is Kamadhenu (the celestial cow); the best among elephants is Ir\={a}vata (the white elephant of Indra); the best among gods is Vishnu; and the best among books is Valluver's Kural.}
\end{quotation}

\section{Philip De Melho. --- பிலீப்பு தெ மெல்லோ.}
\paragraph{}
The name of Philip [Philippus] De Melho, the first native who was admitted to the office of Minister of the Dutch Reformed Church in Ceylon, and who occupied a prominent place in the literary world, stands foremost among the divines and linguists, whether of European or Ceylonese birth, who flourished in the island during the eighteenth century. As an oriental poet also, he ranks high, and his qualifications as a Biblical translator were unsurpassed. Such was in fact his great learning that Governor \textbf{Falk} [Falck], a wise politician and an elegant scholar, complimented him with the title of ``Rabbi De Melho.''

De Melho was born at Colombo, on the 23d of April 1723. He belonged to a Tamil family distinguished for its respectability and opulence, and was the second son of Mr. \textbf{Simon De Melho}, who occupied the important post of Chief Tamil Mudeliar of the Gate,\footnote{In Tamil \emph{V\={a}sul}, signifying the king's Court.} and as such, was attached to the personal staff of no fewer than ten Dutch Governors, including the tyrant \textbf{Peter Vuyst} and the good Baron \textbf{Von Imhoff}. De Melho studied at an early age at the Colombo Seminary, which the Dutch Government founded in 1705, for the purpose of qualifying native youths for ecclesiastical as well as secular professions. Here he was taught Hebrew, Greek, Latin, Dutch, Portuguese and Tamil, as well as Theology. His accademical career was a brilliant one and before his seventeenth year he had become the senior student. At the age of 20 he finished the prescribed course of study and was the most qualified student to be sent to the University of Leyden. The Government desired him to proceed to Holland, which he refused to do. He was therefore on the suggestion of the Rev. J. P. Witzelius, Rector of the Seminary, appointed as a Native Proponent at Colombo, on the 16th April 1744, after he had passed with great credit a public examination in Hebrew and Greek as also in Theology by the learned members of the Consistory. On the 2d of August 1744, De Melho married Miss Magdalene Jurgen Ondaatje, daughter of Mr. \textbf{Philip Jurgen Ondaatje}, Translator to the ``\emph{Hoff van Justitie}'' (the highest tribunal in the island) and sister of the Rev. \textbf{W. Jurgen Ondaatje}, Clergyman of Colombo and Rector of the Seminary. By this union he had a large number of daughters and only two sons who both died at an early age whilst successfully pursuing their studies, one at Amsterdam, and the other in the island.

Though De Melho, as Native Proponent, was to confine his ministrations to Natives, yet on the recommendation of the Consistory the Government privileged him to preach to the Dutch likewise; and thus at the very outset of his ministerial career, he had an auditory for his discourses from the pulpit that might require the full exercise of his eloquence. In 1745, upon the urgent solicitation of the Dutch residents at Negombo who wished to secure for themselves the benefit of his spiritual instructions, he was sent thither, but his stay amongst them appears not to have been long; for we see him again employed in Colombo in the following year, not only as Proponent, but also as one of the Committee for revising the Tamil version of the New Testament from the original text, and by himself as the reviser of the translations of all religious works intended for circulation amongst the natives. In 1746, when the Government established a Normal school at Colombo, De Melho from his known abilities and talents was selected by the Governor as one of its teachers; and in 1747, when owing to a paucity of ordained ministers, it was found necessary to appoint a fixed Assistant Preacher in the Dutch language, De Melho alone amongst the body of Native Proponents, was found worthy of filling that office.

De Melho having served with great efficiency as a Native Proponent and Assistant Dutch Preacher for a period of five years, applied in 1749 for ordination in Ceylon, with a view to enlarge the sphere of his usefulness. Both Governor \textbf{Gollenesse} and the Consistory marked their sense of the appreciation of his eminent abilities and valuable services by willingly and earnestly recommending him to the Supreme Government of Netherland's India, who thereupon permitted his ordination in the Island itself, after previous examination by the Consistory. On the 21st January 1750, De Melho, having been examined by the members of the Consistory assembled in the Fort Church at Colombo on various important theological points, and found to be in every respect a fit candidate for the ministerial office, was with the sanction of the Governor, ordained according to the rules of the Church of Holland, and appointed Minister to the Western District\footnote{Under the Dutch Government, Ceylon, was for ecclesiastical and educational purposes, divided into three Districts, viz. the Northern, Western and Southern districts, the principal stations of which were respectively Jaffna, Colombo and Galle.} with the usual pay and emoluments assigned to a clergyman on the establishment of the East India Company. This is the solitary instance of an individual, being ordained in the island, under the Dutch Government without previous University education. As a reward for his good services in connexion with the Normal School, he was in March following, promoted to be Rector of that institution.

In the same year in which De Melho was ordained, he completed the revision of the Tamil version of the New Testament up to the Epistle to the Ephesians. This was a work which devolved on him exclusively from the want of zeal and co-operation on the part of his colleagues. In the preparation of the version, he was guided by the original Greek and aided by various standard versions of the New Testament. He also consulted many approved theological and philological works, with all of which he appears to have been familiar at an early age. In consequence, however, of a discussion which took place respecting the version, the Government deferred its publication and referred the matter for the decision of the Supreme authorities at Batavia. They felt it their duty to thank the reviser for the zeal and attention he had hitherto devoted to this highly important and difficult undertaking, and to direct the examination of the version by a Committee of Clergymen and Laymen skilled in the languages, and on the receipt of their report, the Ceylon Government ordered it to be printed and circulated amongst the natives. But although the revision of the Tamil version of the New Testament was a work of itself more than sufficient for one individual, how gifted soever he may be, yet it was not the only work in which he was engaged at the time. He also wrote an elaborate work in Dutch under the title of ``The Triumph of Truth, or a Refutation of the Principal Dogmas of the Church of Rome,'' with a view of arresting, as he thought it would, the rapid progress which Catholicity was making in Ceylon about the time through the persevering labours of the Goanese Fathers.\footnote{See Governor Imhoff's Report, in Lee's Rebeyro, p. 176.} This work being approved by the Consistory was afterwards rendered by himself into Tamil, under the title of \emph{Sattiyattin-jeyam} (சத்தியத்தின்செயம்), and published at the expense of Government in 1753, with a dedication in Latin, Dutch and Tamil to the Governor General \textbf{Vander Paar} and the members of the Council of Netherland's India, couched in the most elegant language. At the conclusion of the preface to this work, De Melho announced that he was preparing for publication in Tamil a translation of the Liturgy of the Dutch Reformed Church and that of a selection of the Psalms of David, the Decalogue, the Lord's Prayer, the Creed and the Songs of the Blessed Virgin, Zacharias and Simeon in metre besides a work to be called a Refutation of Hinduism. The two former were published respectively in 1755 and 1760; the latter appears to have remained in manuscript and is not now forthcoming.

In November 1753, De Melho was appointed Minister to the Northern District in the room of the Rev. \textbf{G. Potken}, who was his examiner in the Sacred languages. He did not, however, leave Colombo till the 25th of February 1754, for the scene of his future labours, which with his family he reached on the 12th March following. The duties which devolved on him in this his new appointment were too onerous for any but a man like him of boundless activity; for he had to superintend not only the numerous churches and schools in Jaffna and Wanny, but those at Trincomalie and Batticaloa, which he periodically visited, and in 1758, he was relieved of a portion of these duties, by the appointment as his colleague of the Rev. \textbf{B. Jansze}, a Tamil native, who was educated in Holland.

In October 1756, the death of his father compelled him to visit Colombo, where he remained for some time. Here he was unexpectedly called upon by Governor \textbf{Schreuder} to undertake the translation of the unfinished portion of the Tamil New Testament, viz. from the Epistle to the Ephesians to the end of the Revelation; and his well known zeal prompted him to execute it with alacrity. In 1759 the whole of the New Testament in the Tamil language was for the first time published in Ceylon; and when copies of it were sent to the Supreme Government the Company's Directors increased his salary in acknowledgment of his praiseworthy labours, styling him: ``The Great Labourer,'' as will be seen by their despatch dated 5th of August 1760, and which we quote below.

``To hear of the regular performance of Divine Service at Colombo by the Rev. Messrs. \textbf{Bronsveld}, \textbf{Zybrands} and \textbf{Meyer} and at Wolvendhal, by the Rev. Messrs. \textbf{Ondaatje} and \textbf{Philipsz}, as also at Jaffna by the Rev. Messrs. \textbf{De Melho} and \textbf{Janz} at Galle by the Rev. Messrs. \textbf{Schoults} and \textbf{Smith} and at Trincomalie by the Rev. Mr. \textbf{De Zelun}, and that the Rev. \textbf{Peter Cornelius} who has recently arrived there, is about to be employed by you to the no small benefit of the Church of God, has not been less agreeable to us than that the translation of the New Testament into the Tamil language has at length by the praiseworthy labours of the Rev. Mr. \textbf{De Melho} fully attained its object, as appears from the copies thereof which reached us; and this being now completed, we hope and wish that the blind heathen may be more and more enlightened and conducted to the true knowledge of our Rational Religion. As the only object we have in view in this respect is the promotion of true sanctifying faith, so is it likewise to manifest to that Great Labourer our particular satisfaction, that we have at the day of our sitting granted his Reverence our unsolicited extraordinary augmentation from 80 to 100 florins a month, as an evident token, how readily on our part we reward faithful services, as we in like manner on the same day and for the same reasons granted 100 florins to the Rector of the Seminary \textbf{Dr. Meyer}, in the expectation that they would henceforward redouble their zeal.''

Whilst De Melho was residing at Jaffna, which was and still is the seat of Tamil literature in Ceylon, he met many renowned scholars, poets, and philosophers whose acquaintance he cultivated with a desire to improve his knowledge of philology as well as to gain opportunities which that acquaintance might afford him for bringing under their consideration, with all the force in his power, the arguments in favour of Christianity, and for exposing the absurdities and superstitions of their own system of faith. Among them, we may mention the name of \textbf{K\={u}langkai Tambir\={a}n}, with whom he had frequent religious controversies.\footnote{Vide, p. 47.} While at Jaffna, De Melho enlarged the \emph{S\={u}d\={a}mani Nigandu}, which is the standard Lexicon of the Tamil language, by adding 20 stanzas to the 2d part, and about 100 to the 12th, besides various words and stanzas to others. These additions not only attest his abilities as a philologist, but also his skill as a poet from their splendid versification. They were all unanimously approved by the literati of his day and incorporated with the original; and a portion of them appears in the edition of the ten parts of the work published at the Manipy Press in 1856. The rest may be found in the manuscript copies current in Jaffna. It is not to be supposed that De Melho confined the display of his poetical talents to the above mentioned work; he also composed an elegant panegyric, called \emph{Marudappa Kuravanji} (மருதப்பக்குறவஞ்சி), on \emph{Mardappa Pillei}, Mudeliar of the Gate, under the administration of Governor \textbf{Schreuder}. There are no means of ascertaining the names of the other works of which De Melho was the author. At his death he left behind him besides a highly valuable and extensive library, a large number of manuscripts, including his learned sermons in Dutch, Portuguese and Tamil; but they have all been unfortunately lost, with the exception of his Correspondence on the version of the Scriptures and a corrected draft translation of the Books of Joshua, Judges and Ruth.

Early in 1778, the Rev. Mr. \textbf{Klein}, who was the best Tamil scholar among the Danish Missionaries then residing at Tranquebar, visited Jaffna. Here he heard De Melho's discourse on Tamil philology, in the presence of the Commandeur, as also on several other occasions, and exclaimed with astonishment ``\emph{quantum est quod nescimus}.'' He moreover candidly said, ``I thought I understood the Tamil language, but I must now confess that I am yet a common learner, and wish my circumstances permitted me to remain sometime longer to learn from brother De Melho something more of that language.'' He also requested De Melho to give his brother Missionaries and himself some instructions in the Tamil language by means of correspondence, to correct their writings and to point out the errors committed by them in the translation of the Scriptures, with which request De Melho cheerfully complied.

We will now close the biographical sketch of this laborious and learned man of whom his native country may justly be proud, with a brief notice of his valuable version of the Pentateuch, which was the last work in which he was engaged.

The Tranquebar version of the Old Testament completed and published in 1729, was for good reasons, considered to be imperfect and unadapted to public use; it was therefore in contemplation to prepare a new version in Ceylon for the service of the Tamil Protestant Congregations. With this view the Rev. Mr. \textbf{Ondaatje} (whose name has already been mentioned) executed in 1774, a translation of the Book of Genesis and presented it to Governor \textbf{Falk}, who directed its revision according to the prescribed rule, previous to publication. But as it was deemed preferable to revise the Tranquebar version itself, directions to that effect were conveyed to him. But these were countermanded so soon as it was known in Ceylon that the Tranquebar Missionaries had already begun that work and that the Pentateuch and the Books of Joshua and Judges had been published. On the receipt of this order, Mr. \textbf{Ondaatje} represented to the Governor that he had read the revised version with great care and attention and had found it to be faulty, stating at the same time his opinion that it would be preferable to translate the Old Testament anew. He moreover added ``I have received some portions of the Old Testament rendered by De Melho, which appear superior to the first as well as the second Tranquebar version.'' To this communication, His Excellency replied that both he and De Melho had permission to commence a new version. De Melho accordingly in 1779 and 1780, submitted to the Governor the versions of the Books of Genesis and Exodus made from the original Hebrew with the aid of the Septuagint, the Latin versions of \textbf{Tremellius} and \textbf{Junius}, and the Netherland's State Bible. It was however signified to De Melho that the Government would prefer a revision of the recently corrected Tranquebar version; and they accordingly wished him to prepare such a revision, exhibiting in a separate paper the errors found therein with the required corrections. Pursuant to this order, De Melho prepared and transmitted to Governor \textbf{Falk} a revised version of the Pentateuch, accompanied by a report on the subject and an annotation of the errors in the Tranquebar version. The following passage from that report demonstratively shews the character of that version:---

``The language and style have not been arranged in accordance with the solemnity of such a divine revelation as the Holy Scriptures. The spelling of a great many words is very defective. Foreign words have been introduced, which in Tamil style, appear ungraceful and deform the language and disfigure it. Abundant errors have also crept into it, consisting of unnecessary additions of words which are not in the sacred text and inadmissible omissions of those which are in it, bad and incorrect renderings and incompatible interpretations, instead of translations, and that in so many places as the undersigned has noted down the same, as is to be seen in the accompanying statement of errors in the five Books of Moses alone.''

De Melho's version and his criticisms on the Tranquebar version were sent to the Missionaries there. To the latter they had nothing to say in vindication. With regard to the version, all with whom they consulted pronounced the language employed in it to be excellent and choice; but raised a doubt whether it could be generally understood by the common people. To settle this important point the version was publicly and solemnly read in the Jaffna Fort Church to a large body of learned Tamils and to other auditors and the question put to them whether it was intelligible to the common people. They unanimously replied that it was so, and that the language used therein was matchless, elegant, pathetic and heart cheering, worthy of and becoming Holy Writ, while that of the Tranquebar version was a mixture of all words current on the Coast and was extremely uncouth, barbarous and ridiculous, owing to the grammatical errors and the vulgarisms with which it abounded.\footnote{That unparalleled Tamil scholar of modern times, the learned Beschi, writes in the following strain of the Tranquebar version:--- ``Can those books be fairly called the Word of God, which the Tranquebarians, who do not at all write correctly in Tamil the name of their country, have handed down to us pretending that they have translated the Holy Scriptures in Tamil, whilst ignorant of that language, they have to the bitter paining of our ears, written them in barbarous words. By this means, the truth of God's word has been darkened, and by depriving it of its excellence, been tarnished, even as if a costly bright gem were buried in mire, or poison mixed with ambrosial sweet or a beautiful picture stained with ink.'' \emph{Veder Vilakkam}, Chap. 16.} Reference was also made on this point to De Melho himself, who stated that the words remarked on as high by the Tranquebar Missionaries (who certainly were not competent judges) are no other in reality than pure Tamil words, unintelligible to none but to such as understand no other than lame and bastard Tamil. To obviate however all difficulty on this score and to render his labours generally useful, he prepared an alphabetical glossary of the words in question to be appended to the version, which was accordingly published by the Government in 1790. Though De Melho had now reached his 67th year, yet the energies of his powerful and vigorous mind were unceasingly directed to the translation of the other portions of the Old Testament, and if his life had been spared for a few years more, he would have no doubt completed this great work; but it was ordained otherwise. He died on the 10th of August 1790.

\section{P\={e}i \={A}zhv\={a}r. --- மேயாழ்வார்.}
\paragraph{}
P\={e}i \={A}zhv\={a}r, who was one of the twelve \emph{\={A}zhv\={a}r}, or special votaries of \emph{Vishnu}, enjoyed likewise the reputation of a poet. He was born at \emph{Mailap\={u}r}, in the Carnatic, and having finished his studies, spent all his life in making pilgrimages to the different \emph{Vaishnava} shrines, and in diffusing the \emph{Vaishnava} doctrines in the country. The \emph{N\={a}l\={a}yiraprabhandam} (நாலாயிரபிரபந்தம்) contains a hundred stanzas which he composed at Tirukoval\={u}r, as he was visiting that place in company with \textbf{Poigai \={A}zhv\={a}r} and \textbf{P\={u}dat \={A}zhv\={a}r}.

\section{Poigai \={A}zhv\={a}r. --- பொய்கையாழ்வார்.}
\paragraph{}
The history of Poigai \={A}zhv\={a}r like that of his compeers has been disguised in myths, and it is difficult to extricate it from them. One thing, however, is certain, that he was a native of \emph{K\={a}njipuram}, and ranked high as a poet. He employed his muse solely in hymning the praises of \emph{Vishnu}, and contributed one hundred stanzas to the \emph{N\={a}l\={a}yiraprabhandam} (நாலாயிரபிரபந்தம்).

\section{Poiy\={a}mozhi Pulaver. --- பொய்யாமொழிப்புலவர்.}
\paragraph{}
Poiy\={a}mozhi Pulaver was a celebrated poet who flourished in the reign of the king \textbf{Vanang\={a}mudi P\={a}ndiyen}. Of his compositions, we have met with only an erotic poem in 425 stanzas, under the title of \emph{Tanjeiv\={a}nan-k\={o}vei} (தஞ்சைவாணன்கோவை), which is much valued not only on account of its beautiful versification, but also because it is written in illustration of the rules of \textbf{Natkavir\={a}ja Nambis'} \emph{Agapporul}.

The following anecdote of Poiy\={a}mozhi Pulaver, which we have extracted from the \emph{Dinavartam\={a}ni},\footnote{See No. 46 of 1856.} though tinged with the marvellous may be found interesting.

Poiy\={a}mozhi Pulaver being desirous of having the Madura College re-established, went to the king to speak to him on the subject, but finding the king in the temple lying prostrate at the feet of the image of \emph{Siva} and worshipping it, he addressed to him the following impromptu:---
\begin{quotation}\small
குழற்காலரவிந்தங்‌ கூம்பக்குமுதமுகையவிழ

ஙீழற்கான்மதியமன்றோ நின்றிருக்குலநீயவன்றன்‌

அழற்காலவிர்சடை மீதேயிருந்துமவ்வந்திவண்ணன்‌

கழற்கால்வணங்குதியோ வணங்காமுடிக்கைதவனே.
\end{quotation}
\begin{quotation}\small
\emph{O P\={a}ndiya, who never bows thy head to anyone! Dost thou lie prostrate at the feet of Andivannan (Siva), while the progenitor of thy holy race the cool-beamed moon, to whose rays the lotus-flower closes, and the nymphæ-flower expands, is sitting on his radiant lock of hair.}
\end{quotation}
\paragraph{}
The king upon this rose up and questioned him on the object of his visit, to which he replied that he had come to beg permission for the re-establishment of the Madura College. The king wishing to put to the test his ability for such an undertaking, desired him to improvise a verse which would cause the statues of the professors of the College, which were in the temple to nod. He obeyed and improvised the following verse:---
\begin{quotation}\small
உங்களிலேநானொருவனொவ்வேனோவொப்பேனோ

திங்கட்குலனறியச்‌ செப்புங்கள்‌ --- சங்கத்து

பாடுகின்றமுத்தமிழ்க்கென்‌ பைந்தமிழுமொக்குமோ

ஏடவிழ்தாரேழெழுவீரே.
\end{quotation}
\begin{quotation}\small
\emph{Declare, that the king of the lunar race may know, whether I will be like one of you or not, and whether my lays will equal the three kinds of lays which ye sang in the College or not, O ye the forty-nine (professors) adorned with garlands of full-blown flowers!}
\end{quotation}
\paragraph{}
The statues thereupon nodded by way of affirmation. The king afterwards conducting him to the bank of the \emph{Pott\={a}marei} tank in which the board on which the professors sat was lying plunged, bade him make it float by improvising a verse; and he then improvised the following:---
\begin{quotation}\small
பூவேந்தர்முன்போற்‌ புரப்பாரிலையென்றும்‌

பாவேந்தருண்டென்றும்‌ பான்மையால்‌ --- மாவேந்தன்‌

மாறனறிய மதுராபுரித்தமிழோர்‌

வீறனையே சற்றேமித.
\end{quotation}
\begin{quotation}\small
\emph{Floatest thou awhile, O the board of the Tamil poets of Madura! that the great P\={a}ndiyen may know, that though there are no kings to patronize poetry as in times past, yet there are poets even at the present time.}
\end{quotation}
\paragraph{}
The board came up to the surface of the water immediately as he had uttered the verse, and then sank down again. The king after having such ample proofs of his abilities, nevertheless took very little notice of him; and stung with the neglect he set out to go back to his country. The Queen having heard of it trembled for fear that his malison would ruin the king. She therefore, disguised as a palanquin-bearer, bore his palanquin for some distance, and when he came to know of it, he at once stopped her and wished to know the cause why she had thus condescended to him. She told him that her object was to pacify his wrath, and obtain from him a blessing for the king. He thereupon ejaculated the following verse in her praise:---
\begin{quotation}\small
உமையாளுநீயுமொருங்கொப்பேயொப்‌பே

யுமையாளுக்கங்குண்டோரூனம்‌ --- உமையாடன்‌

பாகந்தோய்ந்தாண்டான்‌ பலிக்குழந்தான்பாண்டியனின்‌

னாகந்தோய்ந்தாண்டானரசு.
\end{quotation}
\begin{quotation}\small
\emph{Uma (Parvati) and thou stand in equal rank; but Uma has a stain; for he who shared with her part of his body, betook himself to begging, while P\={a}ndiyen who pressed thee to his bosom, reigns.}
\end{quotation}

\section{Ponmudiy\={a}r. --- பொன்முடியார்.}
\paragraph{}
Ponmudiy\={a}r figured in the Madura College as one of the forty-nine professors; but nothing further is known concerning him. Of his poetical compositions the only relic is the subjoined impromptu on the \emph{Kural}:---
\begin{quotation}\small
கானின்றதொங்கலாய்‌ காசிபனார்தந்ததுமுன்‌

கூநின்றளந்த குறளென்ப --- நூன்‌முறையான்‌

வானின்றுமண்ணின்றளந்ததேவள்ளுவனார்‌

தாநின்றளந்தகுறள்‌.
\end{quotation}
\begin{quotation}\small
\emph{O thou adorned with the garland of fragrant flowers! It is said that the Kural (meaning Vishnu in his incarnation as a dwarf) produced by Kasypa in times of yore measured the earth\footnote{For the particulars of this myth see Moor's Hindu Pantheon, p. 186.}; but the Kural now produced by Valluver has measured both the earth and the heaven (by its treating of everything which they contain).}
\end{quotation}

\section{P\={o}gar. --- போகர்.}
\paragraph{}
P\={o}gar was one of the sages of antiquity whose biography has merged into fable. He however appears to have been well versed in the art of poetry, as well as that of medicine. Of his works all written in verse, Dr. Ainslie has in his Materia Indica noticed four; viz. 1. \emph{P\={o}gar Y\={o}gam\={a}rgam} (போகர் யோகமார்க்கம்), 2. \emph{P\={o}gar Ezhun\={u}ru} (போகர் எழநூறு), 3. \emph{P\={o}gar Tirumandiram} (போகர் திருமந்திரம்), and 4. \emph{P\={o}gar Nigandu} (போகர் நிகண்டு), all which treat of the preparations of the several kinds of medicines. Some think that he was the same with \textbf{Fo} or \textbf{Fohi}, who introduced Buddhism into China in the latter part of the first century of the Christian era, but this is rather doubtful.

\section{P\={o}kiy\={a}r. --- போக்கியார்.}
\paragraph{}
P\={o}kiy\={a}r was one of the forty-nine professors of the Madura College; but nothing further is known of him, except that he was the author of a stanza merely enumerating the divisions of the \emph{Kural}.

\section{Pugazhendi. --- புகழந்தி.}
\paragraph{}
Pugazhendi was a contemporary poet with \textbf{Kamber} and \textbf{Ottak\={u}ter}, but he was retained at the Court of the king \textbf{Varaguna P\={a}ndiyen} at Madura. He possessed a masterly talent for composing \emph{Venba}, a kind of versification which is so difficult that it has been styled ``the tiger of poets.'' He wrote a poem, entitled \emph{Nalavenba} (நளவெண்பா), in which he has in 417 stanzas of the \emph{Venba} metre narrated the adventures of \textbf{Nala} and \textbf{Damayanti}, and this work has obtained for him a lasting fame, as it is confessedly the most beautiful composition of the kind in the Tamil language. He also wrote a treatise in 71 stanzas, under the title of \emph{Retna-churukam} (இரத்திநச்சுருக்கம்), containing a specification of the metaphors to be used in poetry when describing the perfections of the fair sex. Another poem styled \emph{Alliaras\={a}ni M\={a}lei} (அல்லிஅரசாணிமாலை), treating of the history of the Queen \textbf{Alliaras\={a}ni}, is likewise attributed to him; but it bears internal evidence of its spuriousness and is altogether unworthy of his transcendent genius. Many anecdotes are related of Pugazhendi in proof of his extraordinary wit and skill as an improvisator, and we subjoin one of them. Once when he was visiting the court of \textbf{Kulotunga Chola} at Ureiur, he happened to go out for a walk with the king and his poet \textbf{Ottak\={u}ter}. \textbf{Auveiyar}, who was then sitting in the street with her feet stretched out straight before, withdrew one of them as the king passed by her, and the other as Pugazhendi passed by her; but when she saw \textbf{Ottak\={u}ter}, who followed, she stretched them both again. \textbf{Ottak\={u}ter} offended at the slight offered to him, asked for an explanation, upon which she said, ``I withdrew one foot for the king, because he wears the crown, and both feet for Pugazhendi, because he is a great poet; but as you are nothing but a dunce, I have not withdrawn either of my feet for you, and if you are as able as Pugazhendi, prove it by improvising a verse in which while praising the \emph{Chola} country and its king, the word ``\emph{mathi}'' (மதி)\footnote{The word (மதி) means ``\emph{the moon},'' as well as ``\emph{wit}.''} shall occur three times. He thereupon uttered the following stanza:---
\begin{quotation}\small
வெள்ளத்தடங்காச்‌ சினவாளை வேலிக்கமுகின்மிடறொடித்துத்

தள்ளிமுகிலைக்கிழித்துமழைத்துளியோடிறங்குஞ்சோணாடா

கள்ளர்மறவர்குறும்படக்குங் கண்டாகண்டர்பெருமானே

பிள்ளைமதியாலென்மாது பேதைமதியுமிழந்தனளே.
\end{quotation}
\paragraph{}
As the word ``\emph{mathi}'' (wit) occurred only but twice in the stanza, \textbf{Auveiyar} asked \textbf{Ottak\={u}ter}, pray where is thy other wit? but he was so ashamed that he could not reply. She then addressed herself to Pugazhendi and asked him to try his skill in improvising a verse in praise of the \emph{P\={a}ndiya} country and its king, with the word ``\emph{mathi}'' introduced into the body of the verse three times, and he instantly improvised the subjoined.
\begin{quotation}\small
பங்கப்பழனத்துழுமுழவர்‌ பலவின்கனியைப்பறித்ததென்று

சங்கிட்டெறியக் குரங்கிளநீர்‌ தனைவிட்டெறியுந்தமிழ்நாடா

கொங்கற்கமராபதியளித்த கோவே யமரர்குலதீபா

வெங்கட்பிறைக்குங்‌ கருப்பிறைக்கு மெலிந்தபிறைக்கும்‌ விழிமகவே.
\end{quotation}
\paragraph{}
Applauding him for his successful performance, she again addressed herself to him to improvise another verse, the first line of which should begin with the word ``\emph{kari}'' (கரி)\footnote{The word (கரி) means both ``\emph{Charcoal}'' and ``\emph{an Elephant.}''} and the last end with the word ``\emph{Umi}'' (உமி)\footnote{(உமி) as a noun signifies ``husk;'' but as a verb ``\emph{to spit.}''} and he thereupon uttered the following:---
\begin{quotation}\small
கரியையுரித்தணிந்தாய்‌ காமன்றனையெரித்தாய்

நரியைப்பரியாக்குநாதனே --- விரிசடையில்‌

துண்டுமதியைச்‌ சூடுஞ்சுத்தனேயாழிகட

லுண்டநஞ்சைச்சற்றேயுமி.
\end{quotation}

\section{P\={u}dattazhv\={a}r. --- பூதத்தாழ்வார்.}
\paragraph{}
This poet was a native of \emph{M\={a}valipuram}, near Sadras; but being deified by the \emph{Vaishnavas} to whose sect he belonged, the real particulars of his life have been completely enveloped in fable. Of his compositions, however, there is a series of one hundred stanzas which now form a part of the \emph{N\={a}l\={a}yiraprabhandam} (நாலாயிரபிரபந்தம்).

\section{P\={u}denchendan\={a}r. --- பூதஞ்சேந்தனார்.}
\paragraph{}
Nothing further is known of the personal history of P\={u}denchendan\={a}r than that he was the son of \textbf{Sengund\={u}r kizh\={a}r}, one of the forty-nine professors of the Madura College. He wrote a didactic poem in forty stanzas, under the title of \emph{Iny\={a}n\={a}lpadu} (இனியாநாற்பது), containing an enumeration of everything which is pleasant. It is ranked amongst the works which had received the approbation of the College, and we subjoin a stanza from it as specimen.
\begin{quotation}\small
தங்கண்மரபுடையார்‌ தாம்வாழ்தன்‌முன்னினிதே

யங்கண்விசும்பினகனிலாக்காண்பினிதே

பங்கமில்செய்கையராகிப்பரிந்தியார்க்கு

மன்புடையராதலினிது.
\end{quotation}
\begin{quotation}\small
\emph{Pleasant it is for men to live happily together with their kindred, and pleasant it is to behold the full-moon among the beautiful clouds; but more pleasant is the love kindly bestowed on all by those whose actions are free from guilt.}
\end{quotation}

\section{R\={a}machandra Kavir\={a}yer. --- இராமசந்திரகவிராயர்.}
\paragraph{}
The merits of R\={a}machandra Kavir\={a}yer as an able poet and elegant writer of Dramas in modern times cannot be denied. He was a native of \emph{Rajanell\={u}r}, in the Carnatic, but resided chiefly at Madras, where his high attainments had procured for him the friendship and countenance of Mr. \textbf{Ellis}, the accomplished Orientalist. He wrote the following five dramas: 1. \emph{Sakuntala V\={i}l\={a}sam} (சகுந்தலைவிலாசம்), 2. \emph{P\={a}rada V\={i}l\={a}sam} (பாரதவிலாசம்), 3. \emph{T\={a}ruk\={a} V\={i}l\={a}sam} (தாருகாவிலாசம்), 4. \emph{Iraniya V\={a}sagapp\={a}} (இரணியவாசகப்பர்), and 5. \emph{Rangoon-chandei N\={a}dagam} (இரங்கூன்சண்டைநாடகம்) besides an ode in honour of Mr. \textbf{Ellis}, eulogizing his wise administration of the Collectorate of Madras. The following stanza is reported to have been addressed to the poet by Mr. \textbf{Ellis}:
\begin{quotation}\small
செந்தமிழ்ச்செல்வனு மோராயிரந்தலைச்‌ சேடனும்யாழ்‌

சுந்தரத்தோடிசைவல்லோனும்யாவருந் தோத்திரஞ்செய்‌

கந்தனைச்சொல்லுங்‌ கவிராமசந்திரனைக்‌ கண்டுவெட்கி

யந்தரம்வெற்புழி பாதாளலோகமடைந்தனரே.
\end{quotation}
\begin{quotation}\small
\emph{He who skilfully plays on the lute (meaning N\={a}rada), and he who is enriched with the knowledge of the elegant Tamil (meaning Agastiyer), and the thousand headed Sesha, seeing the poet R\={a}machandra, who chants the praise of Skanda, the object of the praise of all, were ashamed (by a consciousness of inferiority), and quitting the earth, one betook himself to the sky, another to the mountain (Pothiya), and the other to the nether world.}
\end{quotation}

\section{R\={a}m\={a}nuja Kavir\={a}yer. --- இராமானுசக்கவிராயர்.}
\paragraph{}
Of the poets who flourished at Madras in the first half of the present century R\={a}m\={a}nuja Kavir\={a}yer claims a prominent place, on account of his superior talents. He was Moonshee to the Rev. \textbf{W. H. Drew} of the London Missionary Society, and proved himself a valuable helper to that gentleman in his Tamil studies and pursuits. When Mr. \textbf{Drew} published the first volume of his English version of the \emph{Kural} in 1840, R\={a}m\={a}nuja Kavir\={a}yer was alive, but died before the second volume came out in 1852, and Mr. \textbf{Drew} in the notice affixed to the second volume mourns his death ``as a loss not only to himself, but also to the interests of the Tamil literature generally,'' and adds that ``his thorough knowledge of the language; the native vigour of his mind; his power of long continued, patient labour; his delight in work; together with his manliness and integrity of character so remarkable in a Hindu, made him to me a helper in my Tamil studies and pursuits, whose value I could not over-estimate. But he is gone; and I shall enjoy his aid no more.'' He was the author of an ``Amplification of \textbf{Parimelazhagar's} comments on the \emph{Kural},'' which Mr. \textbf{Drew} has published with his English version of that work.

\section{R\={a}sappa Kavir\={a}yer. --- இராசப்பக் கவிராயர்.}
\paragraph{}
This poet was a native of \emph{Kutt\={a}lam} in the Carnatic, and is well known as the author of a popular Drama, under the title of \emph{Kutt\={a}la-kuravanji} (குற்றாலக்குறவஞ்சி), in honour of \emph{Siva} in his character of \emph{Kutt\={a}la N\={a}ther}, as he is worshipped at Kutt\={a}lam. Neither the date of his existence, nor his personal history is forthcoming; but that he was an able dramatist may be inferred from the style of his Drama, which now generally serves as a model.

\section{R\={e}van\={a}tiy\={a}r. --- இரேவனாத்தியர்.}
\paragraph{}
R\={e}van\={a}tiy\={a}r was connected with the learned circle of \emph{Saiva} ascetics at Chillambaram, where he lived and died. He claims notice as the author of a metrical lexicon of Tamil synonyms, which though much briefer than \textbf{Mandalapuruder's} is yet better adapted for memory.

\section{Sampanter. --- சம்பந்தர்.}
\paragraph{}
Sampanter is celebrated as one of the three champions of the \emph{Saiva} faith against the \emph{Jainas} in the reign of the king \textbf{K\={u}na P\={a}ndiyen}. The particulars of his life are related at great length in the \emph{Tirutonder Pur\={a}nam}; but with the exception of what we have here extracted, the rest is too much blended with the marvellous to deserve any notice. It appears that he was born of \emph{Br\={a}hman} parents of the tribe of \textbf{Kauniyer}, at \emph{Shiy\={a}li}, in the Carnatic. Such was the precocity of his genius, that while yet a child he acquired marked distinction as a poet, and visiting the different \emph{Saiva} shrines chanted a \emph{Pathigam} (a poem of ten stanzas) in praise of each. About this time \textbf{K\={u}na P\={a}ndiyen} embraced the \emph{Jaina} religion, and exerted his utmost to establish it throughout his kingdom, but Sampanter having succeeded in curing the king of a fever, which according to the legend had baffled the drugs and spells of the \emph{Jaina} priests, he not only reconverted him to the \emph{Saiva} religion, but at the same time instigated him to impale alive on the banks of the \emph{Vygai} river at Madura no fewer than eight thousand \emph{Jainas} who had disputed with him and been vanquished. Though he had always a strong predilection for the ascetic life, yet in his sixteenth year upon the urgent entreaty of his father, he married the daughter of \textbf{Namb\={i}yand\={a}r Nambi}, a \emph{Brahman} of Tirunell\={u}r; but both he and his wife died on their wedding day, whilst worshipping in the temple at Tirunell\={u}r. The number of \emph{Pathigams} he composed was 16,000, of which only 384 are now extant. Amongst the lost \emph{Pathigams} are reckoned one in honour of the \emph{Saiva} shrine of Tirukkon\={a}malei (\emph{Trincomalie}), and another in honour of the \emph{Saiva} shrine at Tirukethichuram (\emph{Mantotte}).

\section{Sanbagavadivi. --- சண்பகவடிவி.}
\paragraph{}
Sanbagavadivi, styled by way of eminence \emph{Tamil Arival} or one versed in the Tamil language and literature, was a poetess of superior wit and accomplishments. Her history divested of its fabulous additions is as follows: she was the daughter of one of the maidservants (ஆலாத்திப்பெண்கள்) at the court of the king \textbf{Karik\={a}l Chola}, who reigned at Ureiy\={u}r. While yet a child she having recommended herself to the notice of the king by her uncommon genius for belles-lettres, was adopted into the royal family and had a splendid mansion with a retinue of attendants assigned to her. When she reached the age of puberty, her fascinating beauty gained her many admirers; but she resolved not to marry anyone that could not overcome her in poetical contest, which every one who attempted failed, and she retained her solitary dignity until her twenty-fifth year, exulting in her invincibility. At length having excited the indignation of the professors of the Madura College by characterizing them as a set of dunces, in one of her verses, \textbf{Narkira}, the president of that College, encountered her, in the disguise of a wood-seller, and overcame her after a severe and prolonged contest, in which the parties alternately proposed and solved in their verses a series of enigmas as dark as those of a \textbf{Sphinx}. All the verses in which the enigmas were proposed and solved have been preserved, and we subjoin one in which the names of the twelve signs of the Zodiac are concealed.
\begin{quotation}\small
நீரிலானவைதானதுமூன்றுள

நெடியகானில்‌ விளற்கொருமூன்றுள

ஊரிலாடவும்‌ பாடவுங்கொட்டவு

மொன்றிலொன்றொரு பேதையராமெனச்‌

சேரவாங்கு தனுசுடன்கூடவே

தெற்குமாமுனி தன்னையுமாகவே

பாரின்மீது வியலவல்லார்க்கிது

பன்னிரண்டும்‌ வருஷம்பகருமே.
\end{quotation}
\begin{quotation}\small
\emph{Three are found in the water (Pisces, Capricorn and Cancer), three in the forest (Aries, Taurus and Leo) and three in the country (Libra, Gemini and Scorpio), and these including the maiden (Virgo), the bow (Sagittarius), and the southern sage\footnote{Agastiyer, one of whose names is \emph{Kumbamuni}, referring to the fable of his having been born of a water-pot.} (Aquarius), in all twelve, make up the year, if any one is able to understand them.}
\end{quotation}

\section{Sankaranamasiv\={a}yer. --- சங்கரநமச்சிவாயர்.}
\paragraph{}
Neither the date of the birth, nor the death of this poet is known; but we are certain he was alive in \textsc{ce} 1770. He was a native of \emph{Tinnevelly}, and prosecuted his studies under the tuition of \textbf{S\={a}minada Tambir\={a}n}, the author of the \emph{Illakkanakottu}. Upon the request of \textbf{Maruthappa Deven}, the \emph{Poligar} of Utamalei, by whom he was patronized, he wrote a commentary on \textbf{Pavananti's} \emph{Nannul}, which is admitted on all hands to be the best work of the kind.

\section{Saravanaperum\={a}l Aiyer. --- சரவணப்பெருமாள் ஐயர்.}
\paragraph{}
Saravanaperum\={a}l Aiyer was born of a \emph{Jangama} family long settled at \emph{Tirutanigei}, near Madras. He prosecuted his studies under the eye of his father \textbf{Kandappa Aiyer}, who was a man of extensive learning, and being endowed with great talents, he was soon able to assert his superiority over the poets of his day, and to preside over the \emph{Viveha-kalvichch\={a}lei}, a Literary Institution established at Madras by certain Native gentlemen for the promotion of Tamil literature. The first and most prominent of his works was a commentary on the \emph{Kural}, which he published with the text in 1830; and the next, a treatise on grammar, under the title of \emph{Iyattamil-churukkam} (இயற்றமிழ்ச்சுருக்கம்); and another on Rhetoric, under the title of \emph{Aniyiyel-vilakkam} (அணியியல்விளக்கம்). He also wrote a treatise on Geography, entitled \emph{G\={o}lad\={i}pikei} (கோளதீபிகை), in which he exploded the Puranic, and adopted the Copernican system, establishing it by quotations from the ancient Hindu authors, especially with regard to the rotundity of the earth. The latter part of his life was chiefly employed in revising and editing many of the Tamil classical works, which having hitherto been circulated in manuscript, had been rendered erroneous by the carelessness of copyists. Amongst these may be mentioned the \emph{Naidatham} (நைடதம்), \emph{N\={a}ladiy\={a}r} (நாலடியார்), \emph{Tiruv\={a}sagam} (திருவாசகம்), \emph{Tiruvilleiy\={a}del Pur\={a}nam} (திருவிளையாடற்புராணம்), and the moral apophthegms of the poetess \textbf{Auveiyar}.

\section{Saravanamuttu. --- சரவணமுத்து.}
\paragraph{}
Saravanamuttu was the son of \textbf{Manappuli Mudeliar} of \emph{Nellore}, in Jaffna. Being gifted with natural talents and having had very early in life the advantage of pursuing his studies under the guidance of the celebrated Pundit \textbf{Senatharaya Mudeliar}, he was enabled to master the Tamil language and its literature, and to acquire the distinction of being considered as one of the best poets of the day. But unhappily in 1845, when his reputation was at its zenith and much was expected of him, he was hurried away to the grave in the forty-third year of his age. He was a valuable contributor to the Morning Star on several important questions on Grammar, and also took a part in the controversy which was carried through that periodical between \textbf{Vedaghiri Mudeliar} and others, concerning the construction of a verse in the \emph{Naidatham}. He wrote a treatise on Hindu Theology, called \emph{V\={e}d\={a}ntaswayamj\={o}ti} (வேதாந்தசுயஞ்சோதி), and translated into Tamil verse the Sanskrit work on the same subject, entitled \emph{\={A}tmabhoda Prak\={a}sikai} (ஆத்மபோதப்பிரகாசிகை).

\section{Saravanaperum\={a}l Kavir\={a}yer. --- சரவணப்பெருமாள் கவிராயர்.}
\paragraph{}
Saravanaperum\={a}l Kavir\={a}yer was born of a \emph{Vel\={a}la} family at \emph{Nall\={u}r}, in the Carnatic. Having commenced his studies at the proper ago, and mastered the art of poetry, he became a poet at the court of \textbf{Muttur\={a}malinga Setupati}, who ascended the musnud at Ramnad in \textsc{ce} 1795. He was distinguished for his marvellous proficiency in the art of performing at one, and at the same time, eight different things; such as dictating verses, playing at chess, counting separately the grains of paddy and pebbles, as they were thrown on his back, solving Arithmetical problems, etc. His exploits being witnessed by \textbf{Umrutul Umra}, Nabob of Madras, \textbf{R\={a}mavarma R\={a}ja}, king of Travancore, \textbf{Amarasimha Mahar\={a}ja}, king of Tanjore, and \textbf{Raj\={a}dhi R\={a}jasinha}, king of Kandy, procured for him great honors and rich presents. He was the author of a poem in honour of \textbf{Vijaya Regun\={a}tha Tondiman} of Pudukottei, and another in honour of \textbf{Muttirulappa Pillei}, Minister of the \emph{Setupati}. He also wrote many poetical missives which are very much admired for their elegant diction.

\section{Shanmugad\={a}san. --- சண்முகதாசன்.}
\paragraph{}
This poet was well skilled in the composition of songs for the \emph{Nautches}; but his songs, of which there is a collection still extant, though deserving of commendation for their sweetness, are very exceptionable on account of their immoral tendency.

\section{S\={a}min\={a}da D\={e}siker. --- சாமிநாததேசிகர்.}
\paragraph{}
S\={a}min\={a}da D\={e}siker was a celebrated poet who lived in \emph{Tinnevelly} in the early part of the eighteenth century. His parents belonged to the class of \emph{Saiva Vel\={a}las}, and were possessed of some competence; but he entered the \emph{Matum} at Tiruv\={a}vaduturei while yet a youth, and being initiated into the mysteries of the \emph{Saiva} religion, passed through the course of studies required for his profession, and assumed the ascetic life. A learned Tamil poet of the name of \textbf{Myilerum Perum\={a}l}, who became acquainted, with him about this time, perceiving that he was endowed with superior talents and feeling a great interest in him, not only assisted him in acquiring a knowledge of the Tamil grammar and poetry, but sending for a \emph{Br\={a}hman} Pandit named \textbf{Kanakasab\={a}pati Kurukkal} from Sepparai, secured for him his assistance in mastering the Sanskrit language. He continued his studies for twelve years, at the end of which he displayed his abilities, by writing a treatise on the Tamil grammar, under the title of \emph{Illakkanakottu} (இலக்கணக்கொத்து), consisting of 130 aphorisms in metre, with an explanation in prose. Upon the request of \textbf{Myilerum Perum\={a}l}, \textbf{Ambalav\={a}na Tambir\={a}n} the superior of the \emph{Matam} at Tiruvavaduturei, soon afterwards invested him with the office of \emph{Tambir\={a}n}, under the title of \emph{Is\={a}na D\={e}siker}, and appointed him to preside over the \emph{\={I}sana Matam}, belonging to his own community in Tinnevelly.

\section{S\={a}ntalinga D\={e}siker --- சாந்தலிங்கதேிகர்.}
\paragraph{}
S\={a}ntalinga D\={e}siker was an ascetic of the \emph{Vira Saiva} sect, who lived at \emph{Tureiur}, in the Carnatic, and earned for himself much renown by his high attainments both in the \emph{Vedic} and profane literature. He seems to have devoted his attention chiefly to the exposition of the Mystical Philosophy of the Hindus, and he has left no less than six treatises on the subject respectively entitled \emph{Nenjuvidut\={u}thu} (நெஞ்சுவடுதூது), \emph{Ozhivil-odukkam} (ஒழிவிலொடுக்கம்), \emph{Vyr\={a}kya-sadagam} (வைராக்கியசதகம்), \emph{Vyr\={a}kya-d\={i}pam} (வைராக்கியதீபம்), \emph{Kolei-maruttel} (கொலைமறுத்தல்), and \emph{Avir\={o}davundiy\={a}r} (அவிரோதவுந்தியார்), all in verses of different metres.

\section{S\={a}ntalinga Kavir\={a}yer. --- சாந்தலிங்கக்கவிராயர்.}
\paragraph{}
This poet was a native of \emph{Tandaleicheri}, in the Carnatic, and is chiefly known as the author of an ethical poem in a series of 100 stanzas, under the title of \emph{Tandaleiy\={a}r-sadagam} (தண்டலையார்சதகம்), of which we here give the following specimen:---
\begin{quotation}\small
திருவிருக்குந்தண்டலையார்‌ வளநாட்டிலில்வாழ்க்கைச்செலுத்துநல்‌லோர் --- ஒருவிருந்தாகிலுமின்றியுண்டபகற்பகலாமோவுறவாய்வந்த --- பெருவிருந்துக்குபசாரஞ்செய்தனுப்‌பியின்னுமெங்கே பெரியோரென்று --- வருவிருந்தோடுண்பதல்லால்‌ விருந்திலாதுணுஞ்சோறு மருந்துதானே.
\end{quotation}
\begin{quotation}\small
\emph{The worthy householders, who dwell in the country blessed by the sacred Tandaleiyar (a title of Siva), consider the day in which they have no guest as no day: after having treated honourably and dismissed the respectable guest, who has visited them as a friend, they inquire anxiously for virtuous men to partake of their hospitality, for the food taken when there is no guest present, is poison.} \textbf{Ellis.}
\end{quotation}

\section{S\={a}ttan\={a}r. --- சாத்தனார்.}
\paragraph{}
S\={a}ttan\={a}r was one of the forty-nine professors of the Madura College. His father's name \textbf{Seyirkaviriyar} is always prefixed to his own; but the incidents of his history cannot be ascertained. The following is the impromptu which he uttered when the \emph{Kural} passed the College bench:---
\begin{quotation}\small
ஆவனவுமாகாதனவுமறிவுடையார்‌

யாவரும்வல்லாரெடுத்தியம்பத்‌ --- தேவர்

திருவள்ளுவர்தாமுஞ்செப்‌பியவேசெய்வார்‌

பொருவிலொழுக்கம்பூண்டார்‌.
\end{quotation}
\begin{quotation}\small
\emph{The learned will now be able to teach out of the book of the divine Tiruvalluver, that which ought to be done, and that which ought not, and men of good conduct will follow it.}
\end{quotation}

\section{Sengkund\={u}rkizh\={a}r. --- செங்குன்றூர்க்கிழார்.}
\paragraph{}
Sengkund\={u}rkizh\={a}r was one of the forty-nine professors of the Madura College, distinguished by the honorary title of \emph{Tamil-asiriyer}, or Master of Tamil poetry. Nothing more is known of his life, except that he was the author of the encomiastic impromptu on the \emph{Kural}, which we quote below.
\begin{quotation}\small
புலவர்திருவள்ளுவரன்றிப்‌பூமேற்‌

சிலவர்புலவரெனச்செப்ப --- நிலவு

பிறங்கொளிமாலைக்கும்‌ பெயர்மாலைமற்றுங்‌

கறங்கிருண்மாலைக்கும்பெயர்‌.
\end{quotation}
\begin{quotation}\small
\emph{To call any one a poet upon this earth besides the divine Valluver, would be like calling both the evening illumined by the moon, and the evening shrouded in darkness, a fine evening.}
\end{quotation}

\section{Seyal\={u}rkodunchengkannan\={a}r. --- செயலூர்க்கொடுஞ் செங்கண்ணனார்.}
\paragraph{}
This poet was one of the forty-nine professors of the Madura College; but neither his parentage, nor the history of his life is known, and moreover of his compositions only the following stanza, which he uttered on the \emph{Kural}, has been preserved.
\begin{quotation}\small
வேதப்பொருளை விரகால்விரித்துலகோ

ரோதத்தமிழாலுரைசெய்தாராதலா

லுள்ளுநருள்ளும்‌ பொருளெல்லாமுண்டென்ப

வள்ளுவர்வாய்மொழிமாட்டு.
\end{quotation}
\begin{quotation}\small
\emph{Valluver has in Tamil ingeniously explained within the narrow compass of his Kural the substance of the Vedas, that the world may learn it; hence they say that the work contains whatever matter it is possible to think of.}
\end{quotation}

\section{Sekizh\={a}r. --- சேக்கிழார்.}
\paragraph{}
Sekizh\={a}r, called also \emph{Arunmozhitever} (அருண்மொழித்தேவர்), was born of a \emph{Vel\={a}la} family at \emph{Kundat\={u}r}, a village in one of the subdivisions of the \emph{Tonda} country, denominated \emph{Puliy\={u}r-kottam}, and was equally eminent as a poet and politician. The king \textbf{Anabh\={a}ya Chola} having heard of his abilities, made him his Prime Minister, conferring on him at the same time the title of \emph{Uttama-ch\={o}la Pallaver} (உத்தமசோழப்பல்லவர்). While he was employed at Court, he observed that the king though professedly a \emph{Saiva}, yet admiring the \emph{J\={i}vaka Chint\={a}mani} (சீவகசிந்தாமணி), an epic poem belonging to the \emph{Jainas}, made it the sole object of his study. This grieved him very much, and one day having found an opportunity of talking to the king on the subject, he expostulated with him on the impropriety of his conduct in reading an heretical work. The king took his expostulation in good part; but asked him whether he could name any \emph{Saiva} poem which possessed superior merits. He thereupon recited before the king \textbf{Nambi's} \emph{Ant\={a}di} (a poem in 100 stanzas treating of the lives and action of the sixty-three \emph{Tirutonders}, or special votaries of \emph{Siva}). The king was quite pleased with it but thinking it to be too concise, he desired him to enlarge it into an epic poem of as great a length as the \emph{Chint\={a}mani}, under the title of \emph{Tirutonder Puranam} (திருத்தொண்டர்புராணம்). With the king's permission he immediately left the Court and proceeding to Chillambaram took up his abode in the \emph{Mandapam} or open building opposite to the temple there, and began writing the poem and finished it in the course of some time. It must have no doubt been a laborious task, for it consists of 3363 stanzas, arranged in 72 cantos. When the king heard of the completion of the poem, he proceeded in person to Chillambaram, and having had it submitted to the judgment of a learned assembly, which sanctioned it readily, he placed both the author, and the work on a well caparisoned elephant, went forth in royal procession through the town, and afterwards caused a copy of the poem to be engraved on copper plates and deposited in the sanctuary of the temple. It appears that Sekizhar subsequently resigned his office of Prime Minister, and assuming the ascetic life remained at Chillambaram and died there.

\section{S\={e}n\={a}thir\={a}yer. --- சேனாதிராயர்.}
\paragraph{}
S\={e}n\={a}thir\={a}yer was the son of \textbf{Nellein\={a}ther}, the scion of a \emph{Vel\={a}la} family at \emph{Tillipally}, in Jaffna, and traced back his line of ancestors to the time of the \emph{Chakravartis}. He commenced his studies early in life, and distinguished himself as a poet of considerable attainments; but his circumstances in life not permitting him to devote himself exclusively to the Muses, he betook himself to the legal profession, and continued to practice as a Proctor first before the Provincial Court, and afterwards before the District Court of Jaffna, and thereby secured for himself a competency, sufficient to maintain his independence. He was an intimate friend of all the Missionaries in Jaffna, and especially of the late Rev. \textbf{Joseph Knight} of Nellore, and was not unacquainted with the doctrines of Christianity; but such were his prejudices in favour of Hinduism, that he died as he had lived, a Hindu in the widest sense of the word. He is known as the author of a poem, entitled \emph{Nallei Venb\={a}} (நல்லைவெண்பா), which he composed in praise of the \emph{Saiva} shrine at Nellore, and likewise as one of the compilers of the Tamil Dictionary, which was published by the American Missionaries in Jaffna. His death took place in 1840, and a contemporary poet has commemorated that event in the following quatrain:
\begin{quotation}\small
நானாதிராமிடமுநன்னிலக்கண்ணாருறச்செய்‌

சேனாதிராயனையோ செத்ததென்றீர் --- வானாதி

பொன்னிலத்துமுண்மை புகநூலுணர்த்தவனை

யின்னிலத்துவிட்டெடுத்ததே.
\end{quotation}

\section{S\={e}ndan\={a}r. --- சேந்தனார்.}
\paragraph{}
S\={e}ndan\={a}r as represented by himself, was a \emph{Br\={a}hman} of \emph{Ambel}, a town on the banks of the Kaveri, in the Tanjore country, and he acquired the title of \emph{Ub\={a}ya kavi} or the double poet, on account of his skill both in Sanskrit and Tamil poetry. Of his compositions, however, there are none extant, save the \emph{Divakaram}, a Tamil Dictionary comprising 2286 sutras in metre, divided into ten chapters. The period of his existence must be placed before the sixteenth century of the Christian era; for \textbf{Mandalapuruder}, who wrote the \emph{S\={u}d\={a}mani Nigandu} in the early part of that century, professes to have followed the \emph{Divakaram} in composing his work.

\section{Sidambara D\={e}siker. --- சிதம்பரதேசிகர்.}
\paragraph{}
Sidambara D\={e}siker was a \emph{Saiva} ascetic, who studying under the tuition of \textbf{S\={a}ntalinga D\={e}siker} of \emph{Tureiy\={u}r}, became a proficient in the art of poetry, as well as in the science of metaphysics; and soon afterwards had the distinction of being installed into the office of \emph{Tambiran} of the \emph{Matam} at Tiruppor\={u}r. His comments on the five metaphysical treatises of his master are much esteemed; and the poem called \emph{Tiruppor\={u}r sannidhi-murei} (திருப்போரூர்ச்சந்நிதிமுறை), which he wrote in praise of \emph{Skanda}, who is worshipped at Tiruppor\={u}r, is an elegant production.

\section{Sinnatamby Pulaver. --- சின்னத்தம்பிப்புலவர்.}
\paragraph{}
This poet was the son of \textbf{Villavar\={a}ya Mudeliar}, one of the chiefs who revised the \emph{Thesawalamy}, or Codex Jaffanensis, compiled by order of Governor \textbf{Simons} in 1706. His birth took place at \emph{Nellore} in Jaffna. His poetical talents appear to have developed themselves at a very early age; for it is said that while yet a boy, as he was playing a game with areca nuts, with his playmates in the street, a poet, who came over from Southern India to wait on \textbf{Villavar\={a}ya Mudeliar} and to recite before him a poem which he had composed in his praise, not knowing the way to his house, inquired of the children for it; upon which Sinnatamby came forward, and answered him in the following impromptu:---
\begin{quotation}\small
பொன்பூச்சொரியும்‌ பொலிந்தசெழுந்தாதுறைக்கும்‌

நன்பூதலத்தோர்க்கு நன்னிழலா --- மின்‌பிறவை

வீசுபுகழ்நல்லூரான்‌ வில்லவராயன்கனக

வாசலிடைக்கொன்றைமரம்.
\end{quotation}
\begin{quotation}\small
\emph{Lo, there is a Kondei tree\footnote{\emph{Cassia Fistula}.} at the golden gate of Villavar\={a}ya, of Nellore, whose fame is diffused everywhere as vividly as lightening, showering down its golden flowers full of pungent anthers, and affording good shade to the people with its luxuriant branches.}
\end{quotation}
\paragraph{}
He was the author of a poem, entitled \emph{Kalvazheiyant\={a}di} (கல்வழையந்தாதி), consisting of a series of 100 stanzas remarkably beautiful both in their construction, and style. He appears to have been patronized by \textbf{Ganesiyer}, an opulent \emph{Brahman} who lived at Jaffna in his time, and we have the following stanza in which he has commemorated his patron:---
\begin{quotation}\small
கடனந்தனவன கண்டத்தனளனங்கன்கணையால்‌

விடனங்கயிலை மயிலையொத்தாள்விடமிக்கதுத்திப்‌

படனந்தகிமுடி மேல்நின்றநச்சுப்‌ படரரவின்‌

நடனம்புரிதிம்மரசர்கணேசநரேந்திரனே.
\end{quotation}
\begin{quotation}\small
\emph{The damsel with protuberant bosoms and tapering neck when wounded by the shaft of Ananga (Cupid) resembled the peahen of the mount Kailei (meaning Parvati); but thou O Ganesa, the prince of men! resemblest (Krishna), dancing on the extended hood of the venomous serpent.}
\end{quotation}

\section{Sirum\={e}d\={a}viy\={a}r. --- சிறுமேதாவியார்.}
\paragraph{}
The name of this poet occurs as one of the forty-nine professors of the Madura College; but nothing is known of his personal history, and there is no other relic of his poetical effusions than the following impromptu on the \emph{Kural}:---
\begin{quotation}\small
வீடொன்றுபாயிரநான்குவிளங்கற

நாடியமுப்பத்து மூன்றொன்றூழ்‌ --- கூடுபொரு

ளெள்ளிலெழுபதிருபதிற்றைந்தின்பம்‌

வள்ளுவர்சொன்னவகை.
\end{quotation}
\begin{quotation}\small
\emph{Of the number of Chapters in the Kural four pertain to the preface, thirty-three treat of virtue, one of fate, seventy of wealth, and twenty-five of pleasure, and such is the arrangement of Valluver's work.}
\end{quotation}

\section{Sivagn\={a}na D\={e}siker. --- சிவஞானதேசிகர்.}
\paragraph{}
Sivagn\={a}na D\={e}siker, who was the \emph{Tambir\={a}n} of the \emph{Matam} at Tiruv\={a}vaduturei in Tinnevelly, deserves notice both as a metaphysician and poet. The date of his birth as well as that of his death has not been ascertained; but he was contemporary with \textbf{Sankara Namasiv\={a}yer}, and like him prosecuted his studies under the guidance of \textbf{S\={a}minada Desiker}. Of his works the \emph{S\={u}travirutti} (சூத்திரவிருத்தி), an exposition of the \emph{Tolkappiyum}, and the \emph{Nann\={u}l Urei} (நன்னூலுரை), an exposition of the \emph{Nann\={u}l}, are held in high estimation, as the best helps for students in acquiring a correct knowledge of belles lettres. His \emph{Illakanavilakka-ch\={u}r\={a}vali} (இலக்கணவிளக்கச்சூறாவளி), a critique on the \emph{Illakana-vilakkam} of \textbf{Vyttyan\={a}tha N\={a}valer}, is very ably written; but we do not admire the spirit which dictated it, for whilst the faults of the work are exposed with excessive rigour, its merits are studiedly concealed.

\section{Sivaprak\={a}sa D\={e}siker. --- சிவப்பிரகாசதேசிகர்.}
\paragraph{}
Sivaprak\={a}sa D\={e}siker was the eldest of the three sons of \textbf{Kum\={a}rasw\={a}mi Pand\={a}ram}, a \emph{Jangama} priest, who lived at \emph{Kanjipuram} in the early part of the seventeenth century of the Christian era. Losing his father at an early age, and whilst still pursuing his studies, he went with his mother and brothers to Trinomali to settle himself there; but he shortly after set out to Tinnevelly for the purpose of perfecting himself in grammar and poetry. On his way, however, he was delayed at Tirumangalam for a period of two years, in consequence of \textbf{Ann\={a}malei Reddi}, the \emph{Gram\={a}dhipati},\footnote{The head of the village.} of Tirumangalam, having prevailed on him by his entreaties to reside in a \emph{Matam} which he had erected there. Afterwards taking leave of \textbf{Ann\={a}malei Reddi}, who presented him with a sum of 300 pons (equivalent to £52 10), he proceeded to Tinnevelly in company with his brothers. When he reached Tinnevelly, hearing that \textbf{Velliyambala Tambiran}, the head ascetic of the \emph{Saiva Matam} at Sindup\={u}ndurei on the banks of the T\={a}mbraparni, was the best of Tamil scholars, he with his brothers presented himself before him, and asked his instruction in grammar and poetry. The \emph{Tambiran} in order to know the progress he may have already made in composing verses, desired him to compose a verse beginning and ending with the letter (கு) \emph{ku}, and having in the middle the name of \emph{\={U}rudeiyan} (Siva), upon which he uttered the following offhand:---
\begin{quotation}\small
குடக்கோடுவானெயிறு கொண்டார்க்குக்கேழன்‌

முடக்கோடுமுன்னமணிவார்க்கு --- வடக்கோடு

தேருடையான்றெவ்வுக்குத்‌ தில்லைதோன்மெற்கொள்ள

லூருடையானென்னுமுலகு.
\end{quotation}
\paragraph{}
The \emph{Tambiran} was so much satisfied with the performance, that he embraced him, and offering him an honourable seat, told that he did not require any further instruction, but that he would gladly teach his brothers. When his brothers had finished their course of studies, he offered the 300 pons which he had to the \emph{Tambiran} as a remuneration for his trouble, but this offer was declined. The Tambiran only required in return that he should go to Trichend\={u}r, and by his poetry defeat and put to shame a poet, who had set himself up there as the \emph{Tambiran's} rival, and had lampooned him. He gladly undertook the task and repaired to Trichend\={u}r as required, and found the poet in the temple there. No sooner had the object of his visit been reported to the poet than he challenged him to a contest, in which each was to improvise thirty \emph{Yamakama} stanzas (a kind of stanza consisting of four lines beginning with the same word but having a different sense in each line) within a given time, and that the vanquished was to become the slave of the victor. The challenge being accepted, they proceeded to the trial of their skill, but the poet failing delivered himself over as a slave to his competitor, who thereupon conducted him to Sindup\={u}nturei, and consigned him to the \emph{Tambiran}. The \emph{Tambiran} was so much gratified with the success of Sivaprak\={a}sa Desiker, that unwilling to part with him, he endeavoured to make him his associate in the \emph{Matam} at Sindup\={u}nturei; but Sivaprak\={a}sa Desiker preferring to live by himself in solitary dignity took his farewell of the \emph{Tambiran} and returned to his own \emph{Matam} at Tirumangalam. The fame of his superior attainments soon spread over the country, and all persons wishing to avail themselves of his instruction flocked to him from all parts. His time in the latter part of his life appears to have been wholly engrossed in philosophical and metaphysical researches, and in imparting their result to the public by means of a series of poetical treatises. His useful career was however cut short by death, which occurred at Nalkattur in the thirty-second year of his age, as he was on a visitation to the \emph{Matam} there. It is said that he was a staunch defender of celibacy, and that when his friend \textbf{Ann\={a}malei Reddi} proposed to him a marriage, he disdainfully rejected the proposal observing, that he would rather be possessed of a devil than be possessed of a wife.

Sivaprak\={a}sa Desiker was a prolific writer, and the following works all of which are written in well measured poetry, may be mentioned as his productions:

1.\emph{Prabulinga L\={i}la}, பிரபுலிங்கலீலை.

2.\emph{Tiruk\={u}va Puranam}, திருக்கூவப்புராணம்.

3.\emph{Sitt\={a}ndha Sik\={a}mani}, சித்தாந்த சிகாமணி.

4.\emph{V\={e}d\={a}nta Sud\={a}mani}, வேதாந்த சூடாமணி.

5.\emph{Sivaprak\={a}sa Vik\={a}sam}, சிவப்பிரகாசவிகாசம்.

6.\emph{Sivan\={a}ma Mahimei}, சிவநாமமகிமை.

7.\emph{Tarkaparip\={a}shei}, தர்க்கபரிபாஷை.

8.\emph{Sh\={o}nasaila M\={a}lei}, சோணசைலமாலை.

9.\emph{V\={e}ngkai Kalambagam}, வேங்கைக்கலம்பகம்.

10.\emph{Vengkai K\={o}vei}, வெங்கைக்கோவை.

11.\emph{Vengkai Ul\={a}}, வெங்கையுலா.

12.\emph{Vengkai Alangk\={a}ram}, வெங்கையலங்காரம்.

13.\emph{Tiruchendil Ant\={a}di}, திருச்செந்திலந்தாதி.

14.\emph{Sat\={a}mani M\={a}lei}, சதமணிமாலை.

15.\emph{N\={a}lver N\={a}nmani M\={a}lei}, நால்வர் நான்மணிமாலை.

16.\emph{Niranjana M\={a}lei}, நிரஞ்சனமாலை.

17.\emph{Kaitala M\={a}lei}, கைத்தலமாலை.

18.\emph{Abish\={e}ga M\={a}lei}, அபிஷேகமாலை.

19.\emph{Ishtalinga Perungkazhinedil}, இஷடலிங்கப் பெருங்கழில்நெடில்.

20.\emph{Ishtalinga Kurungkazhinedil}, இஷடலிங்கக்குறுங்களில்நெடில்.

21.\emph{Nanneri}, நன்னெறி.

Of these the Nos. 1 and 2 are epic poems, one treating of the history of \emph{Allama Prabhu}, the founder of the \emph{Jangama} sect and the other of the legends of the \emph{Saiva} fane at \emph{Tirukk\={u}vai}. The Nos 3, 4, 5, 6 and 7 treat of the Hindu system of Metaphysical philosophy, and the Nos. 8, 9, 10, 11, 12, 13, 14, 15, 16, 17, 18, 19 and 20 of the praise of certain \emph{Saiva} shrines, while the No. 21 contains a series of moral apothegms, from which we quote below a few as specimen:
\begin{quotation}\small
பெரியவர்தந்நோய்போற்‌ பிறர்நோய்கண்டுள்ள

மெரியினிழுதாவரென்க --- தெரியிழாய்‌

மண்டுபிணியால்‌ வருந்துபிறவுறுப்பைக்

கண்டுகலுழுமேகண்‌.

பேரறிஞர்தாக்கும்‌ பிறர்துயரந்தாங்கியே

வீரமொடுகாக்க விரைகுவார்‌ --- நேரிழாய்‌

மெய்சென்றுதாக்கும்‌ வியன்கோலடிதன்மேற்‌

கைசென்று தாங்குங்கடிது.

தங்குறைவுதீர்வுள்ளார்‌ தளர்ந்துபிறர்க்குறூஉம்

வெங்குகைதீர்க்கிற்பார்‌ விழுமியோர் --- திங்கள்‌

கறையிருளைநீக்கக்‌ கருதாதுலகி

னிறையிருளைநீக்குமேனின்று.
\end{quotation}
\begin{quotation}\small
\emph{The worthy feel the griefs of others as their own griefs as butter melteth in the fire; thus, O thou who art ornamented with choice jewels! the pain caused to a limb by a local disorder the eye beholds and weeps.}

\emph{The truly wise hasten to sustain the griefs sustained by others and to protect them courageously, O thou who art adorned by polished jewels! thus, the arm receives on itself and sustains the blow of the club which the body would otherwise sustain.}

\emph{The worthy think not of supplying their own wants, but endeavour to supply the pressing wants of others; thus the moon, heeding not the removal of the darkness of its own spots, chaseth away the darkness which filleth the world.} \textbf{Ellis.}
\end{quotation}
\paragraph{}
Sivaprak\={a}sa Desiker was also the author of three different panegyrics on \textbf{B\={a}lay\={a}nanda Sw\={a}mi}, who was his \emph{Guru} or spiritual instructor, besides a hymn on each of the \emph{Saiva} shrines in southern India, which were commemorated in the chants of \textbf{Apper}, \textbf{Sundarer} and \textbf{Sampanter} in ancient times. He undertook to perfect the \emph{K\={a}latti Pur\={a}nam} (காளத்திபுராணம்), which his brother \textbf{Karunaiprak\={a}sa Deziker} had left unfinished; but he died before he could execute his task.

\section{Sivav\={a}kiyar. --- சிவவாக்கியர்.}
\paragraph{}
This eminent philosopher is known as the author of a monotheistical treatise bearing his own name, and in which he has attacked and exposed the Hindu religion, and its concomitant superstitions. The era of his existence is variously stated; but it cannot be fixed anterior to the incursion of the Muhammedans into the Dekkan, as he alludes to their creed in his treatise.

As a specimen of the style and tenets of Sivav\={a}kiyar, we subjoin the five following stanzas, and their English translation by the late Mr. \textbf{Ellis.}
\begin{quotation}\small
பண்டுநான்பறித்தெறிந்த பன்மலர்களெத்தனை

பாழிலேசெபித்துவிட்ட மந்திரங்களெத்தனை

மிண்டனாய்த்திரிந்தபோதிறைத்தநீர்களெத்தனை

மீளவுஞ்சிவாலயங்கள்‌ சூழவந்ததெத்தனை

அண்டர்கோனிருப்பிடமறிந்துணர்ந்தஞானிகள்‌

கண்டகோயிற்றெய்வமென்று கையெடுப்பதில்லையே.

நீரையள்ளி நீரில்விட்டு நீர்நினைந்தகாரிய

மாரையுன்னிநீரெல்லாமவத்திலேயிறைக்கிறீர்‌

வேரையுன்னிவித்தையுன்னி வித்திலேமுளைத்தெழுந்த

சீரையுன்னவல்லிராய்‌ சிவபதங்கள்சேரலாம்‌.

அரியுமல்லவரனுமல்லவயனுமல்லவப்புறங்

கருமைவெண்மைசெம்மையுங் கடந்துநின்றகாரணம்‌

பெரியதல்லசிறியதல்ல பெண்ணுமாணுமல்லவே

துரியமுங்கடந்துநின்ற தூரதூரதூரமே.

சாதியாவதேதடா சலந்திரண்டநீரல்லோ

பூதமைந்துமொன்றல்லோ புலன்களைந்துமொன்றல்லோ

காதுவாழிகாறைகம்பி சூடகம்பொன்னொன்றல்லோ

சாதிபேதமோதுகின்ற தன்மையென்னதன்மையே.

கறந்தபால்முலைபுகா கடைந்தவெண்ணெய்மோர்புகா

வுடைந்தசங்கினோசையு முயிர்களுமுடல்புகா

வடிந்தவோரடையினியுதிர்ந்தபூமரம்புகா

விறந்துபோனமானிடரினிப்பிறப்பதில்லையே.
\end{quotation}
\begin{quotation}\small
\emph{Formerly how many flowers have I gathered and scattered.}

\emph{How many prayers have I repeated in a vain worship?}

\emph{While yet in the prime of my life, how much water have I poured out?}

\emph{And, moreover, how often have I encompassed the holy places of Siven,}

\emph{This I have left off, for the wise who know the true God, the Lord of heavenly beings.}

\emph{Believe not the Idol of the temples apparent to the eyes to be God, nor lift up to it their hands.}

\emph{While taking up the water and throwing it again into the water (in performing the sandhya and other rites) what is the object on which you think?}

\emph{On whatsoever you think you have thrown all the water vainly:}

\emph{Think on the root, think on the seed, and on the benefit arising from that seed;}

\emph{When you are thus able to think, you may approach the feet of God.}

\emph{It is not Ari, it is not Aren, it is not Ayen;}

\emph{Far beyond the black (the colour of Vishnu) the white (the colour of Siven) or the red (the colour of Brahma) soars the everlasting cause;}

\emph{It is not great, it is not small, neither is it male nor female:}

\emph{Beyond every state of corporeal being it is farther, farther, and farther still.}

\emph{What, O wretch, is caste? is not water an accumulation of fluid particles?}

\emph{Are not the five elements and the five senses one?}

\emph{Are not the several ornaments for the neck, the breast, and the feet equally gold?}

\emph{What then is the peculiar quality supposed to result from difference in caste?}

\emph{As milk once drawn cannot again enter the udder, nor butter churned be recombined with milk;}

\emph{As sound cannot be produced from a broken conch, nor the life be restored to its body;}

\emph{As a decayed leaf and a fallen flower cannot he reunited to the parent tree;}

\emph{So a man once dead is subject to no future birth.}
\end{quotation}

\section{S\={i}niv\={a}sakar. --- சீனிவாசகர்.}
\paragraph{}
S\={i}niv\={a}sakar was a Brahman poet, who flourished in the Ramnad country. He is celebrated as the composer of a collection of \emph{Nautch} songs, much admired for their melody. Neither the time of his birth nor death is known; but he appears to have been alive in the beginning of the present century, enjoying the patronage of \textbf{Vellei Marudu Poligar} of Sherewele, who was hanged for rebellion against the English Government in \textsc{ce} 1801.

\section{S\={i}ttaleich\={a}ttan\={a}r. --- சீத்தலைச்சாத்தனார்.}
\paragraph{}
This poet was one of the forty-nine professors of the Madura College. It is related that having been in the habit of striking his head with his stylus, whenever he found any fault in the works submitted by their authors to the judgment of the College bench, he brought on a headache which tormented him all his life: hence his fellow professor \textbf{Marutuven T\={a}motharan\={a}r} while applauding the \emph{Kural}, jocosely observed, that other people got rid of their headache by smelling the salt produced from the stalk of the \emph{S\={i}ndil},\footnote{Menispermum cordifolium.} and the sliced dry ginger mixed with honey; but S\={i}ttaleich\={a}ttan\={a}r got rid of his headache, by hearing the \emph{Kural} recited, because it contained no faults. The following is the impromptu which he uttered in praise of the \emph{Kural}:---
\begin{quotation}\small
மும்மலையுமுன்னாடு முந்நதியுமுப்பதியு

மும்முரசுமுத்தமிழு முக்கொடியு --- மும்மாவுந்

தாமுடையமன்னர்‌ தடமுடிமேற்றாரன்‌றோ

பாமுறைதேர்வள்ளுவர்முப்பால்‌.
\end{quotation}
\begin{quotation}\small
\emph{The three parts (of the Kural) composed by Valluver, are each a garland on the crowns of the three kings (Chera, Chola and P\={a}ndiya), who are respectively possessed of the three mountains (Kolli, Neri and Pothiya), the three countries (Chera, Chola and P\={a}ndiya), the three rivers, (Porunai, Caveri and Vygai), the three cities (Kar\={u}r, Ureiy\={u}r and Madura), the three kinds of drums (Mangala-murasu, Vettimurasu and Kodei-murasu) the three dialects of Tamil (colloquial, poetical and dramatic), the three banners (Vitkodi, Pulikodi and Minakodi), and the three horses (Kanavattam, P\={a}dalam and Koram).}
\end{quotation}

\section{Subrahmanya V\={e}diyer. --- சுப்பிரமணியவேதியர்.}
\paragraph{}
This poet was a member of a \emph{Brahman} family at \emph{Alvar Tinnevelly}, in the Carnatic. The time of his birth, as well as that of his death is not known; but he was contemporary with \textbf{Samin\={a}da Tambiran}, who makes honourable mention of him in his Introduction to the \emph{Illakkanakottu}. He wrote the \emph{Pray\={o}ga-Viv\={e}gam} (பிரயோக விவேகம்), a treatise highly esteemed on account of its being a great help in the composition of poems.

\section{Sundarer. --- சுந்தரர்.}
\paragraph{}
Sundarer was a contemporary of \textbf{Apper} and \textbf{Sampanter} and like them a poet, as well as a zealous champion for the \emph{Saiva} faith against the \emph{Jainas}. His real name was \emph{Nambiy\={a}rurer}, which he afterwards changed for that of Sundarer. He was born of a \emph{Br\={a}hman} family of the \emph{Adi Saiva} sect at \emph{Tirunavalur}; in the Carnatic. Whilst yet a child, the king \textbf{Narasinghamuneiyer} with the consent of his parents, took him into his family and adopted him, giving him a good education in all branches of learning. When he came of age, his father by the kings command negotiated a marriage for him with the daughter of a \emph{Brahman} at Putt\={u}r; but just on the day fixed for its celebration he broke off the match, and assuming the ascetic life visited the different \emph{Siva} fanes in Southern India, and hymned their praises in a series of 37,000 \emph{Pathigams}, each consisting of ten stanzas. It is however certain that he did not preserve continence which is manifest from his amorous intrigues with the two courtezans \textbf{Paravei} and \textbf{Sangkili} as related in the \emph{Tirutonder Pur\={a}nam}. He died at \emph{Tiruvanji}, in the Chera country, in the eighteenth year of his age, whilst he was visiting a temple there in company with the king \textbf{Cheraman Peruman}, whose friendship he enjoyed. In our opinion as the date given in the \emph{Cholap\={u}rva Patayam} for the accession of \textbf{Cherama Peruman} seems to admit of no doubt, we may therefore place the period of the existence of Sundarer and his two fellow Champions in the fifth century of the Christian era for a certainty and thereby clear it from the monstrous chronology of the Puranas.

\section{S\={u}dikodutta N\={a}chiy\={a}r. --- சூடிக்கொடுத்தநாய்ச்சியார்.}
\paragraph{}
S\={u}dikodutta N\={a}chiy\={a}r, otherwise called \emph{\={A}nd\={a}l}, was the poetess who composed the poems called \emph{Tirup\={a}vei} (திருப்பாவை), and \emph{Tirumozhi} (திருமொழி), consisting of 173 stanzas in praise of \emph{Vishnu}, and which now form a part of the \emph{N\={a}l\={a}yiraprabhandam} (நாலாயிரபிரபந்தம்). Her parentage is not known; but if we may accept as facts that which appears not improbable in her legend, she was a foundling brought up by a \emph{Vaishnava} devotee at \emph{Villiputt\={u}r}, and having even from her childhood become an enthusiast in the faith of her foster-father, after she had passed her teens she proceeded to Tirupati, and consecrated herself to \emph{Vishnu}, who is worshipped there.

\section{Tandi \={A}siriyer. --- தண்டியாசிரியர்.}
\paragraph{}
Tandi \={A}siriyer was one of the poets who flourished at the Court of \textbf{Bhoja}, king of \emph{Dh\={a}rapuram}, in the twelfth century of the Christian era. As he was skilled both in the Sanskrit and Tamil poetry, he obtained the title of \emph{Ubaya Kavi} or a poet in two languages; and it is said that he ranked foremost of all the poets of his day, except \textbf{K\={a}lid\={a}sa}. Of his works, we have two treatises, one in Sanskrit, entitled \emph{K\={a}vyadersana} (காவியதரிசனம்), and the other in Tamil, under the name of \emph{Alangk\={a}ram} (அலங்காரம்), both of which treat of poetical tropes and figures.

\section{Tatvar\={a}yer. --- தத்வராயர்.}
\paragraph{}
Tatvar\={a}yer was an eminent poet and metaphysician, born of a \emph{Br\={a}hman} family of the \emph{M\={a}dhavach\={a}rya} sect at \emph{Vira}, in the Carnatic, sometime in the sixteenth century of the Christian era. Having very early completed his course of studies in Sanskrit and Tamil in his own village, he assumed the ascetic life at the same time with \textbf{Sv\={a}rup\={a}nander}, who was his kinsman as well as fellow-student. The young men then desirous of having themselves initiated into the mysteries of \emph{Saivaism} proposed to go in search of a \emph{Guru} or teacher; but they agreed before starting, that he who succeeded in his undertaking, should become the \emph{Guru} of the one who failed. Tatvar\={a}yer went towards the north, and \textbf{Sv\={a}rup\={a}nander} towards the south. Tatv\={a}rayer failed in finding a \emph{Guru}; but \textbf{Sv\={a}rup\={a}nander} was fortunate enough to find one in the person of \textbf{Sivaprak\={a}ser}, a \emph{Saiva} recluse who lived in a thicket at Kovattam, and according to the agreement Tatvar\={a}yer became the disciple of \textbf{Sv\={a}rup\={a}nander}. As a poet Tatvar\={a}yer excelled his master, and hence the latter employed him in versifying the substance of the \emph{Vedanta} and \emph{Saiva} S\={a}stras. The following is a list of the works which he composed, viz.

1.\emph{Sivaprak\={a}sa Venb\={a}}, சிவப்பிரகாசவெண்பா.

2.\emph{Tatv\={a}mirtam}, தத்வாமிர்தம்.

3.\emph{Tirutt\={a}l\={a}ttu}, திருத்தாலாட்டு.

4.\emph{Pilleittirun\={a}mam}, பிள்ளைத்திருநாமம்.

5.\emph{Venb\={a}vant\={a}di}, வெண்பாவந்தாதி.

6.\emph{Kalittureiyant\={a}di}, கலித்துறையந்தாதி.

7.\emph{Sinnap\={u} Venb\={a}}, சின்னப்பூவெண்பா.

8.\emph{Tas\={a}ngam}, தசாங்கம்.

9.\emph{Irretteimani Malei}, இரட்டைமணிமாலை.

10.\emph{Mummani K\={o}vai}, மும்மணிக்கோவை.

11.\emph{N\={a}nmani M\={a}lei}, நான்மணிமாலை.

12.\emph{Tiruvadi M\={a}lei}, திருவடிமாலை.

13.\emph{Gn\={a}navin\={o}dan Kalambagam}, ஞானவினோதன்கலம்பகம்.

14.\emph{Ul\={a}}, உலா.

15.\emph{Siledei Ul\={a}}, சிலேடையுலா.

16.\emph{Nenjuvidud\={u}tu}, நெஞ்சுவிடுதூது.

17.\emph{Kalimadel}, கலிமடல்.

18.\emph{Agnaivadeipparani}, அஞ்ஞைவதைப்பரணி.

19.\emph{M\={o}gavadeipparani}, மோகவதைப்பரணி.

20.\emph{Amirtas\={a}ra Venb\={a}}, அமிர்தசாரவெண்பா.

21.\emph{Tiruvarudkazhen M\={a}lei}, திருவருட்கழன்மாலை.

22.\emph{P\={o}ti M\={a}lei}, போற்றிமாலை.

23.\emph{Pugadchi M\={a}lei}, புகழ்ச்சிமாலை.

24.\emph{Sasivarnap\={o}tham}, சசிவர்ணபோதம்.

25.\emph{Peruntirattu}, பெருந்திரட்டு.

26.\emph{Kuruntirattu}, குறுந்திரட்டு.

\section{T\={a}ndavam\={u}rti. --- தாண்டவமூர்த்தி.}
\paragraph{}
T\={a}ndavam\={u}rti was a \emph{Saiva} ascetic who obtained considerable reputation by his skill in the Sanskrit and Tamil languages, and their respective literature. In the latter part of his life he devoted his attention entirely to metaphysical researches and wrote a work on the subject, entitled \emph{Kaivalya-navan\={i}tam} (கைவல்யநவநீதம்), consisting of 293 stanzas.

\section{T\={a}yum\={a}naver. --- தாயுமானவர்.}
\paragraph{}
T\={a}yum\={a}naver was born at \emph{Trichinapoly}, in the Carnatic about 150 years ago. His father \textbf{Kediliyappa Pillei} belonged to the \emph{Vel\={a}la} caste, and was \emph{Samprati} or accountant to the king \textbf{Vijaya Reghunatha Chokalinga Nay\={a}ker}. Having through the endeavours of his father skilled himself in the Sanskrit and Tamil languages, he devoted his attention to the Hindu Philosophy, and under the guidance of a famous \emph{Guru} of the name of \textbf{Maunadesiker}, soon obtained a knowledge of its esoteric doctrines. From his earliest years he had a predilection for a religious life; but on the death of his father he was prevailed upon by the king's importunity to accept the Sampratiship, and employ himself at the Court. He however did not continue long in his office; for when on the death of the king, the queen conceiving a criminal passion for him endeavoured to entice him, he fled from her territory, and took up his residence in the Ramnad country. Some time afterwards he entered into marriage with a female of his own class, but she dying after giving birth to her first child, he assumed the habit of a half-naked Sanniy\={a}si, and led a strolling life until his death; visiting successively the \emph{Saiva} shrines throughout Southern India, and hymning their praises.

As a philosopher T\={a}yum\={a}naver is held in great esteem by the followers of the pantheistical school, and the doctrines which he held respecting God and the universe may be summed up in the following lines of \textbf{Pope}:---
\begin{quotation}\small
``\emph{All are but parts of one stupendous whole,}

\emph{Whose body Nature is, and God the soul;}

\emph{Warms in the sun, refreshes in the breeze,}

\emph{Glows in the stars, and blossoms in the trees,}

\emph{Lives through all life, extends through all extent,}

\emph{Spreads undivided, operates unspent.}''
\end{quotation}
\paragraph{}
As a poet, he certainly claims an undisputed pre-eminence over all his contemporaries. His spiritual songs, of which a collection consisting of 1453 stanzas of various metres was published at Madras in 1836, afford ample proof of his superior talents; for there are scarcely any of the kind in Tamil to equal them in sublimity of style.

\section{Tirumangai \={A}zhv\={a}r. --- திருமங்கையாழ்வார்.}
\paragraph{}
This poet was a native of \emph{Trival\={u}r}, and the \emph{Vaishnavas} venerate him as one of the twelve \emph{\={A}zhv\={a}rs}, who was the foremost in promoting the interests of their sect, in opposition to the \emph{Saivas} and \emph{Jainas}. His father \textbf{Nila} held the post of \emph{Sen\={a}pati} or generalismo of the \emph{Chola} army, and he himself was viceroy of a district; but failing to pay the tribute, and resisting the force which the king had sent to compel payment, he was deprived of his vice-regal authority, and left destitute, upon which he turned out a freebooter, but robbing only those who had not the \emph{Vaishnava} marks on their foreheads. He obtained the title of \emph{N\={a}tkaviperumal} from the skill which he displayed in the composition of the four kinds of poems,\footnote{See Rotler's Tamil and English Dictionary under the word கவி.} and it is said that when \textbf{Sampanter}, the famous champion of the \emph{Saiva} sect, disputed his right to the title he proved himself worthy of it to his satisfaction. His contribution to the \emph{N\={a}l\={a}yiraprabhandam} (நாலாயிரபிரபந்தம்), consists of 1253 stanzas, in which he has not only eulogized \emph{Vishnu}, but has likewise explained the esoteric doctrines of his sect.

\section{Tirumazhishai \={A}zhv\={a}r. --- திருமழிசையாழ்வார்.}
\paragraph{}
The history of Tirumazhishai \={A}zhv\={a}r, who is reckoned as one of the twelve \emph{\={A}zhv\={a}rs} like the preceding, has been altogether rendered a myth. This much however may be certain, that he was a foundling who was brought up in a \emph{Veda} family at \emph{Tirumazhishai}. Having enthusiastically wedded himself to the \emph{Vaishnava} faith in his youth, he continued a rancorous opponent to the \emph{Saivas} throughout his life, and by his skill in polemics made many proselytes amongst them. He appears to have been a poet of superior talents, in proof of which may be mentioned his contribution of 206 stanzas to the \emph{N\={a}l\={a}yiraprabhandam} (நாலாயிரபிரபந்தம்), in praise of \emph{Vishnu}.

\section{Tirum\={u}ler. --- திருமூலர்.}
\paragraph{}
Little more is known of Tirum\={u}ler than that he was a \emph{Saiva} ascetic and philosopher, contemporary with \textbf{Agastiyer}, and like him exerted himself in diffusing the worship of \emph{Siva} amongst the Tamils. All that is related of him in the \emph{Tirutonder Puranam}, in the character of one of the sixty-three special votaries of \emph{Siva}, is manifestly fabulous, except that he was a cow-herd at \emph{S\={a}tan\={u}r}, in the Carnatic, previous to his assuming the ascetic life, and establishing himself in the \emph{Saiva} fane at Tiruvavaduturai. Dr. \textbf{Ainslee} notices a medical treatise of Tirum\={u}ler under the title of \emph{Tirumuler Vytia V\={a}gadam} (திருமூலர்வைத்தியவாகடம்); but he is chiefly celebrated for his \emph{Tirumulamandiram} (திருமூலமந்திரம்), a voluminous work containing an exposition of the tenets of the \emph{Saiva} religion, especially with reference to the four different forms of worship, as enjoined in the \emph{\={A}gamas}. In the following two stanzas which we quote from it, the reader will see the striking coincidence between the doctrine which they inculcate, and that of the Bible:---
\begin{quotation}\small
அன்புஞ்சிவமுமிரண்டென்பாரறிவிலா

ரன்பேசிவமாவதியாருமறிகிலா

ரன்பேசிவமாவதியாருமறிந்தபி

னன்பேசிவமாயமர்ந்திருப்பாரே.

என்பேவிறகாவிறைச்சியறுத்திட்டுப்‌

பொன்போற்கனலிற்‌ பொரியவறுப்‌பினு

மன்போடுருகியகங்குழைவார்க்கன்றி

யென்போன்மணியினையெய்தவொண்ணாதே.
\end{quotation}
\begin{quotation}\small
\emph{The ignorant think that God and love are different; none knows that God and love are the same; did all men know that God and love were the same, they would dwell together in peace, considering love as God.} Compare, 1. John, 4. 16.

\emph{To those of soft hearts, whose minds are melted by divine love, although their flesh be cut off their bones used as fuel, and their moisture dried up by wasting in the golden flame, and to those alone it is not forbidden to approach the God, who is the golden jewel of my soul.} Compare 1. Cor. 13. 1-3.
\end{quotation}

\section{T\={i}run\={i}lakanda Y\={a}zhp\={a}ner. --- தீருநீலகண்ட யாழ்ப்பாணர்.}
\paragraph{}
T\={i}run\={i}lakanda Y\={a}zhp\={a}ner was a minstrel who flourished at \emph{Erukkattampuliy\={u}r}, in the Carnatic, during the reign of \textbf{Kuna P\={a}ndiyen}. The traditions respecting him are tinged with fable; but there is no doubt that he was well skilled in his profession and being a zealous votary of \emph{Siva}, continued all his life in composing verses in praise of his favourite god, and singing them to the lute. It is said that as he went to Madura to exhibit his performance in the temple there, the \emph{Br\={a}hmans} at first would not take any notice of him; but that when he strung his lute and drew forth a few strains, they were so enraptured that they presented him a seat on a golden board, and rewarded him most munificently. He afterwards repaired to Shiyally and joining \textbf{Sampanter} spent the remainder of his days with him.

\section{Tirup\={a}n \={A}zhv\={a}r. --- திருப்பானாழ்வார்.}
\paragraph{}
Tirup\={a}n \={A}zhv\={a}r is celebrated equally as a poet, and as one of the twelve \emph{\={A}zhv\={a}rs} or special votaries of \emph{Vishnu}. When we strip his legend of its disguise, it would appear that his mother having exposed him in a corn-field at \emph{Ureiy\={u}r} soon after his birth, he was found and adopted by a \emph{P\={a}nen} or minstrel, whose profession he afterwards followed. His poetical compositions are not forthcoming at present, except the stanzas which he contributed to the \emph{N\={a}l\={a}yiraprabhandam} (நாலாயிரபிரபந்தம்), portraying the beauties of \emph{Vishnu} from foot to head.

\section{Tiruvalluver. --- திருவள்ளுவர்.}
\paragraph{}
Tiruvalluver deservedly ranks foremost among the poets and moralists of India; for he has had no superior or equal in either character. Some place the date of his existence in the third or fourth and others in the eighth or ninth century of the Christian era, but the last is generally considered as most probable. All accounts concur in representing him as the offspring of a \emph{Brahman} by a \emph{Pariah} woman, brought up by a \emph{Valluven}\footnote{A priest of the Pariah tribe.} at Mailap\={u}r, near Madras, who had found him exposed in an \emph{Illippe} grove in his neighbourhood. It is not known who directed his studies; but we are told that while yet a youth, he had acquired so great a reputation for learning that \textbf{M\={a}rgasagayen}, a wealthy farmer who lived at K\={a}verip\={a}kam, did not disdain to bestow on him the hand of his only daughter named \textbf{V\={a}suki}, and besides the highest personages in the place felt a pride in the enjoyment of his friendship. He, however, appears to have pursued a quiet and unobtrusive course of life, working at the loom for subsistence as the other Pariahs did until he was prevailed upon by his friends to take up the gauntlet against the professors of the Madura College, who had arrogantly defied all other men of learning in the country and claimed all literary honours to themselves. In order to humble the pride of the professors, he wrote an ethical poem, under the title of \emph{Kural} (குறள்), and repairing to the College asked them to review it as judges of poetical composition, that he may have a seat on their College bench in case it was approved. The professors at first scorned at the idea of a Pariah-born poet offering himself as a competitor for a seat with them, but afterwards when he drew them into a contest and completely vanquished them, they not only approved his work and gave him the desired seat, but also conferred on him several honorary titles, such as, \emph{Deivapulaver} (தெய்வப்புலவர்), the divine poet, \emph{Mudetp\={a}valer} (முதற்பாவலர்), the first of poets, \emph{Perun\={a}veler} (பெருநாவலர்), the great poet, etc. A legend would have it that the professors were so sensible of their discomfiture, that unable to survive it, they all drowned themselves in the pond in the neighbourhood of the College; but this is not confirmed by other accounts. Nothing further is known of Tiruvalluver which can be relied upon, excepting that some time after his return to Mailap\={u}r, he was visited with a severe domestic calamity by the death of his wife to whom he was tenderly attached and it so deeply affected his mind that secluding himself from society, he devoted the remainder of his life to religious contemplation. It is said that while lying sleepless and agitated, on the night following his wife's death, he ejaculated extempore the following verse:
\begin{quotation}\small
அடிசிற்கினியாளை யாக்கஞ்செய்வாளைப்‌

படிசொற்படிகடவாளை, யடிவருடிப்‌

பின்றுஞ்சிமுன்னெழூஉம்‌ பேதையைவிட்டையையோ

வென்றுஞ்சுமென்கண்ணிணை.
\end{quotation}
\begin{quotation}\small
\emph{When I have lost a woman who excelled in the knowledge of housewifery, who performed rightly all domestic duties, who never transgressed my word or my door, who chafed my limbs and never slumbering until I slept, arose before I awoke; Alas! Alas! How can my eyes again know sleep!.}\footnote{See Ellis' Kural, p. 168.}
\end{quotation}
\paragraph{}
The time of his death is uncertain; but a tradition would have it that when he died his body was according to his express desire exposed in the open air outside the town to be devoured by crows.

It is difficult to judge from the tenor of his \emph{Kural} to what sect he belonged; for he has entirely avoided in the work every thing that savours of sectarianism in order to harmonize the suffrages of all the sects. The Jainas, however, claim him to belong to their sect from his having used in one of his distichs in praise of God the epithet \emph{Andanan} (அந்தணன்), which is applicable to \emph{Arh\={a}} (அருகன்), the object of their worship.

Independent of the \emph{Kural}, we have no other compositions of Tiruvalluver. The physiological work, entitled \emph{Gn\={a}navettiyan} (ஞானவெட்டியான்), which is commonly ascribed to him, appears to have been written by some person long after his time; for the author, though he assumes the title of Tiruvalluver and speaks in his character in the body of the work, yet betrays the disguise in one of the introductory stanzas, thus:
\begin{quotation}\small
அகமகிழுமம்பிகைப்‌ பெண்ணருளினாலே

யவனிதனில்‌ஞானவெட்டியருள யானும்‌

நிகழ்திருவள்ளுவனயனாருரைத்தவேத

நிரஞ்சனமாநிலவுபொழிரவிகாப்பாமே.
\end{quotation}
\begin{quotation}\small
\emph{Invoking the spotless one whom the holy Valluva has revealed in his incomparable Veda as clear as the light of the sun and moon, and by the grace of the damsel Ambigai (Parvati), who delights the heart, I shall compose the Gn\={a}na Vettiyan.}
\end{quotation}
\paragraph{}
The \emph{Kural} was first brought to the notice of the European literati by Father \textbf{Beschi} in the \emph{Shen Tamil Grammar}, which he published in 1730, and some portions of it have since been translated into Latin, German and English languages. We subjoin the various opinions of the Tamil scholars amongst Europeans on its merits.

The \emph{Kural} ``stands confessed, even to the present day, to be the best and chief of all compositions in the polished dialect.'' \textbf{Taylor's} Oriental Historical Manuscripts, Vol. 1. p. 177.

``Though more than fifteen hundred years old, this rare collection of precepts, conveyed in the style of unequalled poetry, has lost none of its original favour among the people. In one hundred and thirty-three chapters it treats of almost every variety of subjects pertaining to the relations and duties of life, forming a text-book of indisputable authority.'' \textbf{Ward's} India and the Hindus, p. 135.

``The \emph{Kural} of Tiruvalluver is a poetic work on morals, of great merit as a literary performance, and highly esteemed among the Tamil natives, for the beauty of its language, and the truth of its sentiments.'' \textbf{Hoole's} Personal Narrative of a Mission to the South of India.'' Part 2. p. 311.

``It is called the first of works, from which, whether for thought or language, there is no appeal.'' \textbf{Drew's} Preface to the \emph{Kural}, Vol. 1. p. 2.

``Nothing certainly in the whole compass of human language can equal the force and terseness of the sententious distichs in which the author conveys the lessons of wisdom he utters.'' \textbf{Percival's} Land of the Veda, p. 110.

``The \emph{Kural} of Tiruvalluver, a work consisting of 1333 distichs, or poetical aphorisms, on almost every subject connected with morals and political economy, and which is regarded by all Tamilians (and perhaps justly) as the finest composition of which the Tamil can boast appears to be not only the best but the oldest Tamil work of any extent which is now in existence. \textbf{Caldwell's} Dravidian Grammar, Introduction, p. 85.

\section{Tiruv\={e}rk\={a}du Mutti\={a}h Mudeli\={a}r. --- திருவேற்காடு முத்தையா முதலியார்.}
\paragraph{}
Tiruv\={e}rk\={a}du Mutti\={a}h Mudeli\={a}r was born at \emph{Tiruverk\={a}du}, in the Carnatic, in \textsc{ce} 1761. It appears from an amusing account of his own studies, written by himself, and published in \textbf{Dalrymple's} Oriental Repertory, that he had begun to apply himself to his books at the very early age of five, and before he had completed his twenty-first year, he had acquired a knowledge of the Sanskrit, Tamil, Telugu, Maharatta, Persian, English and Latin languages, and became an adept in the mystical philosophy of the Hindus. The fame of his acquirements soon procured for him the patronage of Lord \textbf{Hobart}, the then Governor of Madras, and through his kindness was admitted into the public service. After the capture of Ceylon by the British, he was sent to Colombo as one of the Dubashes, but he remained in the office only for a short time, and then returned to his country.

Professor \textbf{Wilson}, in his Historical Sketch of the kingdom of Pandiya, calls Muttiah ``an ingenious native'' and quotes him as the author of a ``History of the Modern Kings of Madura.''

Although the few occasional effusions of his muse show that he was possessed of abilities to make him distinguished as a poet, yet he seems never to have set his mind towards the composition of a poem of any magnitude. There are extant five acrostic stanzas addressed by him to Lord \textbf{Hobart}, of which the initial letters of the first line of each when taken together form his own name, and the final letters that of his patron.

\section{Tiruv\={e}nkataiyer. --- திருவேங்கடையர்.}
\paragraph{}
This poet was a \emph{Br\={a}hman} who lived at \emph{Villiput\={u}r}, in the Carnatic. His memory has been preserved by his \emph{Uvam\={a}na Sangraham} (உவமானசங்கிரகம்), a short but elegantly written treatise on the metaphors employed in describing the charms of the fair sex, and which greatly facilitates the composition of erotic poems.

\section{Tiruv\={e}nkatasw\={a}mi. --- திருவேங்கடசுவாமி.}
\paragraph{}
Tiruv\={e}nkatasw\={a}mi, who was a \emph{Br\={a}hman} of \emph{M\={a}dei}, in the Carnatic, distinguished himself both as a poet and metaphysician. The date of his existence, as well as the particulars of his personal history is unknown. Of his works the imitation of the \emph{Prab\={o}da Chandr\={o}dayam} (பிரபோதசந்திரோதயம்) of \textbf{Kishna Misra} in Tamil, redounds greatly to his credit. It consists of 2012 stanzas, divided into forty-eight cantos, but unlike the Sanskrit original which is in the form of a drama, it is of an epic character. In this work the \emph{R\={a}m\={a}yana} is explained as an allegory; the triumph of \textbf{R\={a}ma} over \textbf{R\={a}vana} being represented as the triumph of the good principle over the evil.

\section{T\={e}nikudi Kiran\={a}r. --- தேனிக்குடிக்கீரனார்.}
\paragraph{}
Of the history of T\={e}nikudi Kiran\={a}r nothing further is known than that he was one of the forty-nine professors of the Madura College. His name would have long ago been forgotten, but for his encomiastic stanza on the \emph{Kural}, which we quote below.
\begin{quotation}\small
பொய்ப்பாலபொய்யேயாய்ப்போயினபொய்யல்லாத

மெய்ப்பாலமெய்யாய்‌ விளங்கினவே --- முப்பாலிற்‌

றெய்வத்திருவள்ளுவர்‌ செப்பியகுறளால்‌

வையத்துவாழ்வார்மனத்து.
\end{quotation}
\begin{quotation}\small
\emph{By the Kural, the production of the divine Tiruvalluver, the world has been enabled to distinguish truth from falsehood which were hitherto confounded together.}
\end{quotation}

\section{T\={e}reiyer. --- தேரையர்.}
\paragraph{}
T\={e}reiyer, called also \emph{T\={e}rer}, was one of the twelve disciples of \textbf{Agastiyer}. He is described in the \emph{Sadur Agaradi} as ``a poet thoroughly versed in the \emph{\={A}y\={u}r Veda}, or science of medicine.'' There is a tradition that he displayed so much skill in the healing art, that he excited the jealousy of his master by frequently foiling him and was in consequence dismissed with a malediction which rendered his practice unsuccessful for the future; but this, as well as some other incidents reported concerning him, have much the air of fable. He was, however, a prolific writer on subjects connected with his profession; but all his works have perished, except a treatise on medicine, entitled \emph{Sik\={a}mani Venba} (சிகாமணிவெண்பா), a treatise on pulsation, entitled \emph{N\={a}dikottu} (நாடிக்கொத்து), and the fragments of a treatise on hygiene, entitled \emph{N\={o}iyanug\={a}vidhi} (நோயணுகாவிதி).

The following stanzas we quote from Tereiyer's treatise on hygiene, as a specimen of his theories:---
\begin{quotation}\small
திண்ணமிரண்டுள்ளே சிக்கவடக்காமற்‌

பெண்ணின்பாலொன்றைப்‌ பெருக்காம --- லுண்ணுங்கா

னீரருக்கி மோர்பெருக்கி நெய்யுருக்கியுண்பவர்தம்‌

பேருரைக்கிற்போமேபிணி.

உண்பதிருபொழுதொழிய மூன்றுபொழுதுண்ணோ

\hspace*{5mm}முறங்குவதிராவொழியப்‌ பகலுறக்கஞ்செய்யோம்‌

பெண்கடமைத்திங்களுக்கோர்காலன்றிமருவோம்‌

\hspace*{5mm}பெருந்தாகமெடுத்திடினும்‌ பெயர்த்துநீரருந்தோம்‌

மண்பரவுகிழங்குகளிற்‌ கரணையன்றிப்புசியோம்‌

\hspace*{5mm}வாழையிளம்பிஞ்சொழியக் கனியருந்தல்செய்யோம்‌

நண்புபெறவுண்டபின்பு குறுநடையுங்கொள்வோம்‌

\hspace*{5mm}நமனார்க்கிங்கேதுகவை நாமிருக்குமிடத்தே.

ஆறுதிங்கட்கொருதடவை வமனமருந்தயில்வோ

\hspace*{5mm}மடர்நான்குமதிக்கொருக்காற்‌ பேதியுறைநுகர்வோந்

தேறுமதியொன்றரைக்கோர்‌ தரநசியம்பெறுவோந்

\hspace*{5mm}திங்களரைக்கிரண்டுதரஞ்‌ சவளவிருப்புறுவோம்‌

வீறுசதுர்நாட்கொருக்கானெய்முழுக்கைத்தவிரோம்‌

\hspace*{5mm}விழிகளுக்கஞ்சனமூன்று நாட்கொருக்காலிடுவோம்‌

நாறுகந்தம்புட்பமிவை நடுநிசியின்முகரோம்‌

\hspace*{5mm}நமனார்க்கிங்கேதுகவை நாமிருக்குமிடத்தே.
\end{quotation}
\begin{quotation}\small
\emph{He who does not suppress the alvine discharges; he who does not indulge in venery; he who at meals drinks water that has been boiled; he who dilutes well with water his curdled milk; he who uses ghee melted, will enjoy health; for at the very mention of his name disease will depart.}

\emph{We will eat only twice, and not thrice a day; we will sleep only at night, and not in the day; we will hold sexual intercourse but once a month; we will drink water only at meals, although we feel ever so thirsty; we will not eat the bulbous root of any plant, except that of the Dracontium; we will not eat any fruit, except the tender plantain, and we will take a short walk after our meals and then what has Yama (Pluto) to do with us?}

\emph{We will take an emetic once in six months; a purgative once in four months; a Nasya\footnote{A medicine prepared from the juice of certain herbs and injected into the nostrils to promote the discharge of mucus.} once in a month and a half, get shaved (meaning the head) twice in a fortnight; anoint ourselves, and bathe every fourth day; apply collyrium to the eyes every third day; and never smell perfumes, or flowers in the middle of the night; and then what has Yama to do with us?}
\end{quotation}
\paragraph{}
In the \emph{Padh\={a}rtaguna Chintamani} (பதார்த்தகுணசிந்தாமணி), a work on Materia Medica, consisting of 1504 stanzas, and which has been recently recovered from obscurity through the medium of the Press at Madras, by the laudable industry of \textbf{Subramanya Panditer}, we find a great number of the stanzas quoted as Tereiyer's productions. It may be mentioned that this work not only treats of medicinal agents but also of nutrients. It begins with the elements, and climates, and then proceeds to describe the properties of the milk of different kinds of animals and vegetables, the dung and urine of different kinds of animals, the barks, leaves, flowers, fruits, seeds, nuts, gums, and roots of different kinds of trees and plants, the flesh of different kinds of beasts, birds and fish, and the preparations made from different kinds of metals and minerals, and concludes with rules on the mode of living for the preservation of health; and we consider that an English translation of this work will no doubt prove a valuable auxiliary to medical men in England in their researches into the Materia Medica of India.

\section{Todittalei Vizhuttandin\={a}r. --- தொடித்தலைவிழுத்தண்டினார்.}
\paragraph{}
Todittalei Vizhuttandin\={a}r was one of the forty-nine professors of the Madura College. What has contributed to preserve his name from oblivion is his encomiastic stanza on the \emph{Kural}, which we subjoin.
\begin{quotation}\small
அறநான்கறிபொருளேழொன்றுகாமத்‌

திறமூன்றெனப்பகுதிசெய்து --- பெறலரிய

நாலுமொழிந்த பெருநாவலரேநன்குணர்வார்

போலுமொழிந்தபொருள்‌.
\end{quotation}
\begin{quotation}\small
\emph{The great poet has in his Kural divided the topic Virtue into four sections; the topic wealth into seven sections; and the topic pleasure into three sections; including in them the fourth topic, eternal happiness, which is difficult to be obtained: hence the work comprises everything, and if there be anything which it does not comprise, he alone knows it.}
\end{quotation}

\section{Tolk\={a}ppiyan\={a}r. --- தொல்காப்பியனார்.}
\paragraph{}
Tolk\={a}ppiyan\={a}r, otherwise called \emph{Trinadh\={u}m\={a}gni} (திரணதூமாக்கினி), was one of the twelve disciples of \textbf{Agastiyer}, and pre-eminent to them all in learning. He belonged to a \emph{Brahman} family of the tribe of \textbf{Jamadagni}, and \emph{Tolk\={a}ppiyakkudi}, a town situated to the south of Madura, is reported to have been the place of his residence. It is reported of him that having quarrelled with \textbf{Agastiyer}, he set up a rival school, and composed for the use of his scholars a grammar of his own, under the title of \emph{Tolk\={a}ppiyam}, which though based on the grammatical institutes of \textbf{Agastiyer}, yet on the whole deviated from them. \textbf{Athankot\={a}siriyer}, who was also a disciple of \textbf{Agastiyer}, taking the part of his master endeavoured to bring the \emph{Tolkappiyam} into disrepute by severely criticising it, when it was submitted to the judgment of the assembly of the learned in the court of the Pandiya king; but the author defended his work so successfully that it met with due approbation. This work, which is the oldest of the kind in existence, consists of 1276 \emph{S\={u}tras} or rules in verse, and it has had three commentaries written upon it, one by \textbf{Nachchinarkiniyan\={a}r}, another by \textbf{Ilamp\={u}raner}, and the other by \textbf{Senavareiyer} and which are respectively called after the names of the commentators \emph{Nachchin\={a}rkiniyam} (நச்சினார்க்கினியம்), \emph{Ilamp\={u}ranam} (இளம்பூரணம்), and \emph{S\={e}n\={a}vareiyam} (சேனாவரையம்).

\section{Tondaradipp\={o}di \={A}zhv\={a}r. --- தொண்டரடிப்போடியாழ்வார்.}
\paragraph{}
This poet, whose real name was \emph{Vipra Nar\={a}yaner}, ranks high in the estimation of the \emph{Vaishnavas} as one of the twelve special votaries of \emph{Vishnu}. He was born of a \emph{Brahman} family at \emph{Tirumandangkudi}, and after having completed his course of studies in the Vedic and profane literature, became an ascetic, and spent the remainder of his days in ministering in the \emph{Vaishnava} shrine at Sirangam. Of his poetical compositions the \emph{Tirum\={a}lei}, (திருமாலை), and \emph{Tiruppalliezhuchi} (திருப்பள்ளியெழச்சி), which form a part of the \emph{N\={a}l\={a}yiraprabhandam} (நாலாயிரபிரபந்தம்), consist of fifty stanzas in praise of \emph{Vishnu}, and they are deserving of notice.

\section{Ulagan\={a}then. --- உலகநாதன்.}
\paragraph{}
This poet though born of parents who belonged to the \emph{Ambatta}, or barber caste, nevertheless raised himself to eminence by his superior genius. Neither the place of his birth nor the date of his existence is known. He is remembered as the author of two metrical compositions, of which the one entitled \emph{Ulaga Niti} (உலகநீதி), contains a series of moral precepts, and the other, entitled \emph{Jatib\={e}da N\={u}l} (சாதிபேதநூல்), treats of the origin of the divisions, and differences of caste.

\section{Ukraperuvazhuthiy\={a}r. --- உக்கிரப்பெருவழுதியார்.}
\paragraph{}
Ukraperuvazhuthiy\={a}r was one of the forty-nine professors of the Madura College. Some invest him likewise with regal dignity; but no king of such a name appears to have reigned at Madura later than the seventh or sixth century \textsc{bce}, and we cannot therefore admit it without involving ourselves into a serious anachronism. Tradition has handed down the following impromptu, in which he eulogized both \textbf{Tiruvalluver} and his \emph{Kural}.
\begin{quotation}\small
நான்மறையின்மெய்ப்பொருளை முப்பொருளாநான்முகத்தோன்‌

றான்மறைந்து வள்ளுவனாய்த்‌ தந்‌துரைத்த --- நூன்முறையை

வந்திக்க சென்னிவாய்‌ வாழ்த்துகநன்னெஞ்சஞ்‌

சிந்திக்கக் கேட்கச்செவி.
\end{quotation}
\begin{quotation}\small
\emph{The four-faced (Brahma) disguising himself as Valluver has imparted the truths of the four Vedas in the three parts of the Kural, which is therefore to be adored by the head, praised by the mouth, pondered by the mind, and heard by the ears.}
\end{quotation}

\section{Um\={a}pati Siv\={a}chariy\={a}r. --- உமாபதி சிவாசாரியர்.}
\paragraph{}
This celebrated poet and philosopher flourished in the seventeenth century of the Christian era. He belonged to the \emph{Vaishnava Brahmans} resident at \emph{Chillambaram}, in the Carnatic; but being dissatisfied with the philosophy of the \emph{Vaishnava} school, he embraced that of the \emph{Saiva} which he studied thoroughly under the guidance of the \emph{Saiva Guru} \textbf{Marreign\={a}na Sampanter} at Maruth\={u}r. His principal work is the \emph{Sivaprakasam} (சிவப்பிரகாசம்), consisting of a centum of stanzas, illustrative of the doctrines as taught in the \emph{\={A}gamas} respecting the existence of God, his essence and attributes, and of the nature of the human soul, and the organization of the human body. We are glad to notice an abstract of its contents translated into English by the Rev. \textbf{Henry R. Hoisington}, and published in the fourth volume of the Journal of the American Oriental Society. He was also the author of the \emph{Tirutonder Puranas\={a}ram} (திருத்தொண்டர்புராணசாரம்), in which he has, in the narrow compass of seventy-five stanzas, given the substance of the legends of the sixty-four votaries of \emph{Siva} as related by \emph{Sekilar} in 2255 stanzas. He wrote as well the \emph{S\={e}kilar Puranam} (சேக்கிழார்புராணம்), which treats of the life and character of \emph{Sekilar} himself in 100 stanzas, and the \emph{K\={o}yit Puranam} (கோயிற்புராணம்), which treats of the legends of the shrine of \emph{Siva} at Chillambaram in 410 stanzas. The time of his existence is not known; but we find his name mentioned in the introduction to the \emph{Sidambara Pur\={a}nam} (சிதம்பரபுராணம்), which dates \textsc{ce} 1513.

\section{Uppei. --- உப்பை.}
\paragraph{}
This poetess was one of the sisters of \textbf{Tiruvalluver}, and having been exposed by her mother, as soon as born, in a choultry at \emph{Utukad\={u}} in the Arcot district, she was found and brought up by a washerman. She distinguished herself not only by her proficiency in the art of poetry but likewise in the science of moral philosophy, of which she wrote a treatise. The Hindus having deified her after her death, under the title of \emph{M\={a}ri Ammei}, the further particulars of her life have become blended with fable.

\section{Urei\={u}r Muthuk\={u}tanar. --- உறையூர்முதுகூற்றனார்.}
\paragraph{}
The name of this poet has been handed down by tradition as one of the forty-nine professors of the Madura College; but we have no particulars of his life, or any of his compositions, save the subjoined impromptu on the \emph{Kural}.
\begin{quotation}\small
தேவிற்சிறந்த திருவள்ளுவர்குறள்வெண்

பாவிற்சிறந்திடுமுப்பால்‌ பகரார்‌ --- நாவிற்‌

குயலில்லைசொற்சுவையோர்வில்லைமற்றுஞ்‌

செயலில்லையென்னுந்திரு.
\end{quotation}
\begin{quotation}\small
\emph{They who have not studied the Kural of the divine Valluver are incapable of good actions; neither their tongues have expressed what is sweet in language, nor their minds understood what is sublime in sense.}
\end{quotation}

\section{Uruttirasanmer. --- உருத்திரசன்மர்.}
\paragraph{}
This poet was the son of an opulent merchant at Madura. He is said to have been dumb from his birth until he passed his childhood, and then became supernaturally possessed of poetic spirit. Such was his fame as a critic, that when a dispute arose amongst the professors of the Madura College as to superiority, in consequence of each of them having composed a poem, and each vaunted the merits of his own composition, they chose him as arbitrator to settle that dispute. Having heard the professors each recite his poem, he adjudged some to be good in style, but bad in sense; and some to be good in sense but bad in style; and nodded approbation only to those composed by \textbf{Kapiler}, \textbf{Paraner}, and \textbf{Narkirer}, as perfect both in style, and in sense; and the professors were so much satisfied with his judgment that they at once admitted him as their associate. Moreover when the professors were in doubt, as to who amongst them was worthy of sitting on the College bench, and of hearing \textbf{Tiruvalluver} recite his \emph{Kural} it is stated that an aerial voice declared Urithirasanmer to be the person.

The following is his impromptu on the merits of the \emph{Kural}:
\begin{quotation}\small
மணற்கிளைக்கநீரூறு மைந்தர்கள்வாய்வைத்

துணச்சுரக்குந்தாய்முலையொண்பால்‌ --- பிணக்கிலா

வாய்மொழிவள்ளுவர்‌ முப்பான்மதிப்புலவோர்

காய்தொறுமூறுமறிவு.
\end{quotation}
\begin{quotation}\small
\emph{Water will spring forth when the earth is dug, milk will spring forth when the child sucks the mother's breast; but knowledge will spring forth when the poets study Valluver's Kural.}
\end{quotation}

\section{Uruvei. --- உறுவை.}
\paragraph{}
This poetess was a sister of \textbf{Tiruvalluver}, and was found when she was exposed by her mother soon after her birth, and brought up by a man of the \emph{Sanar} caste at \emph{K\={a}veripatnam}. She is said to have been well versed in poetry, and to have composed several poems on different subjects, which, however, have not reached the present time. The Hindus have deified her, as they have done her sister \textbf{Uppei}, and she is now an object of worship at Tiruv\={a}langk\={a}du.

\section{Vannakkanch\={a}ttan\={a}r. --- வண்ணக்கஞ்சாற்றனார்.}
\paragraph{}
Vannakkanch\={a}ttan\={a}r was one of the forty-nine professors of the Madura College.

The following is his stanza on the merits of the \emph{Kural}.
\begin{quotation}\small
ஆரியமுஞ்‌ செந்தமிழுமாராய்ந்திதனினிது

சீரியதேன்றொன்றைச்‌ செப்பரிதா --- லாரியம்‌

வேதமுடைத்துத்‌ தமிழ்‌ திருவள்ளுவனா

ரோதுகுறட்பாவுடைத்து.
\end{quotation}
\begin{quotation}\small
\emph{It is difficult to say whether the Sanskrit or the Tamil is the best, they are perhaps on a par, since the Sanskrit possesses the Veda, and the Tamil the Kural, composed by the divine Valluver.}
\end{quotation}

\section{Varantaruv\={a}n. --- வரந்தருவான்.}
\paragraph{}
This poet was the son of Villiputt\={u}rer; but of his life no memorials have been preserved. It would appear, however, that he was the author of the prefatory stanzas which are affixed to his father's imitation of the \emph{Mah\={a}bh\={a}rat}.

\section{Varathakavir\={a}yer. --- வரதகவிராயர்.}
\paragraph{}
Varathakavir\={a}yer was a poet, who lived at \emph{Chunn\={a}gam}, in Jaffna, in the seventeenth century of the Christian era. He belonged to a \emph{Br\={a}hman} family and was well skilled both in the Vedic and profane literature. His principal works were the \emph{Sivar\={a}tri Puranam} (சிவராத்திரிபுராணம்) and the \emph{Ēk\={a}d\={e}si Puranam}, (ஏகாதேசிப்புராணம்), both which are epic poems, recounting the legends connected with the vigil observed by the \emph{Saivas} on the night preceding the new moon towards the end of February or beginning of March, and the fast on the eleventh day after the full or new moon in every month.

\section{Varatunga R\={a}ma P\={a}ndiyen. --- வரதுங்கராமபாண்டியன்.}
\paragraph{}
Varatunga R\={a}ma P\={a}ndiyen was the younger brother of \textbf{Athivira R\={a}ma P\={a}ndiyen}, and was like him a poet of great attainments. His imitation of the \emph{Bram\={o}thara K\={a}ndam} (பிரமோத்திரகாண்டம்) of the \emph{Skanda Purana} is a work of sterling merit, in so far as its poetry is concerned, and it appears to have been approved of as such by the learned professors of the Madura College to whose criticism it was subjected. It consists of 1323 stanzas, arranged in twenty-two cantos, and has for its subject the legends of certain votaries of \emph{Siva} who had obtained beatitude by bearing his emblems, or uttering his name by accident.

\section{Vellivithiy\={a}r. --- வெள்ளிவீதியார்.}
\paragraph{}
Vellivithiy\={a}r was one of the forty-nine professors of the Madura College. He was present when the \emph{Kural} was exhibited and subjected to the criticism of the College bench, and the subjoined is the stanza which he uttered on the merits of that work.
\begin{quotation}\small
செய்யாமொழிக்குந் திருவள்ளுவர்மொழிந்த

பொய்யாமொழிக்கும்‌ பொருளொன்றே --- செய்யா

வதற்குரியாரந்தணரே யாராயினேனை

யிதற்குரியாரல்லாதாரில்‌.
\end{quotation}
\begin{quotation}\small
\emph{The uncreated word (meaning the Veda) and the unfailing word of Tiruvalluver both have the same meaning; but the former is the exclusive property of the Brahmans (for they alone are allowed to study it), while the latter is the common property of all.}
\end{quotation}

\section{Venb\={a}puli. --- வெண்பாப்புலி.
\paragraph{}
Venb\={a}puli was born of the \emph{Kollery} tribe at Ramnad, in the beginning of the eighteenth century of the Christian era. His real name is not known, the one under which he is here noticed being merely titular, denoting his adroitness in the composition of verses of the \emph{Venba} metre. He appears to have been a very eccentric character; for we are told that shortly after he had finished his studies and obtained a reputation as a poet, he became an ascetic and lived upon alms until he was persuaded by the Raja of Sivaganga to take up his residence in his palace. Whether he composed any larger poems or not is unknown, and we have met with only a few of his stray pieces. It is related of him, that one day as he was carrying a bundle of betel leaves, Tamba Mudeli, the betel renter, met him on the road, and not knowing who he was, asked him where he had stolen the betel leaves; upon which he uttered the following verse extempore:---
\begin{quotation}\small
தம்பாமுதலி சகலகலைக்கியானமுள்ள

வெண்பாப்புலியென்றால்வேற்குமோ

குத்தகையில்வெற்றிலையைக்‌ கொண்டுவந்தாலென்னபயம்‌

புத்திகெட்டமாடே நீ போ.
\end{quotation}
\paragraph{}
\textbf{Tamba Mudeli} immediately knowing to whom he was speaking apologized for his conduct and they both parted friends.

\section{V\={e}daghiri Mudeliar. --- வேதகிரிமுதலியார்.}
\paragraph{}
Vedaghiri Mudeliar was born of a \emph{Tuluva Vel\={a}la} family at \emph{Kulatt\={u}r}, in the Carnatic, \textsc{ce} 1795. Having made some progress in the study of grammar and poetry in his own native place, he repaired to Madras, where he had the advantage of perfecting them under the able tuition of \textbf{Ramanuja Kavir\={a}yer}, and where, he soon afterwards became the head of a Tamil Literary Institution which his master had originated. Sometime afterwards he was invited over to Madura, by the principal native gentlemen of that town, to whom his rising fame had reached, and there continued to preside over a Tamil Literary Institution for a period of seven years, when owing to a failure in his health he was obliged to relinquish his post and repair to Pondicherry, where for some time he was employed as a Tamil Pundit in the Catholic College. He afterwards went a second time to Madras, and having established a Press of his own, commenced publishing through its medium several valuable classical works, such as the \emph{Kural}, \emph{Naidatham}, \emph{Illakkanakottu}, etc. with a learned commentary of his own in prose on each; and while thus usefully engaged, his career was closed by death in the year 1852, at the age of fifty-seven.

Vedaghiri Mudeliar was a frequent contributor to the Jaffna \emph{Morning Star}, from 1841 to 1843, and his contributions, which were chiefly on Grammar, and Classics, tended very much to enhance the value of that periodical, and drew forth to him from the Editors their high tribute of thanks through the following stanza:---
\begin{quotation}\small
ஞானவடிவானசுடர்‌ பாதமலர்சூடியவனாமமறவாதபுனிதன்‌ --- நாதனருள்வேதமறை யோதுமுறைபோதருவினோதநிறைஞானசரிதன்‌ --- நாளுமநுநீதிவழுவாதுபரனேசவிசுவாசரெவர்நாடுமறிஞன்‌ --- நாவலர்சிரோமணி யெனாவொளிருபோசனுயர்‌ ஞாலமதிராசகவிஞன் --- தேனமிர்தமாமதுரபாடலெனுநீரிலகு தேவதமிழாயகடலே --- சீரியகுசாநுனிய னாமதியெனோடமிசை தேடியதனாலுகரையுந் --- தேறிமணிநேருகலை மாலைபுனைவீரனுயர்‌ மேருகுணதீரசதுரன்‌ --- தேர்புலவர்யாருமொரு தூயவடிவாயதென வேதிகழ்விலாசகலைஞன்‌ --- வானகரமேயதரு நாணவளமாரியிடுதாரனறியாதெதிர்புகும்‌ --- வாதிகளைமோதுமது சூதனிகர்வேதகிரி மாதவனினாதுரமுறா --- வாலியதொனூல்பலவினாலருமைசாலும்விதி வாய்மையொடாராய்வுபுரியா --- மாமதியைநாளும்வளர்‌ தாரகையினாலுலக வானிருளையேகடிசெய்தான்‌ --- மேனிகரின்ஞானதபனோதயனுமாநிலமின்மீதுபிரகாசமதுவாய்‌ --- மேயினசெயாதருளுமோருதவியீததனுண்‌ மேதகையபோதவுதவி --- மீனவனினானதமிழ்‌ மானைமணமேவுபுய வேதகிரிராசனுதவி --- மேதினியில்‌யாதுநிகரோதமிளிர்பூதலமு மீதலமுமேனுமரிதே.

\hspace*{10mm}See Morning Star, vol. 1. p. 280.
\end{quotation}
\paragraph{}
He did not, however, compose any epic poems. All that he has left as specimens of his poetical talents are four didactic poems, respectively entitled \emph{Manun\={i}ti Sadagam} (மனுநீதிசதகம்); \emph{Manvikiy\={a}na Sadagam} (மனுவிக்கியானசதகம்), \emph{N\={i}ti Chint\={a}mani} (நீதிசிந்தாமணி) and \emph{Sanm\={a}rgas\={a}ram} (சன்மார்க்கசாரம்), of which the first three contain a hundred stanzas each, and the last a hundred and ten stanzas. These poems were inscribed by him to a friend of his in Ceylon, whose name he has introduced at the end of every stanza. The Addendum to the Tamil Dictionary published at the American Mission Press at Manipy in 1842, was his production as well as the stanzas marked with asterisks in the 11th part of \textbf{Mandalapuruder's} \emph{Nigandu} published at the same Press in 1843.

\section{V\={e}laya D\={e}siker. --- வேலையதேசிகர்.}
\paragraph{}
V\={e}laya D\={e}siker was born at \emph{Kanjipuram}. After the death of his father \textbf{Kumarasw\={a}mi Pand\={a}ram}, he repaired with his brothers \textbf{Sivaprak\={a}sa} and \textbf{Karunaiprak\={a}ser} to Tinnevelly, where under the tuition of \textbf{Velliambala Tambir\={a}n} the head ascetic of the \emph{Matam} of Sindupoonturei, he studied the Tamil Grammar and Classics, and became a renowned poet. He composed the following poems, viz: the \emph{Nall\={u}r Pur\={a}nam} (நல்லூர்ப்புராணம்), \emph{V\={i}rasing\={a}dana Puranam} (வீரசிங்காதனபுராணம்), \emph{Ishtalingakaittala M\={a}lei} (இஷடலிங்கக்கைத்தலமாலை), \emph{Namasiv\={a}ya L\={i}lei} (நமசிவாயலீலை), \emph{Krishnacharitrap\={a}ris\={a}da L\={i}lei} (கிருஷணசரித்திரபாரிசாதலீலை), and \emph{Mailattirattei M\={a}lei} (மயிலத்திரட்டைமாலை), besides completing the \emph{K\={a}lati Pur\={a}nam} (காளத்திபுராணம்) which his brothers had left unfinished. He died at Perum\={a}tur in the seventy-second year of his age.

The following is an elegiac verse which he recited extempore on the death of his two brothers:---
\begin{quotation}\small
அல்லிமலர்ப்பண்ணவனுமாராய்ந்தறிகவிதை

சொல்லுமிருவரிடைத் தோன்றியயான்‌ --- முல்லை

யரும்பிற்பொலியுமணிமுறுவனல்லாய்‌

கரும்பிற்கணுநிகர்த்தேன்காண்‌.
\end{quotation}
\begin{quotation}\small
\emph{O Damsel, whose teeth vie with the Jessamine flower in beauty! Born as I was second of three brothers, whose poetry, even the God enthroned on the lotus-flower could not easily understand, alas! I have now become like a joint in the stalk of a sugar-cane (cut out and thrown aside as refuse).}
\end{quotation}

\section{Villip\={a}kam T\={a}ndavar\={a}ya Mudeliar. --- தாண்டவராயமுதலியார்.}
 \paragraph{}
This poet was first a student and afterwards a Teacher in the College of Fort St. \textbf{George}, and having been well versed both in the Tamil and Maharatta languages enjoyed much reputation on that account. In 1825, by order of the Board of College he wrote a Tamil Grammar in a catechetical form under the title of \emph{Illakkana Vin\={a}videi} (இலக்கணவினாவிடை) and in the following year translated the \emph{Panchatantra} stories (பஞ்சதந்திரக்கதை) from Maharatta into Tamil, besides himself producing a series of instructive fables called \emph{Kat\={a}manjari} (கதாமஞ்சரி). He also revised and edited the three first parts of \textbf{Beschi's} \emph{Sadur Agar\={a}di}, the ten first parts of the \emph{S\={u}damani Nigandu}, and the first eight parts of the \emph{Div\={a}karam}, which had hitherto been circulated in manuscript replete with errors of the copyists.

\section{Vidv\={a}n S\={a}min\={a}da Pillei. --- வித்துவான் சாமிநாதபிள்ளை.}
\paragraph{}
Vidv\={a}n S\={a}min\={a}da Pillei, who distinguished himself by his literary attainments, and enjoyed a prominent position among the poets of his day, was born of a \emph{Vel\={a}la} family at Pondicherry in the latter part of the last century. His parents not only secured for him a classical education, but likewise brought him up in the Catholic religion which they professed. Having removed his residence to Madras, he was introduced to the notice of Mr. \textbf{Ellis}, and was employed by that celebrated orientalist as his preceptor in the Tamil language. His juvenile productions the \emph{Nazarai Kalambagam} (நசரைக்கலம்பகம்), and \emph{Samin\={a}den Pilleitamil} (சாமிநாதன்பிள்ளைத்தமிழ்), though possessed of poetic beauties did not however attract much notice, being merely panegyrics. But the epic which he wrote in his maturer years, under the title of \emph{Gn\={a}n\={a}diku R\={a}yer K\={a}ppiyam} (ஞானாதிக்கராயர்காப்பியம்), commanded general admiration, and spread his fame throughout Southern India. In the \emph{Gn\={a}n\={a}diku R\={a}yer K\={a}ppiyam}, he has introduced many historical passages from the Bible, by way of episodes in the life of the hero of the poem, in a style and language highly creditable to him. As usual with the writers of epics in Tamil, he opens his poem after the invocation with an eulogy on rain, and the following two stanzas as quoted by Mr. \textbf{Ellis}, in his comments on the \emph{Kural}, form the first part of it:
\begin{quotation}\small
வானமாலியுங்‌ காய்வதேதெனப்புவிவதைப்ப

மானமாவுற வெஞ்சுடர்மாரியேபடைகொண்

டூனமாக்கெனவுச்சியிலம்பிடிபொழிந்து

தானமேகரைவாக்கினசரிவுநீர்தோற்றம்‌.

பல்லுயிர்க்கொரு தந்தையாம்பருதியவ்வுயிரைக்‌

கொல்லுருத்தலைக்கூவெனுந்தாயுளத்திரங்கிச்‌

செல்லியைதிரைத்தெரிசனமறைக்கவுமகர்க்குக்

சல்லியன்‌முலைப்பால்‌ சுரந்தளித்தெனக்காண்டும்.
\end{quotation}
\begin{quotation}\small
\emph{The earth, impatient of the scorching heat of the sun of heaven, to annoy him assumed the form of a stupendous mountain; the bright burning sun, calling forth the armies of the clouds, ordered them to reduce the rebellious earth; then the thunderbolts darted forth like arrows and like water flowing down seemed to dissolve the entire mountain.}

\emph{When the sun, the father of various living beings, destroys in his burning wrath his offspring, their mother earth, inwardly pitying them, hides her children from his view by a propitious veil of clouds and seems to cherish them with milk springing from her breasts, the rocky mountains.}
\end{quotation}

\section{Villiputt\={u}rer. --- வில்லிபுத்தூரர்.}
\paragraph{}
Villiputt\={u}rer was born of a \emph{Br\={a}hman} family of the \emph{Vaishnava} sect at \emph{Saniy\={u}r}, in the Coimbatore country, in the fifth century of the Christian era. He was no less distinguished as one of the \emph{\={A}zhv\={a}rs} or special votaries of \emph{Vishnu} than as a poet. Having refused to give his brother his share of inheritance from the paternal estate, the brother applied to the king for justice against him; but the king without pronouncing any judgment on the matter directly sent for him, and requested him to imitate in Tamil, the Mahabh\={a}rat. On the day the work was finished and recited in the royal presence his brother by desire of the king repeated his complaint. The king thereupon feigning to be angry looked fiercely at the poet, and asked him ``Oh! is this the perception of the great man who wrote the \emph{Mahabh\={a}rat}?'' The king's question struck his mind with shame, for the subject of the \emph{Mahabh\={a}rat} is the feud which originated in the elder brother refusing to give the younger his paternal share of inheritance and he thereupon delivered up the whole of the patrimony to his brother and retiring to Sirangam spent the remainder of his life in devotional exercise. His imitation of the \emph{Mahabh\={a}rat} includes only the first ten \emph{Parvas}, or books of that work, and consists of 3373 stanzas. It is decidedly a masterpiece, for scarcely any Tamil epics, except perhaps \textbf{Kamber's} \emph{Ramayanam}, can vie with it either in ``vivida vis animi'' or charms of language and versification.

\section{V\={i}rap\={a}ndiya Pulaver. --- வீரபாண்டியப்புலவர்.}
\paragraph{}
This poet who lived in the latter part of the eighteenth century of the Christian era, was the son of a tailor at Tutocoryn. Imbibing a dislike to his father's profession, he applied himself to the study of poetry, and obtained some reputation in the literary circle of his day. He was first patronized by \textbf{Katabomma Nayaker}, the Poligar of P\={a}njalamkurichi, and afterwards by \textbf{Don Gabriel Vaz Faldano}, the Parava Chieftain of Tutocoryn. He was a great wit, in proof of which many anecdotes are related of him. It is said that on a certain day as he went to witness a Nautch, a poet of the shepherd tribe, pointing at the dancing girl, observed to him, ``\emph{Taiyel nalla taiyel}'' (a phrase which may mean either, ``the girl is a fine girl,'' or ``the sewing is good sewing),'' and he immediately understanding that to be a hint at his caste, made this repartee ``\emph{Attukku tirrama taiyel}'' (a phrase which may mean either ``Is the girl better than the dance?'' or ``is the sewing better than the sheep?'') and completely outwitted him. Of his compositions there are only a few stray pieces, chiefly devoted to the praises of his benefactors.

\section{V\={i}rakavir\={a}yer. --- வீரகவிராயர்.}
\paragraph{}
V\={i}rakavir\={a}yer was a distinguished poet who lived at \textbf{Nell\={u}r}, in the Carnatic. In \textsc{ce} 1524 he wrote an epic poem, called \emph{Arichandra Pur\={a}nam} (அரிச்சந்திரபுராணம்), which drew forth much applause from the savants at Madura, when it was submitted to their inspection. This poem consists of 1212 stanzas, arranged into twelve cantos, and recounts the adventures of the king \emph{Harischandra}. Though it cannot stand on a par with the \emph{Naidatham} in the charms of language and versification; yet it is far superior to the ordinary strain of current poetry and bears the stamp of no common talent. The following stanzas, which we have extracted from the canto, entitled \emph{May\={a}na K\={a}ndam}, containing the lamentations of the Queen \emph{Chandramati} on finding her son lying dead in a wood by the bite of a serpent, are very pathetic and affecting:---
\begin{quotation}\small
பனியானனைந்து வெயிலாலுலர்ந்து பசியாலலைந்துமுலவா

வனியாயவெங்கணரவாலிறந்தவதிபாவமென்கொலறியேன்‌

தனியேகிடந்து விடநோய்செறிந்து தரைமீதுருண்டமகனே

யினியாரைநம்பியுயிர்வாழ்வமென்றனிறையோனும்யானுமவமே.

நிறையோசையுற்ற பறையோசையற்றுநிரையாய்நிறைந்தகழுகின்‌

சிறையோசையுற்ற செடியூடிறக்க விதியாரிழைத்தசெயலோ

மறையோனிரக்கவளநாடனைத்தும்‌ வழுவாதளித்தவடிவே

லிறையோனளித்தமகனேயுனக்குமிதுவோவிதித்தவிதியே.

வானின்றிழிந்துசொரிகின்றதாரைமழைபோலவீழவிழிநீ

ரூனின்றிமேனியுதிரங்கள்‌சிந்தவுயிரின்றிவெம்புதழலின்‌

மேனின்றுவெந்ததளிர்போலுயங்கிவிதியாரைநொந்துதனியே

யானின்றிரங்கவேனென்பதில்லையிதுவோவுனீதிமகனே.

செங்கோலறத்தின்‌ முறையேசெலுத்துதிறலோனெவர்ககுமுரவோன்‌ (ள்‌

வெங்கோபயானைவிறன்மன்னனம்மைவிடுவிக்கவெண்ணிவருநா

பங்கேருகத்துமலர்போல்‌விளங்குவதனாமகிழ்ந்தபரிவா

லெங்கேயெனாசைமகனென்றுரைக்கினினியேதுசொல்வன்மகனே.
\end{quotation}
\begin{quotation}\small
\emph{I do not know what heinous sin I have committed that thou shouldst meet with thy death from the cruel serpent with protuberant eyes and which suffers neither from being drenched with dew, nor scorched with the sun, nor from want of food. Oh! my son, who, when affected by the venom didst lay thyself all alone and writhe in agony on the ground! To whom shall my lord and myself look for help hereafter?}

\emph{Wert thou fated to die in a wood, where instead of the sound of musical instruments, only the sound of the flappings of the wings of eagles is heard? Oh my son! why was this fate ordained for thee, the son of the king, who wields the bright javelin, and who made over his fertile kingdom to the suppliant Brahman.}

\emph{Whilst with tears streaming from mine eyes, like rain from the clouds, I stand by thy side as a green twig burning in the fierce fire, and lament thy fate, Oh my son! is it proper that thou shouldst not respond to my call?}

\emph{Oh! my son, whose face resembles the lotus-flower! When thy father, who sways the sceptre with justice and equals the fierce elephant in strength, shall return to redeem us and ask me ``where is my darling son?'' what shall I say to him?}
\end{quotation}
\section{Visvan\={a}tha S\={a}striy\={a}r. --- விசுவநாதசாஸ்திரியார்.}
\paragraph{}
Visvan\={a}tha S\={a}striy\={a}r was both a poet and astronomer who flourished among the \emph{Brahmans} at Ar\={a}ly, a village of \emph{Batticotta}, in Jaffna. Sir \textbf{Emerson Tennent}, in his ``Christianity in Ceylon,'' states, that Visvan\={a}tha S\={a}striy\={a}r, or as he calls him \emph{Visuvenathen}, was ``the most celebrated of the astronomers of Ceylon, inasmuch as he concentrated in his own person the accumulated science of his ancestors, who for nine generations had been cultivators of the same study.''\footnote{Christianity in Ceylon, p. 139.} His Almanac, which he published every year until his death, was preferred to all others as being more accurate; and Sir \textbf{Alexander Johnston} while Chief Justice of Ceylon, procured for him from \textbf{George} the 4th, the honorary distinction of Almanac Maker to his Majesty.''\footnote{Pridham's Ceylon and its Dependencies, vol. 2. p. 516.} But it should be mentioned, that in his controversy with the American Missionaries at Batticotta, respecting the beginning and extent of the obscuration and the ultimate duration of the eclipse of the moon on the 21st of March 1828, he did not come off triumphant; for that eclipse began fifteen minutes sooner and continued for twenty-four minutes longer than his calculation and it obscured only three-eighths and not five-eighths of the disc as he had predicted. As a poet he has distinguished himself by the production of two poems, the \emph{M\={a}veikuravanji} (மாவைக்குறவஞ்சி), a drama founded on the legend of the horse-faced \emph{Chola} princess who is fabled to have been cured of her deformity by bathing in the sea, near Kirimalei, and the \emph{Kurun\={a}ther Killeividud\={u}tu} (குருநாதர் கிள்ளைவிடுதூது), a panegyric on \emph{Skanda}.

\section{Vytian\={a}tha N\={a}valer. --- வைத்தியநாதநாவலர்.}
\paragraph{}
This poet was a native of \emph{Tiruv\={a}r\={u}r}, a town in the Tanjore country, near Kumbakonam, and lived at the same time with \textbf{S\={a}min\={a}da Tambiran}, whose friendship he enjoyed. He is celebrated for his high attainments in philology, and we have a work of his on Grammar, entitled \emph{Ilakkana Vilakkam} (இலக்கணவிளக்கம்), written chiefly as a critique on the \emph{Nann\={u}l}.

\section{Y\={a}lp\={a}na N\={a}yan\={a}r. --- யாழ்ப்பாணநாயனார்.}
\paragraph{}
Y\={a}lp\={a}na N\={a}yan\={a}r, otherwise called \emph{Virar\={a}gaven}, was a minstrel who lived in the \emph{Chola} country. Being blind, he depended for his subsistence entirely on the earnings of his wife. One day his wife having delayed to serve out to him his meals at the proper hour, he quarrelled with her on that account, and quitted his house saying, that he was going to Ceylon, upon which she sneeringly observed. ``Ah! you are going to Ceylon to obtain a tusked elephant and a fertile land.'' When he arrived in Ceylon, he was refused admittance into the king's presence, as it was considered ominous for a king to see a blind man; but it was afterwards arranged that the king should stand behind a curtain, and hear the blind minstrel's song. The king divesting himself of his royal garments, put on a pair of short drawers like those worn by an archer, and stood behind the curtain with a bow in his hand, and bade him sing. He, somehow or other, being informed of the king's disguise improvised the following stanza characterizing the king as \textbf{R\={a}m\={a}} of old:---
\begin{quotation}\small
வாழுமிலங்கைக்கோ மானில்லை மானில்லை

யேழுமராமரமுமிங்கில்‌லை --- ஆழி

யலையடைத்ததிண்புயாவாதித்தாவுன்கை

சிலையெடுத்தவாறேதுசெப்பு.
\end{quotation}
\begin{quotation}\small
\emph{O Aditya! (meaning R\={a}ma) whose potent arms stemmed the waves of the sea, pray why have thy hands assumed the bow? for the prosperous Lanka has neither her king (meaning R\={a}vana), nor the stag (meaning M\={a}richi,) who in the disguise of a stag decoyed away R\={a}ma from Sita), nor the seven Mara trees\footnote{Diospyros Ebenaster.} (meaning those which R\={a}ma shot through with a single arrow) to shoot at.}
\end{quotation}
\paragraph{}
The king immediately laid aside the garb of the archer, and successively putting on the robes of the Queen, and his own, desired him to sing. He thereupon composed extempore a lyric in ten stanzas to suit the respective guises of the king, and sang it to the lute which he himself played. The king being greatly pleased with his performance, honoured him by the gift of a tusked elephant, and by the donation of a land on the northern extremity of the Island in perpetuity; and thus was realized what his wife had said in bitter irony.

The land which he obtained from the king was no other than the present peninsula of Jaffna. It was then uninhabited, and covered with jungle, but he had it cleared, and having induced a colony of Tamils from Southern India to settle in it, soon rendered it a rich country which he called after his own professional name \emph{Y\={a}lp\={a}na nadu}, i. e. the minstrel's country. He did not, however, assume any sovereignty over the country but inviting over from the Continent a youth, who was the natural son of a \emph{Chola} king by a \emph{Brahman} woman, installed him as its king, under the title of \emph{Sing\={a}riya Chakravarti}. The installation of \emph{Sing\={a}riya Chakravarti} is stated in the \emph{Kailasa M\={a}lei} to have occurred in the \emph{Kali} year 3000 (\textsc{ce} 101); and if this date be correct, the \emph{Singhalese} monarch, who patronized the blind minstrel, must have been \emph{Wasabha}, who reigned at Anuradhapura from \textsc{ce} 62 to 106. In the lyric, which has been noticed above, the minstrel does not celebrate his patron under his proper name, he styles him merely \emph{Parar\={a}jasingan}, an appellation signifying one who is as a lion to other kings.

\textbf{Finis.}

\end{document}
